% English translation, made from our own transcription (original.tex), not from any existing
% translation. Every mathematical expression, formula, symbol, \origpage{N} marker, tag,
% environment, and structural anchor is reproduced unchanged; only the prose is translated.
% The editorial notes (\ednote) of the transcription are carried over unchanged, since they are
% already in English and some carry mathematics.
%
% TERMINOLOGY. Poincaré's own terms are rendered literally and consistently:
%   * variété = "manifold"; variété à n dimensions = "manifold of n dimensions"
%   * domaine = "domain"; frontière (complète) = "(complete) boundary"
%   * finie / continue = "finite" / "continuous"
%   * limitée / illimitée = "limited" / "unlimited" (a manifold with / without a boundary);
%     fermée = "closed"
%   * homéomorphe, homéomorphisme = "homeomorphic", "homeomorphism"
%   * variétés opposées = "opposite manifolds"
%   * unilatère / bilatère = "one-sided" / "two-sided"
%   * homologie, équivalence = "homology", "equivalence"; nombres de Betti = "Betti numbers"
%   * groupe fondamental = "fundamental group"; équivalences / homologies fondamentales =
%     "fundamental equivalences / homologies"
%   * déterminant fonctionnel = "functional determinant"; Tableau = "array" (§ 1) or "Table"
%   * coupure = "cut"; lèvres = "lips"; nappe = "sheet"; lacet = "loop"; contour fermé =
%     "closed contour"; faces conjuguées = "conjugate faces"; aster = "aster" (Poincaré's word)
%   * quadrillage = "gridding", mailles = "meshes"; annexer = "annex"
%   * isomorphisme holoédrique / mériédrique = "holohedral / merihedral isomorphism"
%   * C. Q. F. D. = "Q. E. D."; the combinatorial "m à m" = "m by m"
%   * the French space before a colon (~:) is set tight, as in English and in poincare-1910
%
% PRINTER'S ERRORS. The transcription's hedged \ednote on each apparent misprint is carried over
% verbatim. Where a misprint is a misspelt word whose literal rendering would be meaningless in
% English (p. 40 "on" for "ou", p. 46 "précisément" for "préciser", p. 62 "et" for "à", p. 63
% the doubled "que", p. 110 "naisasnce", p. 114 "relations" for "relatives"), the intended word
% is translated and the note records what is printed. Misprints that still make sense in English
% (p. 28 "one-sided" for "two-sided", p. 96 "equally distinct") are translated as printed.
%   * \emph{Analysis Situs} is kept in Latin, as in the original.
%   * "M.~X" before a name = "Mr.~X", as in the poincare-1910 translation.
%
% Notation and layout decisions shared across pages are recorded in notation.md.
\documentclass{article}
\usepackage{readmasters}
\begin{document}

\origpage{1}

\section*{Analysis situs;}

By Mr.~H.~Poincaré.

\section*{Introduction.}

Geometry of $n$ dimensions has a real object; nobody doubts this today. The beings of hyperspace are susceptible of precise definitions, like those of ordinary space, and if we cannot represent them to ourselves, we can conceive and study them. If, therefore, Mechanics of more than three dimensions, for example, must be condemned as devoid of any object, the same is not true of Hypergeometry.

Geometry, indeed, does not have as its sole reason for being the immediate description of the bodies that fall under our senses: it is above all the analytic study of a group; nothing prevents us, consequently, from taking up other groups, analogous and more general.

But why, it will be said, not keep the analytic language rather than replace it by a geometric language, which loses all its advantages as soon as the senses can no longer intervene. It is because this new language is more concise; and next because the analogy with ordinary Geometry can create fruitful associations of ideas and suggest useful generalizations.
\origpage{2}

Perhaps these reasons are not sufficient? It is not enough, indeed, that a science be legitimate: its usefulness must be beyond dispute. So many diverse objects solicit our attention that only the most important have the right to obtain it.

So there are parts of Hypergeometry in which there is no reason to take much interest: such are, for example, the researches on the curvature of surfaces in space of $n$ dimensions. One is sure in advance of obtaining the same results as in ordinary Geometry, and one does not undertake a long voyage to find again sights just like those one meets at home.

But there are problems where the analytic language would be altogether inconvenient.

One knows the usefulness of geometric figures in the theory of imaginary functions and of integrals taken between imaginary limits, and how much one misses their help when one wishes to study, for example, the functions of two complex variables.

Let us try to account for the nature of this help; figures first make up for the infirmity of our mind by calling our senses to its aid; but that is not all. It has very often been repeated that Geometry is the art of reasoning well on badly made figures; yet these figures, if they are not to deceive us, must satisfy certain conditions; the proportions may be grossly altered, but the relative positions of the various parts must not be overturned.

The use of figures therefore has above all the aim of making known to us certain relations between the objects of our studies, and these relations are those with which a branch of Geometry is concerned that has been called \emph{Analysis Situs}, and which describes the relative situation of the points of lines and surfaces, without any consideration of their magnitude.

There are relations of the same nature between the beings of hyperspace; there is therefore an \emph{Analysis Situs} of more than three dimensions, as Riemann and Betti have shown.

This science will make known to us this kind of relations, although this knowledge can no longer be intuitive, since our senses fail
\origpage{3}
us. It will thus, in certain cases, render us some of the services that we ordinarily ask of the figures of Geometry.

I shall confine myself to three examples.

The classification of algebraic curves into genera rests, according to Riemann, on the classification of real closed surfaces, made from the point of view of \emph{Analysis Situs}. An immediate induction makes us understand that the classification of algebraic surfaces and the theory of their birational transformations are intimately bound up with the classification of real closed surfaces of space of five dimensions from the point of view of \emph{Analysis Situs}. Mr.~Picard, in a Memoir crowned by the Académie des Sciences, has already insisted on this point.

On the other hand, in a series of Memoirs published in the \emph{Journal de Liouville}, and entitled: \emph{Sur les courbes définies par les équations différentielles}, I employed ordinary \emph{Analysis Situs} of three dimensions in the study of differential equations. The same researches have been pursued by Mr.~Walther Dyck. One easily sees that generalized \emph{Analysis Situs} would allow one to treat in the same way the equations of higher order and, in particular, those of Celestial Mechanics.

Mr.~Jordan has determined analytically the groups of finite order contained in the linear group in $n$ variables. Mr.~Klein had previously, by a geometric method of rare elegance, solved the same problem for the linear group in two variables. Could one not extend Mr.~Klein's method to the group in $n$ variables, or even to an \emph{arbitrary continuous group}? I have not so far been able to achieve this, but I have reflected a great deal on the question, and it seems to me that the solution must depend on a problem of \emph{Analysis Situs}, and that the generalization of Euler's celebrated theorem on polyhedra must play a part in it.

I do not therefore believe that I have done a useless work in writing the present Memoir; I only regret that it is too long; but when I tried to restrain myself, I fell into obscurity; I preferred to pass for a little talkative.
\origpage{4}

\section*{I.}

\subsection*{\S~1. --- First definition of manifolds.}

Let $x_{1}, x_{2}, \ldots, x_{n}$ be $n$ variables that I may regard as the coordinates of a point in space of $n$ dimensions.

Until further notice, I shall suppose that these $n$ variables always remain real.

Any system of values of $n$ variables will be called \emph{a point}.

Let us consider the following system, formed of $p$ equalities and $q$ inequalities:
\[
\left\{
\begin{aligned}
&F_{1}(x_{1}, x_{2}, \ldots, x_{n}) = 0,\\
&F_{2}(x_{1}, x_{2}, \ldots, x_{n}) = 0,\\
&\ldots\ldots\ldots\ldots\ldots\ldots,\\
&F_{p}(x_{1}, x_{2}, \ldots, x_{n}) = 0,\\
&\varphi_{1}(x_{1}, x_{2}, \ldots, x_{n}) > 0,\\
&\varphi_{2}(x_{1}, x_{2}, \ldots, x_{n}) > 0,\\
&\ldots\ldots\ldots\ldots\ldots\ldots,\\
&\varphi_{q}(x_{1}, x_{2}, \ldots, x_{n}) > 0.
\end{aligned}
\right.
\tag{1}
\]

I shall suppose that the functions $F$ and $\varphi$ are uniform and continuous and that they have continuous derivatives; I shall suppose moreover that, if one forms the array,
\[
\begin{matrix}
\frac{dF_{1}}{dx_{1}}, & \frac{dF_{1}}{dx_{2}}, & \ldots, & \frac{dF_{1}}{dx_{n}},\\[1ex]
\frac{dF_{2}}{dx_{1}}, & \frac{dF_{2}}{dx_{2}}, & \ldots, & \frac{dF_{2}}{dx_{n}},\\[1ex]
\ldots, & \ldots, & \ldots, & \ldots,\\[1ex]
\frac{dF_{p}}{dx_{1}}, & \frac{dF_{p}}{dx_{2}}, & \ldots, & \frac{dF_{p}}{dx_{n}},
\end{matrix}
\]
and one forms the determinants obtained by taking any $p$ columns
\origpage{5}
in this array, I shall suppose, I say, that these determinants never all vanish at once.

I shall say that the set of points that satisfy the conditions (1), if there are any, which I suppose, forms a \emph{manifold of $n-p$ dimensions}. (If in particular $p = 0$, so that there are only inequalities, I shall have a manifold of $n$ dimensions which will be nothing other than a portion of space of $n$ dimensions; I shall ordinarily designate this manifold by the name of \emph{domain}.)

Two cases may present themselves. Let $\alpha_{1}, \alpha_{2}, \ldots, \alpha_{n}$; $\beta_{1}, \beta_{2}, \ldots, \beta_{n}$ be two points satisfying the conditions (1). Either one will be able to vary $x_{1}, x_{2}, \ldots, x_{n}$ in a continuous manner from $\alpha_{1}, \alpha_{2}, \ldots, \alpha_{n}$ to $\beta_{1}, \beta_{2}, \ldots, \beta_{n}$ without the relations (1) ceasing to be satisfied, and this whatever the values of $\alpha_{1}, \alpha_{2}, \ldots, \alpha_{n}$ and of $\beta_{1}, \beta_{2}, \ldots, \beta_{n}$, provided these values satisfy the conditions (1); or else this will not always be possible.

In the first case, we shall say that the manifold defined by the conditions (1) is \emph{continuous}.

In what follows I shall consider only continuous manifolds and, as regards those that are not continuous, I shall confine myself to observing that they can always be decomposed into a finite or infinite number of continuous manifolds.

Let there be, for example, the manifold
\[
x^{2}_{2} + x^{4}_{1} - 4x^{2}_{1} + 1 = 0.
\]

Here $n = 2$ and the point $x_{1}, x_{2}$ is a point of the plane; our manifold is then nothing other than a curve of the 4th degree; but, since this curve consists of two closed branches of curve, our manifold is not continuous.

But we can decompose it into two others, namely:
\[
x^{2}_{2} + x^{4}_{1} - 4x^{2}_{1} + 1 = 0, \qquad x_{1} > 0
\]
and
\[
x^{2}_{2} + x^{4}_{1} - 4x^{2}_{1} + 1 = 0, \qquad x_{1} < 0
\]
\origpage{6}
and each of them, reducing to a single closed branch of curve, is continuous.

I shall say that a manifold is \emph{finite} if all its points satisfy the condition
\[
x^{2}_{1} + x^{2}_{2} + \ldots + x^{2}_{n} < K^{2},
\]
$K$ being a given constant.

Let us now consider the system of relations
\[
\left\{
\begin{aligned}
&F_{\alpha} = 0 &&(\alpha = 1, 2, \ldots, p),\\
&\varphi_{\beta} = 0,\\
&\varphi_{\gamma} > 0 &&(\gamma \gtrless \beta)
\end{aligned}
\right.
\tag{2}
\]
which consists of $p + 1$ equalities and $q - 1$ inequalities.

It may happen either that there is no point satisfying the conditions (2), or that there are some, and in that case they will form a manifold, continuous or not, which will have fewer than $n - p$ dimensions.

The set of points that satisfy one of the $q$ systems of relations:
\[
\left\{
\begin{aligned}
&F_{\alpha} = 0, &&\varphi_{1} = 0, &&\varphi_{\gamma} > 0 &&(\gamma \gtrless 1),\\
&F_{\alpha} = 0, &&\varphi_{2} = 0, &&\varphi_{\gamma} > 0 &&(\gamma \gtrless 2),\\
&\ldots\ldots, &&\ldots\ldots, &&\ldots\ldots &&\ldots\ldots,\\
&F_{\alpha} = 0, &&\varphi_{q} = 0, &&\varphi_{\gamma} > 0 &&(\gamma \gtrless q)
\end{aligned}
\right.
\tag{3}
\]
will be called the \emph{complete boundary} of the manifold defined by the conditions (1). But we shall sometimes place ourselves at another point of view, and shall regard as a true boundary manifold only those that have $n - p - 1$ dimensions.

It may be that there is no manifold of $n - p - 1$ dimensions that satisfies one of the $q$ systems (3). In that case the manifold defined by the conditions (1) will be called \emph{unlimited}. In the contrary case, it will be called \emph{limited}.
\origpage{7}

If a manifold is at once finite, continuous and unlimited, it will be called \emph{closed}.

To shorten the language a little, we shall give the name of \emph{surfaces} to manifolds of $n - 1$ dimensions, except in the case $n = 2$; and we shall give, in all cases, the name of \emph{curves} to manifolds of one dimension.

\subsection*{\S~2. --- Homeomorphism.}

Let us consider a substitution that changes $x_{1}, x_{2}, \ldots, x_{n}$ into $x_{1}, x'_{2}, \ldots, x'_{n}$\ednote{Printed $x_{1}$ without a prime; expected $x'_{1}$. The prime may be a misprint or lost ink.}, which I subject only to the following conditions.

We have
\[
x'_{i} = \varphi_{i}(x_{1}, x_{2}, \ldots, x_{n}) \qquad (i = 1, 2, \ldots, n).
\tag{4}
\]

In a certain domain, the functions $\varphi_{i}$ are uniform, finite and continuous; they have continuous derivatives, and their functional determinant does not vanish.

If one solves the equations (4) with respect to $x_{1}, x_{2}, \ldots, x_{n}$, there results
\[
x_{k} = \varphi'_{k}(x'_{1}, x'_{2}, \ldots, x'_{n}) \qquad (k = 1, 2, \ldots, n)
\]
and the functions $\varphi'_{k}$ satisfy the same conditions as the functions $\varphi_{i}$.

It is clear that the set of substitutions that satisfy these conditions forms a group, and this group is one of the most general that one can imagine. The Science whose object is the study of this group and of some others analogous to it has received the name of \emph{Analysis Situs}.

It is clear that a substitution of this group will transform a manifold of $m$ dimensions into a manifold of $m$ dimensions, and that the transformed manifold will be continuous, or finite, or unlimited if the given manifold is so itself (and conversely).

Let there be two manifolds of the same number of dimensions, $V$ and $V'$, defined respectively by the conditions
\[
\left\{
\begin{aligned}
&F_{\alpha} = 0 &&(\alpha = 1, 2, \ldots, p),\\
&\varphi_{\beta} > 0 &&(\beta = 1, 2, \ldots, q),
\end{aligned}
\right.
\tag{1}
\]
\origpage{8}
and
\[
\left\{
\begin{aligned}
&F'_{\alpha} = 0 &&(\alpha = 1, 2, \ldots, p),\\
&\varphi'_{\beta} > 0 &&(\beta = 1, 2, \ldots, q').
\end{aligned}
\right.
\tag{1 \textit{bis}}
\]

Let us suppose that one can make correspond to a point $x_{1}, x_{2}, \ldots, x_{n}$ of the manifold $V$ a point $x'_{1}, x'_{2}, \ldots, x'_{n}$ of the manifold $V'$, in such a way that one has
\[
x'_{k} = \psi_{k}(x_{1}, x_{2}, \ldots, x_{n}) \qquad (k = 1, 2, \ldots, n).
\tag{5}
\]

I consider the domain $D$ defined by the inequalities
\[
F_{\alpha} > -\varepsilon, \qquad F_{\alpha} < \varepsilon, \qquad \varphi_{\beta} > 0.
\]

The manifold $V$ is evidently contained entirely in the domain $D$.

I suppose that \emph{in the domain} $D$ the functions $\psi_{k}$ are finite, continuous and uniform, that they have continuous derivatives and that their functional determinant is never zero.

Let us now solve the equations (5); we shall find
\[
x_{k} = \psi'_{k}(x'_{1}, x'_{2}, \ldots, x'_{n}) \qquad (k = 1, 2, \ldots, n).
\tag{6}
\]
\ednote{In the parenthesis of (6) the scan shows only a single bar where $=$ is expected, apparently the lower stroke of $=$ with the upper stroke unprinted.}

I consider the domain $D'$ defined by the inequalities
\[
F'_{\alpha} > -\varepsilon, \qquad F'_{\alpha} < \varepsilon, \qquad \varphi'_{\beta} > 0,
\]
and I suppose that \emph{in the domain} $D'$ the functions $\psi'_{k}$ are finite, continuous and uniform, that they have continuous derivatives and that their functional determinant is not zero.

It follows from these hypotheses that to every point of $V$ there corresponds one point of $V'$ and only one, and conversely; to every manifold $W$ contained in $V$ there will correspond a manifold $W'$, of the same number of dimensions, contained in $V'$; if $W$ is continuous, finite or unlimited, the same will be true of $W'$, and conversely.
\origpage{9}

If all these conditions are fulfilled, we shall say that the two manifolds $V$ and $V'$ are equivalent from the point of view of \emph{Analysis Situs}, or, to shorten the language, that they are \emph{homeomorphic}, that is to say of like form.

I may also say that two more complicated figures, composed of any number of manifolds, are homeomorphic when one passes from the one to the other by a transformation of the form (5).
Thus two polygons with the same number of sides will be homeomorphic; two polyhedra whose faces are equal in number, have the same number of sides and are similarly arranged, will be homeomorphic, etc.

\subsection*{\S~3. --- Second definition of manifolds.}

Manifolds can be defined in an entirely different way. Let us consider $n$ equations
\[
\left\{
\begin{aligned}
&x_{1} = \theta_{1}(y_{1}, y_{2}, \ldots, y_{m}),\\
&x_{2} = \theta_{2}(y_{1}, y_{2}, \ldots, y_{m}),\\
&\ldots\ldots\ldots\ldots\ldots\ldots,\\
&x_{n} = \theta_{n}(y_{1}, y_{2}, \ldots, y_{m}).
\end{aligned}
\right.
\tag{8}
\]

It is clear that these equations (if the $y$ are regarded as independent variables) represent a manifold of $m$ dimensions.

One will still have a manifold of $m$ dimensions if one adjoins to the equations (8) a certain number of inequalities of the form
\[
\psi(y_{1}, y_{2}, \ldots, y_{m}) > 0,
\]
which limit the field of variation of the variables $y$.

I shall suppose that the functions $\theta$ are finite and continuous; but I can make one more hypothesis without restricting the generality in an essential way (this is, moreover, what I could equally well have done with my first definition of manifolds).
\origpage{10}

I can suppose, I say, that the functions $\theta$ are analytic. If, indeed, these functions $\theta$ are finite and continuous, one will be able to find functions $\theta'$ which will be analytic and which will differ from the $\theta$ as little as we please.

I shall suppose, moreover, that the functional determinants of $m$ of the functions $\theta$ with respect to the $y$ never all vanish at once.

One will obtain a manifold which will not differ from the first by replacing the $y$ in the equations (8) by any $m$ analytic functions of $m$ variables $z_{1}, z_{2}, \ldots, z_{m}$.

The scope of this new definition would be rather restricted if one could not enlarge it by the process of \emph{analytic continuation}.

Let us consider two manifolds $V$ and $V'$ defined by equations analogous to (8). Let there be, for example, the equations
\[
x_{i} = \theta_{i}(y_{1}, y_{2}, \ldots, y_{m})
\tag{8}
\]
for $V$, and the equations
\[
x_{i} = \theta'_{i}(y'_{1}, y'_{2}, \ldots, y'_{m})
\tag{8 \textit{bis}}
\]
\ednote{In (8 bis) the scan shows only a single bar at mid-height where $=$ is expected, apparently one stroke of $=$ with the other unprinted, as in (6) on p. 8.}
for $V'$.

It may happen that these two manifolds have a common part $V''$ also having $m$ dimensions; that is to say, all the points of $V''$ will belong at once to $V$ and to $V'$.

In that case, in the interior of $V''$, the $y$ will be analytic functions of the $y'$ and conversely.

One will then say that the two manifolds $V$ and $V'$ are the \emph{analytic continuation} of one another.

One can thus form a chain of manifolds
\[
V_{1}, \qquad V_{2}, \qquad \ldots, \qquad V_{n},
\]
in such a way that each of them is the analytic continuation of the
\origpage{11}
preceding one and that between two consecutive manifolds of the chain there is a common part. This is what I shall call a \emph{continuous chain}.

It may also happen that the chain is closed, that is to say that $V_{n}$ does not differ from $V_{1}$.

One can also have a network of manifolds, that is to say a set of manifolds such that each of them is the continuation of several others and that one can pass from any one of them to any other of them by analytic continuation; this is what I shall call a \emph{continuous network}.

One will then be able to consider the set of all the manifolds of one chain or of one network as forming a single manifold.

This is a definition wider still than the first.

There are, indeed, manifolds (and we shall see examples of them further on) which can be decomposed into a certain number of partial manifolds, forming a continuous chain or network, and such that each of them can be defined by equations of the form (8) (manifolds which, consequently, come under our second definition), and which nevertheless cannot be defined by relations of the form (1) and do not come under our first definition.

On the contrary, every manifold that satisfies the first definition also satisfies the second.

We can indeed, by a well-known theorem, if $y_{1}, y_{2}, \ldots, y_{n}$ are defined by $n$ relations of the form
\[
\left\{
\begin{aligned}
&y_{1} = F_{1}(x_{1}, x_{2}, \ldots, x_{n}),\\
&y_{2} = F_{2}(x_{1}, x_{2}, \ldots, x_{n}),\\
&\ldots\ldots\ldots\ldots\ldots\ldots,\\
&y_{n} = F_{n}(x_{1}, x_{2}, \ldots, x_{n}),
\end{aligned}
\right.
\tag{$\alpha$}
\]
if, in a certain domain, the $F$ are holomorphic functions of the $x$ whose functional determinant does not vanish; we can, I say, solve these equations and derive from them
\[
x_{i} = \theta_{i}(y_{1}, y_{2}, \ldots, y_{n}),
\tag{$\beta$}
\]
\origpage{12}
the $\theta_{i}$ being functions of the $y$, holomorphic in a certain domain.

Now let there be a manifold $V$ satisfying our first definition, that is to say defined by the relations
\[
F_{\alpha} = 0, \qquad \varphi_{\beta} > 0.
\tag{1}
\]

Let us then take up the equations ($\alpha$) again; let us take in them, for $F_{1}, F_{2}, \ldots, F_{p}$, the first members of the $p$ equalities (1). As for the other functions
\[
F_{p+1}, \qquad F_{p+2}, \qquad \ldots, \qquad F_{n},
\]
they will be any $n - p$ holomorphic functions of the $x$. I shall subject them to one condition only.

Let $x^{0}_{1}, x^{0}_{2}, \ldots, x^{0}_{n}$ be any point $M_{0}$ of the manifold $V$. I shall arrange things in such a way that the functional determinant of the $n$ functions $F$ does not vanish for
\[
x_{i} = x^{0}_{i}.
\]

This is evidently possible, since I have supposed that the functional determinants of the $p$ functions
\[
F_{1}, \qquad F_{2}, \qquad \ldots, \qquad F_{p},
\]
with respect to any $p$ of the variables $x$, never all vanish at once.

I can suppose further that $F_{p+1}, F_{p+2}, \ldots, F_{n}$ vanish at the point $M_{0}$, that is to say for
\[
x_{i} = x^{0}_{i}.
\]

Then, by the theorem cited above, one can solve the equations ($\alpha$), and one finds that the $x_{i}$ are expressed in series ordered according to the powers of
\[
y_{1}, \qquad y_{2}, \qquad \ldots, \qquad y_{n},
\]
and convergent provided that these quantities satisfy certain inequalities.
\origpage{13}

Let, then,
\[
x_{i} = \theta_{i}
\tag{$\beta$}
\]
be these equations, and
\[
\lambda_{k}(y_{1}, y_{2}, \ldots, y_{n}) > 0
\]
the conditions that the $y$ must satisfy for the series to be convergent.

Let us now put
\[
y_{1} = y_{2} = \ldots = y_{p} = 0;
\]
the first $p$ equations ($\alpha$) will no longer differ from the $p$ equations (1); and the functions $\theta_{i}$, now depending only on $n - p = m$ variables, will be of the form (8).

Then the set of relations
\[
x_{i} = \theta_{i}, \qquad \varphi_{\beta} > 0, \qquad \lambda_{k} > 0
\]
will represent a manifold $v$, which will be defined in the second manner and which, not differing from
\[
F_{\alpha} = 0, \qquad \varphi_{\beta} > 0, \qquad \lambda_{k} > 0,
\]
will form part of $V$.

The point $M_{0}$, which is an \emph{arbitrary} point of $V$, forms part of $v$. One can therefore construct, around any point of $V$, a manifold analogous to $v$.

The simplest case is that in which the conditions of convergence $\lambda_{k} > 0$ are a consequence of the inequalities $\varphi_{\beta} > 0$.

Then $v$ does not differ from $V$, and one can be content, to define $v$, with the equalities (8) and the inequalities
\[
\psi_{\beta} > 0,
\tag{9}
\]
\origpage{14}
which one obtains by substituting the $\theta$ in place of the $x$ in the inequalities
\[
\varphi_{\beta} > 0.
\]

Let us remark in passing that the functional determinants of $m$ of the functions $\theta$ with respect to the $y$ will never all vanish at once.

If the conditions $\lambda_{\beta} > 0$\ednote{Printed $\lambda_{\beta}$; the convergence conditions are written $\lambda_{k}$ on p. 13, so $\lambda_{k}$ is apparently meant.} were not a consequence of the inequalities $\varphi_{\beta} > 0$, one would divide the manifold $V$ into partial manifolds; thus, for example, the manifold $V$ can evidently be decomposed into two partial manifolds
\[
F_{\alpha} = 0, \qquad \varphi_{\beta} > 0, \qquad \psi > 0
\]
and
\[
F_{\alpha} = 0, \qquad \varphi_{\beta} > 0, \qquad \psi < 0,
\]
$\psi$ being any function of $x_{1}, x_{2}, \ldots, x_{n}$.

This being so, we shall always be able to divide $V$ into partial manifolds, or better to construct on $V$ a certain number of partial manifolds, overlapping one another and small enough that one can, for each of them, find a system of auxiliary variables which allows this partial manifold to be represented by equalities and inequalities of the form (8) and (9), satisfying all the conditions I have just stated. Every point $M_{0}$ of $V$ will belong to one of these partial manifolds, and the set of these manifolds will form a continuous network. The first definition is thus reduced to the second.

Yet there remains a remark to be made; it may happen that to two different systems of values of $y_{1}, y_{2}, \ldots, y_{m}$ there corresponds one and the same system of values of the functions $\theta_{1}, \theta_{2}, \ldots, \theta_{n}$ and, consequently, one and the same point of the manifold $V$.

In that case it is appropriate to adjoin to the inequalities (9) other inequalities
\[
\psi_{\gamma} > 0,
\tag{10}
\]
choosing them in such a way that, of the various systems of values of the
\origpage{15}
variables $y$ which correspond to one and the same point of $V$, there is always one and only one that satisfies the inequalities (9) and (10).

Let there be, for example, a torus, whose equation will be
\[
(x^{2}_{1} + x^{2}_{2} + x^{2}_{3} + R^{2} - r^{2})^{2} - 4R^{2}(x^{2}_{1} + x^{2}_{2}) = 0.
\]
\ednote{The minus sign before $r^{2}$ prints as broken short strokes; read as a single $-$.}

Let us put
\[
\begin{aligned}
&x_{1} = (R + r\cos y_{1})\cos y_{2},\\
&x_{2} = (R + r\cos y_{1})\sin y_{2},\\
&x_{3} = r\sin y_{1}.
\end{aligned}
\]

We see that to one and the same point of the torus there will correspond an infinity of systems of values of the $y$, comprised in the formula
\[
y_{1} + 2K_{1}\pi, \qquad y_{2} + 2K_{2}\pi,
\]
where the $K$ are integers.

But if we subject the $y$ to the conditions
\[
0 < y_{1} < 2\pi, \qquad 0 < y_{2} < 2\pi,
\]
we shall have only one system of values of the $y$ that corresponds to a point of the torus.

\subsection*{\S~4. --- Opposite manifolds.}

With the one definition as with the other, there is occasion to make a distinction whose importance will be understood later.

Let us first suppose a manifold $V$ defined in the first manner, that is to say by the relations
\[
F_{\alpha} = 0, \qquad \varphi_{\beta} > 0.
\]

One will take account of the order in which the equations $F_{\alpha} = 0$ are arranged; if two of these equations are permuted, we shall agree to say
\origpage{16}
that the system of relations no longer represents the manifold $V$, but a \emph{manifold opposite} to $V$.

We can now replace the equations
\[
F_{\alpha} = 0 \qquad (\alpha = 1, 2, \ldots, p)
\]
by the following:
\[
\begin{aligned}
&\Phi_{1} = A_{1.1}F_{1} + A_{1.2}F_{2} + \ldots + A_{1.p}F_{p} = 0,\\
&\Phi_{2} = A_{2.1}F_{1} + A_{2.2}F_{2} + \ldots + A_{2.p}F_{p} = 0,\\
&\ldots\ldots\ldots\ldots\ldots\ldots,\\
&\Phi_{p} = A_{p.1}F_{1} + A_{p.2}F_{2} + \ldots + A_{p.p}F_{p} = 0,
\end{aligned}
\]
the $A_{i.k}$ being any functions of the $x$.

If the determinant $\Delta$ of the coefficients $A_{i.k}$ does not vanish in the domain considered, the equations $\Phi_{\alpha} = 0$ will be equivalent to the equations $F_{\alpha} = 0$, and consequently, if one adjoins to them a certain number of inequalities, they will still represent the manifold $V$ or an opposite manifold.

We shall agree to say that if $\Delta$ (which always has the same sign, since I have supposed that it does not vanish) is positive, these equations represent the manifold $V$, and that if $\Delta$ is negative, they represent the opposite manifold.

Likewise with the second definition. Let $V$ be a manifold defined by the equations
\[
x_{i} = \theta_{i}(y_{1}, y_{2}, \ldots, y_{m}).
\]

One will have to take account of the order of the variables $y$ and, if two of these variables are permuted, we shall agree to say that the equations no longer represent $V$, but the opposite manifold.

Let us imagine that the $y$ are replaced by $m$ analytic functions of the $m$ new variables $z_{1}, z_{2}, \ldots, z_{m}$, in such a way that one has
\[
x_{i} = \theta'_{i}(z_{1}, z_{2}, \ldots, z_{m}),
\]
these new equations will still represent $V$ or the opposite manifold.
\origpage{17}

We must suppose that the functional determinant $\Delta$ of the $y$ with respect to the $z$ does not vanish, so that to one system of values of the $y$ there corresponds a single system of values of the $z$. It will therefore always be of the same sign.

We shall agree to say that, if $\Delta$ is positive, the new equations still represent $V$, and that if $\Delta$ is negative, they represent the opposite manifold.

Let us now see what happens when one passes from one definition to the other.

Let there be a manifold $V$ defined by
\[
F_{1} = F_{2} = \ldots = F_{p} = 0,
\]
and by certain inequalities.

Let us adjoin to these $p$ equations the following
\[
y_{1} = F_{p+1}, \qquad y_{2} = F_{p+2}, \qquad \ldots, \qquad y_{n-p} = F_{n},
\]
$F_{p+1}, F_{p+2}, \ldots, F_{n}$ being any $n - p$ functions of the $x$.

We have seen that if the functional determinant $\Delta$ of the $n$ functions $F_{1}, F_{2}, \ldots, F_{n}$ is not zero, one can solve these $n$ equations with respect to the $x$, and that one thus finds $n$ equations
\[
x_{i} = \theta_{i}(y_{1}, y_{2}, \ldots, y_{n-p}),
\]
which represent a manifold of $n - p$ dimensions. But one may ask whether they will represent $V$ or the opposite manifold.

We shall agree to say that they represent $V$ if $\Delta$ is positive and the opposite manifold if $\Delta$ is negative.

Let us consider the manifold of $n - p$ dimensions
\[
F_{\alpha} = 0, \qquad \ldots, \qquad \varphi_{\beta} > 0,
\]
which I call $V$. We have seen that the manifold of $n - p - 1$ dimensions
\origpage{18}
\[
F_{\alpha} = 0, \qquad \varphi_{\gamma} = 0, \qquad \varphi_{\beta} > 0, \qquad (\beta \gtrless \gamma),
\]
which I call $v$, forms part of the boundary of $V$.

But it is important to arrange the equations that define $v$ in the following order:
\[
F_{1} = 0, \qquad F_{2} = 0, \qquad \ldots, \qquad F_{p} = 0, \qquad \varphi_{\gamma} = 0,
\]
for if two of them were permuted, we should agree to say that these equations no longer represent the boundary of $V$ or a part of this boundary, but a manifold opposite to this boundary.

\subsection*{\S~5. --- Homologies.}

Let us consider a manifold $V$ of $p$ dimensions; now let $W$ be a manifold of $q$ dimensions ($q \leqq p$) forming part of $V$. Let us suppose that the complete boundary of $W$ consists of $\lambda$ continuous manifolds of $q - 1$ dimensions
\[
v_{1}, \qquad v_{2}, \qquad \ldots, \qquad v_{\lambda}.
\]

We shall express this fact by the notation
\[
v_{1} + v_{2} + \ldots + v_{\lambda} \backsim 0.
\]

More generally the notation
\[
k_{1}v_{1} + k_{2}v_{2} \backsim k_{3}v_{3} + k_{4}v_{4},
\]
where the $k$ are integers and the $v$ manifolds of $q - 1$ dimensions, will signify that there exists a manifold $W$ of $q$ dimensions forming part of $V$ whose complete boundary will consist of $k_{1}$ manifolds little different from $v_{1}$, of $k_{2}$ manifolds little different from $v_{2}$, of $k_{3}$ manifolds little different from the manifold
\origpage{19}
opposite to $v_{3}$, and of $k_{4}$ manifolds little different from the manifold opposite to $v_{4}$.

Relations of this form may be called \emph{homologies}.

Homologies can be combined like ordinary equations. We shall also employ the following notation; let us suppose that one has
\[
k_{1}v_{1} + k_{2}v_{2} + \ldots + k_{p}v_{p} \backsim w_{1} + w_{2} + \ldots + w_{q},
\]
and that the manifolds $w_{1}, w_{2}, \ldots, w_{q}$ form part of the boundary of $V$; we shall sometimes write
\[
k_{1}v_{1} + k_{2}v_{2} + k_{p}v_{p} \backsim \varepsilon.
\]
\ednote{The terms $+ \ldots$ between $k_{2}v_{2}$ and $k_{p}v_{p}$, present in the preceding homology, are not printed here; an apparent misprint.}

\subsection*{\S~6. --- Betti numbers.}

We shall say that the manifolds
\[
v_{1}, \qquad v_{2}, \qquad \ldots, \qquad v_{\lambda},
\]
of the same number of dimensions and forming part of $V$, are \emph{linearly independent} if they are bound by no homology with integer coefficients.

If there exist $P_{m} - 1$ \emph{closed} manifolds of $m$ dimensions forming part of $V$ and linearly independent, and if there exist only $P_{m} - 1$ of them, we shall say that the order of connection of $V$ with respect to manifolds of $m$ dimensions is equal to $P_{m}$.

Thus there are defined, as regards a manifold $V$ of $m$ dimensions, $m - 1$ numbers which I shall call
\[
P_{1}, \qquad P_{2}, \qquad \ldots, \qquad P_{m-1},
\]
and which are the orders of connection of $V$ with respect to manifolds of
\[
1, \qquad 2, \qquad \ldots, \qquad m - 1
\]
dimensions.
\origpage{20}

I shall call them in what follows the \emph{Betti numbers}.

Let us make these definitions clearer by an example:

Let $D$ be a domain forming part of ordinary space and bounded by $n$ closed surfaces
\[
S_{1}, \qquad S_{2}, \qquad \ldots, \qquad S_{n},
\]
which do not intersect.

This domain is a manifold of three dimensions. It therefore admits two Betti numbers $P_{1}$ and $P_{2}$.

This manifold is then defined by the inequalities
\[
\varphi_{1} > 0, \qquad \varphi_{2} > 0, \qquad \ldots, \qquad \varphi_{n} > 0,
\]
if the equations of the $n$ surfaces $S$ are
\[
\varphi_{1} = 0, \qquad \varphi_{2} = 0, \qquad \ldots, \qquad \varphi_{n} = 0.
\]

Since the surfaces do not intersect, there are no values of $x_{1}, x_{2}, x_{3}$ that satisfy two of these equations at once,
\[
\varphi_{i} = 0, \qquad \varphi_{k} = 0.
\]

The surfaces $S_{1}, S_{2}, \ldots, S_{n}$, having two dimensions, will have only a single Betti number, which will be Riemann's order of connection; let
\[
2Q_{1} + 1, \qquad 2Q_{2} + 1, \qquad \ldots, \qquad 2Q_{n} + 1,
\]
be the orders of connection (which are odd, since the surfaces are closed) of the $n$ surfaces
\[
S_{1}, \qquad S_{2}, \qquad \ldots, \qquad S_{n},
\]
then one will have
\[
P_{2} = n, \qquad P_{1} = Q_{1} + Q_{2} + \ldots + Q_{n} + 1.
\]

Thus, for the region interior to a sphere
\[
P_{2} = 1, \qquad P_{1} = 1;
\]
\origpage{21}
for the region comprised between two spheres
\[
P_{2} = 2, \qquad P_{1} = 1;
\]
for the region interior to a torus
\[
P_{2} = 1, \qquad P_{1} = 2,
\]
for the region comprised between two tori
\[
P_{2} = 2, \qquad P_{1} = 2.
\]

\subsection*{\S~7. --- Use of integrals.}

Let us consider a manifold $V$ which can be represented by the inequalities and equalities (8), (9) and (10) in such a way as to satisfy all the conditions stated above.

One then knows what is to be understood by the multiple integral of order $m$
\[
\int F\,dy_{1}, dy_{2}, \ldots, dy_{m}
\]
extended over the manifold $V$; I denote, of course, by $F$ a given function of the $y$. The integration must be carried out successively with respect to the $m$ variables $y$, and the limits of integration will be defined by the inequalities (9) and (10).

This being so, I am going to define the following integral:
\[
\int \sum X_{\alpha_{1}\alpha_{2}\ldots\alpha_{m}}\,dx_{\alpha_{1}}\,dx_{\alpha_{2}} \ldots dx_{\alpha_{m}}. \tag{11}
\]

The differentials $dx_{\alpha_{1}}, dx_{\alpha_{2}}, \ldots, dx_{\alpha_{m}}$ are any $m$ of the $n$ differentials $dx_{1}, dx_{2}, \ldots, dx_{n}$. The functions $X_{\alpha_{1}\alpha_{2}\ldots\alpha_{m}}$ are given functions of $x_{1}, x_{2}, \ldots, x_{n}$, and there are as many of them as there are possible combinations of the indices
\[
\alpha_{1}, \qquad \alpha_{2}, \qquad \ldots, \qquad \alpha_{m},
\]
\origpage{22}
that is to say as there are combinations of $n$ letters $m$ by $m$. One must agree that the function $X$ is zero if two of its indices are equal and that it changes sign when two of its indices are permuted.

This being so, the integral (11) will by definition be equal to the integral of order $m$
\[
\int \sum X_{\alpha_{1}\alpha_{2}\ldots\alpha_{m}} \frac{\partial(x_{\alpha_{1}}, x_{\alpha_{2}}, \ldots, x_{\alpha_{m}})}{\partial(y_{1}, y_{2}, \ldots, y_{m})}\,dy_{1}, dy_{2}, \ldots, dy_{m}.
\]

If now the manifold $V$ were not susceptible of being represented by relations of the form (8), (9) and (10) satisfying all the stated conditions, one would decompose the manifold $V$ into partial manifolds small enough to be susceptible of this mode of representation, and the integral (11), extended over the total manifold $V$, would by definition be the sum of the integrals (11) extended over the various partial manifolds.

This definition nevertheless still leaves an ambiguity.

Indeed, if one permutes two of the letters $y_{1}$ and $y_{2}$, the integral changes sign; it is therefore important to be given the order of these letters, and the permutation of two of these letters would be equivalent to a change of the sense of integration in the study of simple integrals. I shall therefore say the sense of integration when speaking of the order in which one agrees to arrange the letters $y_{1}, y_{2}, \ldots, y_{m}$.

I have had occasion to concern myself with an analogous question in a Memoir \emph{sur les résidus des intégrales doubles}, published in Volume IX of the \emph{Acta mathematica}, and in particular in \S~3 of that Memoir, entitled: \emph{Conditions d'intégrabilité}.

I sought in which cases these \emph{conditions of integrability} are fulfilled, that is to say in which cases the integral (11) is zero every time it is applied to a closed manifold.

Here is what I found; let us write, to shorten the writing,
\[
(\alpha_{1}, \alpha_{2}, \ldots, \alpha_{m})
\]
instead of $X_{\alpha_{1}, \alpha_{2}, \ldots, \alpha_{m}}$ and $[\alpha_{p}]$ instead of $x_{\alpha_{p}}$.
\origpage{23}

Our conditions of integrability will be written
\[
\left\{
\begin{aligned}
&\frac{d(\alpha_{1}, \alpha_{2}, \ldots, \alpha_{m})}{d[\alpha_{m+1}]} \pm \frac{d(\alpha^{2}, \alpha^{3}, \ldots, \alpha_{m-1})}{d[\alpha_{1}]}\\
&\pm \frac{d(\alpha_{3}, \alpha_{1}, \ldots, \alpha_{m}, \alpha_{m+1}, \alpha_{1})}{d[\alpha_{2}]} \pm \ldots \pm \frac{d(\alpha_{m+1}, \alpha_{1}, \alpha_{2}, \ldots, \alpha_{m-1})}{d[\alpha_{m}]} = 0.
\end{aligned}
\right. \tag{12}
\]
\ednote{In the second term of (12) the print sets $\alpha^{2}, \alpha^{3}, \ldots, \alpha_{m-1}$, with superscripts and a final $\alpha_{m-1}$, where the cyclic pattern requires $\alpha_{2}, \alpha_{3}, \ldots, \alpha_{m+1}$; and the second index of the third term reads as $\alpha_{1}$ where $\alpha_{4}$ is expected, possibly a damaged 4. Apparent misprints, reproduced as printed.}

Here is the law according to which the signs $\pm$ must be chosen. One will always take the sign $+$ if $m$ is even, and alternately the sign $+$ and the sign $-$ if $m$ is odd.

There will be as many equations (12) as there are possible systems of indices
\[
\alpha_{1}, \qquad \alpha_{2}, \qquad \ldots, \qquad \alpha_{m}, \qquad \alpha_{m+1},
\]
that is to say, since these indices must be chosen among the letters
\[
1, \qquad 2, \qquad \ldots, \qquad n,
\]
as many as there are combinations of $n$ letters $m + 1$ by $m + 1$.

Let us now suppose that the conditions (12), instead of being satisfied for all possible values of the $n$ variables
\[
x_{1}, \qquad x_{2}, \qquad \ldots, \qquad x_{n},
\]
are satisfied only for certain values of these variables. For example, let us consider a manifold $V$ defined by the conditions
\[
F_{\alpha} = 0, \qquad \varphi_{\beta} > 0.
\]
Next let there be a domain $D$ comprising all the points neighbouring those of $V$ and defined, for example, by the conditions
\[
-\varepsilon < F_{\alpha} < \varepsilon, \qquad \varphi_{\beta} > -\varepsilon,
\]
$\varepsilon$ being a very small positive number.
\origpage{24}

Let us suppose that the conditions (12) are satisfied for all the points of the domain $D$.

By repeating the reasoning of the Memoir cited and modifying it suitably, one will prove the following: Let there be a manifold $V'$ of $m + 1$ dimensions forming part of $V$ (the number $m + 1$ must therefore be less than or at most equal to the number of dimensions of $V$).

Let us suppose that the complete boundary of $V'$ consists of the $k$ manifolds of $m$ dimensions
\[
W_{1}, \qquad W_{2}, \qquad \ldots, \qquad W_{k},
\]
so that $W_{1} + W_{2} + \ldots + W_{k} \backsim 0$.

Then, if the integral (11) satisfies the conditions (12) in the domain $D$, the algebraic sum of the integrals (11) extended over the manifolds $W_{1}, W_{2}, \ldots, W_{k}$ is zero. One must, of course, pay attention for each of them to the sense of integration.

The conditions (12) are sufficient for this to be so, but they are not necessary; these conditions, as we have seen, are equal in number to the combinations of $n$ letters $m + 1$ by $m + 1$; for the result I have just stated to remain exact, it would suffice that the integral (11) satisfy, for all the points of $V$, certain conditions equal in number to the combinations of $n - p$ letters $m + 1$ by $m + 1$, $n - p$ being the number of dimensions of $V$.

These conditions would be easy to form, but that would lead me too far from my subject.

If, then, an integral (11) satisfies the conditions (12) in the domain $D$, the various values that this integral can take when it is extended over various closed manifolds of $m$ dimensions forming part of $V$ will be linear combinations with integer coefficients of a certain number of them, which may be called the periods of the integral (11).

The maximum number of periods is equal to $P_{m} - 1$; for, if one considers any $P_{m}$ closed manifolds of $m$ dimensions, there will always be a manifold of $m + 1$ dimensions that admits as complete boundary
\origpage{25}
these $P_{m}$ manifolds or some of them. There will therefore always be, between the $P_{m}$ corresponding integrals, a linear relation with integer coefficients. One could moreover show that there always exist integrals of the form (11) for which the maximum number of periods is attained. This way of making the definition of the Betti numbers understood was employed by Betti himself for the first and the last of these numbers, that is to say for $P_{1}$ and $P_{m-1}$; but we have just seen that it is easy to do the same for the other Betti numbers.

\subsection*{\S~8. --- One-sided and two-sided manifolds.}

Let us consider a manifold $V$ defined in the second manner, that is to say formed of a chain or a network of partial manifolds, each of which is itself defined by relations of the form (8) and (9).

Let $v_{1}$ be a partial manifold defined by the conditions
\[
\begin{aligned}
&x_{i} = \theta_{i}(y_{1}, y_{2}, \ldots, y_{m}),\\
&|y_{k}| < \beta_{k}.
\end{aligned}
\]

Let $v_{2}$ be another partial manifold, defined by the conditions
\[
\begin{aligned}
&x_{i} = \theta'_{i}(z_{1}, z_{2}, \ldots, z_{m}),\\
&|z_{k}| < \gamma_{k}.
\end{aligned}
\]

Let us suppose that these two manifolds have a common part $v'$ forming a continuous manifold. I say that in the interior of this manifold the functional determinant
\[
\Delta = \frac{\partial(y_{1}, y_{2}, \ldots, y_{m})}{\partial(z_{1}, z_{2}, \ldots, z_{m})}
\]
is always of the same sign.

Indeed, it could not change sign without vanishing or without becoming
\origpage{26}
infinite. Now one has
\[
\Delta = \frac{\partial(x_{1}, x_{2}, \ldots, x_{m})}{\partial(z_{1}, z_{2}, \ldots, z_{m})} : \frac{\partial(x_{1}, x_{2}, \ldots, x_{m})}{\partial(y_{1}, y_{2}, \ldots, y_{m})},
\]
so that $\Delta$ is itself the ratio of two functional determinants; since these two functional determinants are essentially finite, $\Delta$ could vanish only if one had
\[
\frac{\partial(x_{1}, x_{2}, \ldots, x_{m})}{\partial(z_{1}, z_{2}, \ldots, z_{m})} = 0,
\]
and since nothing distinguishes the first $m$ variables $x_{1}, x_{2}, \ldots, x_{m}$ from the $n - m$ others $x_{m+1}, x_{m+2}, \ldots, x_{n}$, the functional determinant of any $m$ of the $x$ with respect to the $z$ would have to be zero.

All the functional determinants of the functions $\theta'$ would therefore have to vanish at once, which by hypothesis never happens. $\Delta$ therefore cannot vanish, and one would see in absolutely the same way that it cannot become infinite either.

$\Delta$ is therefore always of the same sign, and we shall be able to choose the order of the variables $z$ in such a way that this sign is always positive.

A difficulty might present itself in certain cases; let us suppose that the part common to $v_{1}$ and $v_{2}$, instead of reducing to a single continuous manifold $v'$, consists of several continuous manifolds $v'$, $v''$, $v'''$; in each of them the sign of $\Delta$ will remain constant, but it may change when one passes from one to another. In this case we shall say that the manifold $V$ is \emph{one-sided}.

Let us suppose that this circumstance does not present itself, and let us consider a sequence of partial manifolds forming a closed chain
\[
v_{1}, \qquad v_{2}, \qquad \ldots, \qquad v_{q}, \qquad v_{1}.
\]
Let
\[
y^{i}_{1}, \qquad y^{i}_{2}, \qquad \ldots, \qquad y^{i}_{m}
\]
be the $m$ variables that play the part of $y_{1}, y_{2}, \ldots, y_{m}$ with respect to $v_{i}$.
\origpage{27}

Let us suppose that $v_{1}$ and $v_{2}$ have a common part $v'_{1}$, $v_{2}$ and $v_{3}$ a common part $v'_{2}$, \ldots, $v_{q-1}$ and $v_{q}$ a common part $v'_{q-1}$, and finally $v_{q}$ and $v_{1}$ a common part $v'_{q}$.

Let $\Delta_{i}$ be the functional determinant of
\[
y^{i+1}_{1}, \qquad y^{i+2}_{2}, \qquad \ldots, \qquad y^{i+1}_{m}
\]\ednote{Printed $y^{i+2}_{2}$ instead of $y^{i+1}_{2}$; an apparent misprint.}
with respect to
\[
y^{i}_{1}, \qquad y^{i}_{2}, \qquad \ldots, \qquad y_{m}.
\]\ednote{Printed $y_{m}$ without the superscript; $y^{i}_{m}$ is expected. The superscript may have failed to print.}

This determinant will be defined in the interior of $v'_{i}$.

I define likewise in the interior of $v'_{q}$ the functional determinant $\Delta_{q}$ of
\[
y^{1}_{1}, \qquad y^{1}_{2}, \qquad \ldots, \qquad y^{1}_{m}
\]
with respect to
\[
y^{q}_{1}, \qquad q^{q}_{2}, \qquad \ldots, \qquad y^{q}_{m}.
\]\ednote{Printed $q^{q}_{2}$ instead of $y^{q}_{2}$; an apparent misprint.}

By what we have just seen, we can always choose the order of the variables in such a way that
\[
\Delta_{1}, \qquad \Delta_{2}, \qquad \ldots, \qquad \Delta_{q-1}
\]
are always positive. On the other hand, $\Delta_{q}$ is always of the same sign, but is this constant sign the sign $+$ or the sign $-$?

If it is the sign $-$, I shall again say that the manifold $V$ is \emph{one-sided}. I may also say that the manifolds $v_{1}, v_{2}, \ldots, v_{q}$ form a one-sided chain.

Let us now suppose that one has constructed a certain continuous network of partial manifolds
\[
v_{1}, \qquad v_{2}, \qquad \ldots, \qquad v_{q}, \tag{4}
\]
in such a way that every point of $V$ is in the interior (I exclude the boundary) of one of the manifolds (4) or of several of them. If, in the part common to any two of the manifolds (4), the determinant $\Delta$ is positive, I shall say that the manifold $V$ is \emph{two-sided}. If it were not so, it is clear
\origpage{28}
that one could always, with some of the manifolds (4), form a one-sided chain, and that the manifold $V$ would be one-sided. I may also say that the network of the manifolds (4) forms a two-sided system.

But to justify our definition completely, it must be shown that $V$ cannot be at once one-sided and two-sided. It is clear first that in the two-sided system (4) one cannot find a one-sided chain.

But it remains to show that the system (4) will remain one-sided\ednote{Printed ``unilatère''; apparently for ``bilatère'', which is what the proof on p. 29 concludes.} when one adjoins to it any manifold $v_{q+1}$ forming part of $V$.

Let $v'_{i}$ be the part common to $v_{i}$ and to $v_{q+1}$; every point of $v_{q+1}$ will belong to at least one of the manifolds $v'_{i}$, and since $v_{q+1}$ is continuous, if I consider two points $M_{1}$ and $M_{k}$ of $v_{q+1}$ belonging respectively to $v'_{1}$ and to $v'_{k}$, one will be able to find intermediate manifolds which I may call (since the numbering remains arbitrary):
\[
v'_{1}, \qquad v'_{2}, \qquad \ldots, \qquad v'_{k},
\]
and which will form a chain.

I shall always be able to choose the order of the variables analogous to the $y$ that define the manifold $v_{q+1}$ in such a way that, in $v'_{1}$, the determinant analogous to $\Delta$ and relative to $v_{1}$ and $v_{q+1}$ is positive; I shall call it $\Delta_{1}$.

Let us call likewise $\Delta_{i}$ the determinant analogous to $\Delta$ and relative to $v_{i}$ and $v_{q+1}$; $\Delta_{i}$ is defined in the interior of $v'_{i}$.

Next let $\Delta'$ be the determinant relative to $v_{1}$ and $v_{2}$; it will be defined in the whole part common to these two manifolds and, in particular, in the part common to $v'_{1}$ and $v'_{2}$; it will be positive since the system (4) is two-sided.

Then, in the part common to $v'_{1}$ and $v'_{2}$, $\Delta_{1}$ and $\Delta'$ will be positive; it follows that their ratio
\[
\Delta_{2} = \frac{\Delta_{1}}{\Delta'}
\]
is positive, and since it always keeps the same sign, it will be positive for all the points of $v'_{2}$.

One would prove step by step that $\Delta_{3}, \ldots, \Delta_{k}$ are likewise positive.
\origpage{29}

The system will therefore remain two-sided after the adjunction of $v_{q+1}$; it will still remain so if, in the expression of the equations of one of the manifolds $v_{i}$, one replaces the variables $y_{1}, y_{2}, \ldots, y_{m}$ by holomorphic functions of new variables $y'_{1}, y'_{2}, \ldots, y'_{m}$ whose functional determinant never vanishes (which is necessary if one wishes that to one system of values of the $y$ there correspond a single system of values of the $y'$). One must, of course, choose the order of the new variables $y'$ in such a way that this functional determinant is positive.

Since the system always remains two-sided, one will not be able to construct a one-sided chain in it, so that a manifold cannot be at once two-sided and one-sided. Q.~E.~D.

Everyone knows the example of a one-sided surface that one obtains by folding a rectangle of paper $ABCD$ (whose opposite sides are $AB$ and $CD$ on the one hand, $BC$ and $DA$ on the other), and then gluing together the sides $AB$ and $CD$ in such a way that $A$ is glued to $C$ and $B$ to $D$.

Examples of two-sided manifolds are easier still to form. Thus in space of $n$ dimensions:

1° Every domain of $n$ dimensions is two-sided;

2° Every curve of 1 dimension is two-sided;

3° Every \emph{closed} surface of $n - 1$ dimensions is two-sided.

But one can go further.

Let us consider a manifold $V$ defined \emph{in the first manner}, that is to say by equalities and inequalities of the form (1). I say that it will always be two-sided.

Let, indeed,
\[
F_{1} = F_{2} = \ldots = F_{p} = 0
\]
be the $p$ equalities which, with some inequalities that we shall not write, define $V$. Let us divide $V$ into a certain number of partial manifolds $v$, defined by relations of the form (8) and (9). Let
\[
v_{1}, \qquad v_{2}, \qquad \ldots, \qquad v_{q}, \qquad v_{1}
\]
\origpage{30}
be a certain number of partial manifolds forming a continuous chain, that is to say such that each of them has a common part with the next. I say that this chain is always two-sided.

Indeed, we supposed above that the functional determinants of the $p$ functions $F$ with respect to any $p$ of the variables $x$ never all vanish at once. I can therefore suppose that the manifolds $v$ are small enough that, in the interior of $v_{1}$, one of the functional determinants (which I shall call $\Delta_{1}$) does not vanish; that, in the interior of $v_{2}$, another of these determinants (which I shall call $\Delta_{2}$) does not vanish, and so on.

If this were not so, one would only have to subdivide each of the manifolds $v$ into smaller manifolds.

Let, for example,
\[
\begin{aligned}
&\Delta_{1} = \frac{\partial(F_{1}, F_{2}, \ldots, F_{p})}{\partial(x_{\alpha_{1}}, x_{\alpha_{2}}, \ldots, x_{\alpha_{p}})},\\
&\Delta_{2} = \frac{\partial(F_{1}, F_{2}, \ldots, F_{p})}{\partial(x_{\beta_{1}}, x_{\beta_{2}}, \ldots, x_{\beta_{p}})},\\
&\ldots\ldots\ldots\ldots\ldots\ldots,
\end{aligned}
\]

Since the order of the letters $x_{\alpha_{1}}, x_{\alpha_{2}}, \ldots, x_{\alpha_{p}}$ remains arbitrary and $\Delta_{1}$ keeps the same sign in the interior of $v_{1}$, I shall always be able to suppose that $\Delta_{1}$ is positive in the interior of $v_{1}$. Likewise $\Delta_{2}$ will be positive in the interior of $v_{2}$; and so on.

Let us then put, in the interior of $v_{1}$,
\[
\left\{
\begin{aligned}
&F_{1} = 0, &&F_{2} = 0, &&\ldots, &&F_{p} = 0,\\
&y^{1}_{1} = x'_{\alpha_{1}}, &&y^{1}_{2} = x'_{\alpha_{2}}, &&\ldots, &&y^{1}_{m} = x'_{\alpha_{m}}.
\end{aligned}
\right.
\tag{13}
\]

Here $m = n - p$; the variables $x'_{\alpha_{i}}$ are the $m = n - p$ variables $x$ that remain when the $p$ variables $x_{\alpha_{i}}$ have been suppressed.

Let us put likewise, in the interior of $v_{2}$,
\[
\left\{
\begin{aligned}
&F_{1} = 0, &&F_{2} = 0, &&\ldots, &&F_{p} = 0,\\
&y^{2}_{1} = x'_{\beta_{1}}, &&y^{2}_{2} = x'_{\beta_{2}}, &&\ldots, &&y^{2}_{m} = x'_{\beta_{m}}.
\end{aligned}
\right.
\tag{14}
\]
\origpage{31}

The variables $x'_{\beta_{i}}$ are the $m = n - p$ variables $x$ that remain when the $p$ variables $x_{\beta_{i}}$ have been suppressed.

And so on.

I shall suppose that the order of the variables $x'_{\alpha_{i}}$ has been chosen in such a way that one can pass from the normal order of the variables $x$, that is to say from the order
\[
x_{1}, \qquad x_{2}, \qquad \ldots, \qquad x_{n}
\]
to the order
\[
x_{\alpha_{1}}, \qquad x_{\alpha_{2}}, \qquad \ldots, \qquad x_{\alpha_{p}}, \qquad x'_{\alpha_{1}}, \qquad x'_{\alpha_{2}}, \qquad \ldots, \qquad x'_{\alpha_{m}},
\]
by a substitution of the alternating group.

Then, by solving the equations (13), one would see that in the interior of $v_{1}$ the $x$ are holomorphic functions of
\[
y^{1}_{1}, \qquad y^{1}_{2}, \qquad \ldots, \qquad y^{1}_{m}
\]
and so on. In general, in the interior of $v_{i}$, the $x$ will be holomorphic functions of
\[
y^{i}_{1}, \qquad y^{i}_{2}, \qquad \ldots, \qquad y^{i}_{m}.
\]

It is then a matter of calculating the determinant
\[
\frac{\partial(y^{2}_{1}, y^{2}_{2}, \ldots, y^{2}_{m})}{\partial(y^{1}_{1}, y^{1}_{2}, \ldots, y^{1}_{m})},
\]
or, for short,
\[
\frac{\partial y^{2}_{i}}{\partial y^{1}_{i}}
\]
in the interior of the part common to $v_{1}$ and to $v_{2}$.

For this, let us replace the equations (13) and (14) by the more general equations
\[
\left\{
\begin{aligned}
&F_{1} = \lambda_{1}, &&F_{2} = \lambda_{2}, &&\ldots, &&F_{p} = \lambda_{p},\\
&y^{1}_{i} = x'_{\alpha_{i}}
\end{aligned}
\right.
\tag{13 \textit{bis}}
\]
\origpage{32}
and
\[
F_{k} = \lambda_{k}, \qquad y^{2}_{i} = x'_{\beta_{i}}.
\tag{14 \textit{bis}}
\]

We shall then put $\lambda_{k} = 0$.

Let us solve the equations (13~\emph{bis}); we shall have the $x$ as functions of the $\lambda_{k}$ and of the $y^{1}_{i}$, holomorphic in the interior of $v_{1}$, and the functional determinant of
\[
x_{1}, \qquad x_{2}, \qquad \ldots, \qquad x_{n}
\]
with respect to
\[
\lambda_{1}, \qquad \lambda_{2}, \qquad \ldots, \qquad \lambda_{p}, \qquad y^{1}_{i}, \qquad y^{2}_{i}, \qquad \ldots, \qquad y^{m}_{i}
\]
\ednote{Printed $y^{1}_{i}, y^{2}_{i}, \ldots, y^{m}_{i}$; the variables of $v_{1}$ are apparently meant, i.e. $y^{1}_{1}, y^{1}_{2}, \ldots, y^{1}_{m}$ as in the denominator of the determinant below. Apparently a transposition of the indices.} will evidently be
\[
\frac{1}{\Delta_{1}}.
\]

Let us solve likewise the equations (14~\emph{bis}); we shall have the $x$ as functions of the $\lambda_{k}$ and of the $y^{2}_{i}$, holomorphic in the interior of $v_{2}$, and the functional determinant will be $\frac{1}{\Delta_{2}}$.

There will then result, for the points common to $v_{1}$ and to $v_{2}$,
\[
\frac{\partial y^{2}_{i}}{\partial y^{1}_{i}} = \frac{\partial(\lambda_{1}, \lambda_{2}, \ldots, \lambda_{p}, y^{2}_{1}, y^{2}_{2}, \ldots, y^{2}_{m})}{\partial(\lambda_{1}, \lambda_{2}, \ldots, \lambda_{p}, y^{1}_{1}, y^{1}_{2}, \ldots, y^{1}_{m})} = \frac{\Delta_{1}}{\Delta_{2}}.
\]

It then only remains to put $\lambda_{k} = 0$ in the expressions of the $x$.

One would find likewise, in the part common to $v_{h}$ and $v_{h+1}$,
\[
\frac{\partial y^{h+1}_{i}}{\partial y^{h}_{i}} = \frac{\Delta_{h}}{\Delta_{h+1}},
\]
and in the part common to $v_{q}$ and to $v_{1}$
\[
\frac{\partial y^{1}_{i}}{\partial y^{q}_{i}} = \frac{\Delta_{q}}{\Delta_{1}}.
\]

We have seen above that one can always suppose that $\Delta_{i}$ is,
\origpage{33}
in the interior of $v_{i}$, positive. All these ratios are therefore positive, and the chain is two-sided. Q.~E.~D.

All the manifolds that come under the first definition are therefore two-sided and, since I cited above an example of a one-sided manifold satisfying the second definition, we must conclude that the manifolds that satisfy the second definition do not all satisfy the first. This is what I had announced above.

Every one-sided manifold is by definition \emph{opposite to itself}.

\subsection*{\S~9. --- Intersection of two manifolds.}

Let $V$ and $V'$ be two manifolds defined in the second manner, the one having $p$ dimensions, the other $n-p$ dimensions, and let $M_{0}$ and $M'_{0}$ be two points belonging respectively to $V$ and to $V'$ and having respectively as coordinates
\[
x^{0}_{1}, \qquad x^{0}_{2}, \qquad \ldots, \qquad x^{0}_{n}
\]
and
\[
x'^{0}_{1}, \qquad x'^{0}_{2}, \qquad \ldots, \qquad x'^{0}_{n}.
\]

Around the point $M_{0}$, let us construct on $V$ a partial manifold $v$, analogous to those we considered in \no 3, in such a way that at a point $x_{1}, x_{2}, \ldots, x_{n}$ of $v$ one has
\[
x_{i} = \theta_{i}(y, y_{2}, \ldots, y_{p}) \qquad (i = 1, 2, n).
\]
\ednote{The first argument is printed $y$ with no index, apparently for $y_{1}$; the index may have failed to print. The range is printed $i = 1, 2, n$, apparently for $i = 1, 2, \ldots, n$.}

Likewise around $M'_{0}$, let us construct on $V'$ a manifold $v'$ such that at a point $x'_{1}, x'_{2}, \ldots, x'_{n}$ of $v'$ one has
\[
x'_{i} = \theta'_{i}(y'_{1}, y'_{2}, \ldots, y'_{n-p}) \qquad (i = 1, 2, n).
\]
\ednote{The range is printed $i = 1, 2, n$, apparently for $i = 1, 2, \ldots, n$.}
\origpage{34}

Let us consider the determinant
\[
\begin{vmatrix}
\frac{dx_{1}}{dy_{1}} & \frac{dx_{2}}{dy_{1}} & \ldots & \frac{dx_{n}}{dy_{1}} \\[1ex]
\frac{dx_{1}}{dy_{2}} & \frac{dx_{2}}{dy_{2}} & \ldots & \frac{dx_{n}}{dy_{2}} \\[1ex]
\ldots & \ldots & \ldots & \ldots \\[1ex]
\frac{dx_{1}}{dy_{p}} & \frac{dx_{2}}{dy_{p}} & \ldots & \frac{dx_{n}}{dy_{p}} \\[1ex]
\frac{dx'_{1}}{dy'_{1}} & \frac{dx'_{2}}{dy'_{1}} & \ldots & \frac{dx'_{n}}{dy'_{1}} \\[1ex]
\ldots & \ldots & \ldots & \ldots \\[1ex]
\frac{dx'_{1}}{dy'_{n-p}} & \frac{dx'_{2}}{dy'_{n-p}} & \ldots & \frac{dx'_{n}}{dy'_{n-p}}
\end{vmatrix}.
\]

It is a function of $y_{1}, y_{2}, \ldots, y_{p}, y'_{1}, y'_{2}, \ldots, y'_{n-p}$; this function will therefore depend on the position of the points $M$ and $M'$, with coordinates $x_{1}, x_{2}, \ldots, x_{n}$ and $x'_{1}, x'_{2}, \ldots, x'_{n}$, on the manifolds $v$ and $v'$. I shall therefore call it
\[
f(M, M').
\]

I can change variables by replacing $y_{1}, y_{2}, \ldots, y_{p}$ by holomorphic functions of $p$ new variables $z_{1}, z_{2}, \ldots, z_{p}$, chosen in such a way that to one system of values of $y$ there corresponds a single system of values of the $z$. For this the functional determinant of the $y$ with respect to the $z$ must not vanish, and I can always suppose, by arranging the $z$ in a suitable order, that this determinant is positive.

The function $f(M, M')$ is then multiplied by this functional determinant and consequently keeps its sign.

The same is again true when one makes an analogous change of variables on the manifold $v'$.

Let us now consider another manifold $v_{1}$, analogous to $v$, constructed on $V$, and a manifold $v'_{1}$, analogous to $v'$, constructed on $V'$.

Let $M_{1}$ be a point of $v_{1}$ and $M'_{1}$ a point of $v'_{1}$. We shall be able to construct a function analogous to $f(M, M')$, which I shall call $f_{1}(M_{1}, M'_{1})$.
\origpage{35}

Let us now suppose that the manifolds $v_{1}$ and $v$ (as well as the manifolds $v'_{1}$ and $v'$) have a common part, and that the points $M$ and $M_{1}$ coincide, as well as the points $M'$ and $M'_{1}$.

Let us compare the two functions
\[
f(M, M'), \qquad f_{1}(M, M').
\]

Let us suppose that the two manifolds $V$ and $V'$ are two-sided. We shall then be able to suppose (by suitably choosing the order of the variables $y$ relative to $v_{1}$) that the determinant analogous to the one we called $\Delta$ in \no 8, and relative to $v$ and $v_{1}$, is positive in the part common to $v$ and $v_{1}$; likewise the determinant analogous to $\Delta$ and relative to $v'$ and $v'_{1}$ will also be positive.

Then $f$ and $f_{1}$ will be of the same sign.

Let then
\[
S(M, M')
\]
be a function which will be equal to $+1$, to $-1$ or to 0, according as $f(M, M')$ is positive, negative or zero. This function will then be perfectly determined; it will depend only on the position of the points $M$ and $M'$ on $V$ and $V'$; it will not depend on the way in which the manifolds $v$ and $v'$ have been constructed around $M$ and $M'$.

\emph{This would no longer be true if one of the manifolds $V$ and $V'$ were one-sided, which I do not suppose}.

Let us suppose, in particular, that the points $M$ and $M'$ coincide and, consequently, that the point $M$ is a point of intersection of $V$ and $V'$; the consideration of the function $S(M, M')$ is then of great interest.

Let us suppose that one forms the function $S(M, M)$ for all the points of intersection $M$ of $V$ and $V'$, and let us take the sum of all the functions $S$ thus obtained; I shall denote this sum by
\[
N(V, V').
\]
\origpage{36}

This definition of $N(V, V')$ applies when the manifolds $V$ and $V'$ have been defined as in \no 3. But it can be simplified when $V$ and $V'$ have been defined as in \no 1.
Let then
\[
F_{\alpha} = 0, \qquad \varphi_{\beta} > 0 \qquad (\alpha = 1, 2, \ldots, p; \beta = 1, 2, \ldots, q)
\]
be the conditions that define $V$, and let
\[
F'_{\gamma} = 0, \qquad \varphi'_{\delta} > 0 \qquad (\gamma = 1, 2, \ldots, n-p; \delta = 1, 2, \ldots, q'),
\]
be those that define $V'$. $V$ will have $n-p$ dimensions and $V'$ will have $p$.

Let us consider a point $M$ of $V$ having as coordinates $x_{1}, x_{2}, \ldots, x_{n}$ and a point $M'$ of $V'$ having as coordinates $x'_{1}, x'_{2}, \ldots, x'_{n}$, and let us form the determinant
\[
\begin{vmatrix}
\frac{dF_{1}}{dx_{1}} & \frac{dF_{2}}{dx_{1}} & \ldots & \frac{dF_{p}}{dx_{1}} & \frac{dF'_{1}}{dx'_{1}} & \frac{dF'_{2}}{dx'_{1}} & \ldots & \frac{dF'_{n-p}}{dx'_{1}} \\[1ex]
\ldots & \ldots & \ldots & \ldots & \ldots & \ldots & \ldots & \ldots \\[1ex]
\frac{dF_{1}}{dx_{n}} & \frac{dF_{2}}{dx_{n}} & \ldots & \frac{dF_{p}}{dx_{n}} & \frac{dF'_{1}}{dx'_{n}} & \frac{dF'_{2}}{dx'_{n}} & \ldots & \frac{dF'_{n-p}}{dx'_{n}}
\end{vmatrix}
\]
which I shall call $\psi(M, M')$. What relation is there between $\psi(M, M')$ and $S(M, M')$?

To account for it, let us put
\[
x'_{i} = x_{i} + h_{i} \qquad (i = 1, 2, \ldots, n);
\]
let us then vary the $x'_{i}$ and the $x_{i}$ at once, but in such a way that their difference $h_{i}$ remains constant.

Let us put
\[
\begin{aligned}
&F_{\alpha}(x_{i}) = y_{\alpha},\\
&F'_{\gamma}(x_{i} + h_{i}) = y'_{\gamma}.
\end{aligned}
\]

If the determinant $\psi(M, M')$ is not zero, we shall be able to solve these
\origpage{37}
equations with respect to the $x_{i}$, and we shall derive from them
\[
x_{i} = \theta_{i}(y_{\alpha}, y'_{\gamma}).
\]

The functional determinant of the $\theta_{i}$ with respect to the $y_{\alpha}$ and the $y'_{\gamma}$ will then be $\frac{1}{\psi(M, M')}$.

If we put $y_{\alpha} = 0$, we shall have
\[
x_{i} = \theta_{i}(0, y'_{\gamma});
\]
this will be the equation of a manifold of $n-p$ dimensions that will form part of $V$; likewise, by putting
\[
x'_{i} = h_{i} + \theta_{i}(y_{\alpha}, 0),
\]
one will have the equation of a manifold of $p$ dimensions that will form part of $V'$.

If one defines in this way the manifolds $v$ and $v'$ of the beginning of this number, one sees that the functional determinant of the $\theta_{i}$ is nothing other than $f(M, M')$, whence
\[
f(M, M')\psi(M, M') = 1.
\]

Let us now suppose that the two manifolds $V$ and $V'$ both form part of a manifold $W$ of a greater number of dimensions. Let, for example,
\[
F_{\alpha} = 0, \qquad \varphi_{\beta} > 0 \qquad (\alpha = 1, 2, \ldots, p; \beta = 1, 2, \ldots, q) \tag{1}
\]
be the relations that define $W$, which will thus have $n-p$ dimensions.

Let
\[
F'_{\gamma} = 0, \qquad \varphi'_{\delta} > 0 \qquad (\gamma = 1, 2, \ldots, p; \delta = 1, 2, \ldots, q') \tag{2}
\]
\ednote{The range of $\gamma$ is printed as ending at $p$; apparently for $p'$, since $V$ is said to have $n-p-p'$ dimensions and the next system runs to $p''$.} be the relations which, joined to (1), define $V$, which will thus have $n-p-p'$ dimensions.

Let finally
\[
F''_{\varepsilon} = 0, \qquad \varphi''_{\zeta} > 0 \qquad (\varepsilon = 1, 2, \ldots, p''; \zeta = 1, 2, \ldots, q'') \tag{3}
\]
\origpage{38}
be the relations which, joined to (1), will define $V'$, which will thus have $n-p-p''$ dimensions. I suppose $p + p' + p'' = n$.

Let $M$ be a point of intersection of $V$ and $V'$.

I shall form the functional determinant of the $F_{\alpha}$, the $F'_{\gamma}$ and the $F''_{\varepsilon}$ with respect to the $x$, and it is this determinant that I shall call $\psi(M)$.

Then $S(M)$ will by definition be equal to $+1$ or to $-1$, according as $\psi(M)$ is positive or negative.

I shall next denote by
\[
N(V, V')
\]
the sum of all the quantities $S(M)$ relative to all the points of intersection of $V$ and $V'$.

It is clear that, if one of the manifolds $V$ or $V'$ is replaced by the opposite manifold, the expression $N(V, V')$ changes sign.

Now let $V$ be a closed manifold of 1 dimension and let $W$ be a domain of $n$ dimensions whose complete boundary consists of the manifolds of $n-1$ dimensions
\[
V_{1}, \qquad V_{2}, \qquad \ldots, \qquad V_{k},
\]
so that
\[
V_{1} + V_{2} + \ldots + V_{k} \backsim 0.
\]

I say that
\[
N(V, V_{1}) + N(V, V_{2}) + \ldots + N(V, V_{k}) = 0.
\]

Indeed, the manifold $V$ will consist of a certain number of continuous manifolds of 1 dimension, that is to say of a certain number of closed lines. The equations of one of these lines can be put in the form
\[
x_{i} = \theta_{i}(t),
\]
$t$ being a variable that we shall make increase from $-\infty$ to $+\infty$.

Since the line is closed, we shall have
\[
\theta_{i}(-\infty) = \theta_{i}(+\infty).
\]
\origpage{39}

Let us now consider a point $M$ of intersection of one of these lines with the boundary of $W$; if, when $t$ is made to increase, this line passes from the interior of $W$ to the exterior, one will have
\[
S(M) = 1.
\]
If it passes from the exterior to the interior, $S(M)$ will be equal to $-1$. But the line, being closed, will return to its point of departure, so that the sum of all the $S(M)$ must be zero. Q.~E.~D.

Let us now suppose that all the manifolds we are considering, for example $V$, $W$, $V_{1}, V_{2}, \ldots, V_{k}$, form part of a manifold $U$, and that it is this manifold $U$ whose orders of connection we are studying and with respect to which we take the homologies.

(Here is what I mean; when, in \no 5, I defined homologies, there was a certain manifold $V$ that played an important part in that definition; well, here it is $U$ that will play that part.)

I suppose that $U$ has $h$ dimensions, $V$ one dimension, $W$ $h$ dimensions, $V_{1}, V_{2}, \ldots, V_{k}$ $h-1$ dimensions.

There is nothing to change in what precedes, and one sees that, if $V$ is closed and of one dimension and
\[
V_{1} + V_{2} + \ldots + V_{k} \backsim 0,
\]
one will have
\[
N(V, V_{1}) + N(V, V_{2}) + \ldots + N(V, V_{k}) = 0.
\]

I say that this will still be true if the manifold of one dimension $V$ is no longer closed, but has its two extremities on the boundary of $U$, and if $W$ has no point on the boundary of $U$.

Indeed, $U$ is divided into two regions, continuous or not, namely: $W$ and the region $R$ exterior to $W$. By hypothesis, the boundary of $U$ belongs entirely to $R$.

The line $V$ not being closed, the initial point corresponding to $t = -\infty$ will no longer coincide with the final point corresponding to $t = +\infty$; but,
\origpage{40}
by hypothesis, these two points both belong to the boundary of $U$ and consequently to $R$.

Our line will therefore pass exactly as many times from $R$ into $W$ as from $W$ into $R$; hence the sum of the $S(M)$ is zero. Q.~E.~D.

I am now going to establish the converse.

Let $V_{1}, V_{2}, \ldots, V_{k}$ be $k$ closed manifolds of $h-1$ dimensions situated in $U$, $U$ having no point in common with the boundary of $V_{i}$. Let us suppose that one does not have
\[
V_{1} + V_{2} + \ldots + V_{k} \backsim 0,
\]
I say that one will always be able to find a line $V$, situated in $U$, which will be closed, or\ednote{Printed ``on''; apparently for ``ou'', a turned letter.} will have its two extremities on the boundary of $U$, and which will be such that
\[
N(V, V_{1}) + N(V, V_{2}) + \ldots + N(V, V_{k}) \gtrless 0.
\]

Indeed, either the manifolds $V_{1}, V_{2}, \ldots, V_{k}$ will not divide $U$ into several distinct regions. Then one will be able to go from any point of $U$ to any other point of $U$ without meeting any of the manifolds $V_{i}$.

Let us then consider a very small line meeting $V_{1}$ in a point $M$ and meeting none of the other manifolds $V_{i}$; we shall be able to join the two extremities of this line by another continuous line that will meet none of the manifolds $V_{i}$; we shall thus have constructed a closed line that will meet the manifolds $V_{i}$ in only a single point. The sum of the $S(M)$ therefore cannot be zero.

Let us now suppose that the manifolds $V_{i}$ divide $U$ into several regions $R$, but that they are nevertheless linearly independent. In this case the complete boundary of none of the regions $R$ can be constituted by a part of the manifolds $V_{i}$, otherwise there would be a homology between these manifolds, which we do not suppose.

The complete boundary of $R$ will therefore consist of some of the
\origpage{41}
manifolds $V_{i}$ and of a part of the boundary of $U$. It follows that from any point of $R$ one will be able to go to the boundary of $U$ without leaving $R$ and without meeting any of the manifolds $V_{i}$.

Let us then consider one of the manifolds $V_{i}$, $V_{1}$ for example, and let us consider a small line cutting $V_{1}$ in a point $M$; let $R$ and $R'$ be the regions to which the two extremities $\mu$ and $\mu'$ of this small line belong. From $\mu$ one will be able to go by a continuous line to the boundary of $U$ without leaving $R$; from $\mu'$ one will be able to go by a third continuous line to the boundary of $U$ without leaving $R'$.

The set of these three continuous lines will go from the boundary of $U$ to the boundary of $U$, and if I call it $V$, one will have
\[
N(V, V_{1}) = 1, \qquad N(V, V_{i}) = 0 \qquad (i > 1).
\]

One can therefore choose $V$ in such a way that all the $N(V, V_{i})$ are zero, except one of them which will be equal to 1.

Let us suppose finally that there are between the $V_{i}$ a certain number of homologies, 3 for example,
\[
\Sigma k_{i}V_{i} \backsim \Sigma k'_{i}V_{i} \backsim \Sigma k''_{i}V_{i} \backsim 0. \tag{$\alpha$}
\]

We shall then be able, among the $k$ manifolds $V_{i}$, to find $k-3$ that are linearly independent; let them be, for example, $V_{1}, V_{2}, \ldots, V_{k-3}$.

We can choose $V$ in such a way that all the $N(V, V_{i})$, where $i = 1, 2, \ldots, k-3$, are zero with the exception of any one of them, which will be equal to 1. The values of
\[
N(V, V_{k-2}), \qquad N(V, V_{k-1}), \qquad N(V, V_{k})
\]
will be deduced from them by means of the relations
\[
\Sigma k_{i}N(V, V_{i}) = \Sigma k'_{i}N(V, V_{i}) = \Sigma k''_{i}N(V, V_{i}) = 0 \tag{$\beta$}
\]
which are a necessary consequence of our homologies ($\alpha$).
\origpage{42}

It cannot then happen that one has
\[
\Sigma N(V, V_{i}) = 0, \tag{$\gamma$}
\]
whichever of the $k-3$ quantities
\[
N(V, V_{i}) \qquad (i = 1, 2, \ldots, k-3)
\]
is equal to 1. That could happen, indeed, only if the equation ($\gamma$) were a necessary consequence of the equations $\beta$; but then the homology
\[
\Sigma V_{i} \backsim 0 \tag{$\delta$}
\]
would be a necessary consequence of the homologies ($\alpha$). Now we supposed at the outset that this homology ($\delta$) did not hold.

One will therefore always be able to choose $V$ in such a way that the equality $\gamma$ does not hold. Q.~E.~D.

We shall designate by the name of \emph{cut} every manifold contained in $U$: 1° if it is closed; 2° if it is not closed, but its boundary forms part of the complete boundary of $U$.

We can then state the following theorem, which sums up the preceding discussion:

\emph{The necessary and sufficient condition} (if the manifolds $V_{i}$ are closed and of $h-1$ dimensions) \emph{for one to be able to choose the cut} $V$ \emph{in such a way that the equality}
\[
\Sigma k_{i}N(V, V_{i}) = 0
\]
\emph{does not hold, is that the homology}
\[
\Sigma k_{i}V_{i} \backsim 0
\]
\emph{does not hold}.
\origpage{43}

Let us now seek to extend this theorem to the case where $V$ has $p$ dimensions and $V_{1}, V_{2}, \ldots, V_{k}$ have $h-p$.

Let
\[
F_{\alpha} = 0, \qquad \varphi_{\beta} > 0 \qquad (\alpha = 1, 2, \ldots, n-h)
\]
be the equations of $U$; let
\[
F_{\alpha} = 0, \qquad F'_{\gamma} = 0, \qquad \varphi_{\beta} > 0 \qquad (\gamma = 1, 2, \ldots, h-p)
\]
be those of $V$.

As for $V_{1}, V_{2}, \ldots, V_{k}$, we shall define them in the following manner. We shall always be able to find $p-1$ equations
\[
\Phi_{1} = \Phi_{2} = \ldots = \Phi_{p-1} = 0
\]
which all the points of $V_{1}, V_{2}, \ldots, V_{k}$ satisfy; thus, if $h = 3$, $p = 2$, if $U$ is ordinary space and if $V_{1}, V_{2}, \ldots, V_{k}$ are curves, one will always be able to pass a surface through these curves, whatever they may be.

To define $V_{i}$ we shall adjoin to these a $p^{\text{th}}$ equality
\[
F''_{i} = 0.
\]

Let then $U'$ be the manifold of $h-p+1$ dimensions
\[
F_{\alpha} = 0, \qquad \Phi_{\nu} = 0, \qquad \varphi_{\beta} > 0,
\]
and $V'$ the manifold of 1 dimension
\[
F_{\alpha} = F'_{\gamma} = \Phi_{\nu} = 0, \qquad \varphi_{\beta} > 0.
\]

One will have first
\[
N(V', V_{i}) = N(V, V_{i}).
\]

On the other hand, if the homology
\[
\Sigma k_{i}V_{i} \backsim 0
\]
\origpage{44}
holds with respect to $U'$, it will hold with respect to $U$. It is true that the converse is not true, and that this homology may hold with respect to $U$ without holding with respect to $U'$; but, if it holds with respect to $U$, one will always be able, by suitably choosing the functions $\Phi$, to find a manifold $U'$ of $h-p+1$ dimensions with respect to which it holds.

\emph{We must therefore conclude that the theorem is still true when} $V$ \emph{has more than one dimension}.

If, therefore, $V_{1}$ and $V_{2}$ are two closed manifolds of $h-p$ dimensions, and if for all the cuts $V$ of $p$ dimensions one has
\[
N(V, V_{1}) = N(V, V_{2}),
\]
one will also have
\[
V_{1} \backsim V_{2}
\]
and conversely.

Let us confine ourselves to the case where $U$ is closed. Then $U$ has no boundary, and all the cuts are closed.

The number of linearly independent cuts of $p$ dimensions is $P_{p} - 1$.

Let
\[
C_{1}, \qquad C_{2}, \qquad \ldots, \qquad C_{\lambda} \qquad (\lambda = P_{p} - 1)
\]
be these cuts.

Next let
\[
V_{i} \qquad (i = 1, 2, \ldots, \mu)
\]
be $\mu$ closed manifolds of $h-p$ dimensions.
The necessary and sufficient condition for a homology to hold between them,
\[
\Sigma k_{i}V_{i} \backsim 0,
\]
is that one has:
\[
\Sigma k_{i}N(C_{1}, V_{i}) = \Sigma k_{i}N(C_{2}, V_{i}) = \ldots = \Sigma k_{i}N(C_{\lambda}, V_{i}) = 0.
\]
\origpage{45}

Now, if the number $\mu$ of the $V_{i}$ is greater than $\lambda$, one will always be able to find integers $k_{i}$ satisfying these conditions, since we shall have $\mu$ arbitrary $k_{i}$ and $\lambda$ conditions to fulfil. If, therefore, the $V_{i}$ are linearly independent, one will have
\[
\mu \leqq \lambda.
\]

Hence
\[
P_{h-p} \leqq P_{p};
\]
but, changing $p$ into $h-p$, one finds
\[
P_{p} \leqq P_{h-p}.
\]

Hence
\[
P_{p} = P_{h-p}.
\]

\emph{Consequently, for a closed manifold, the Betti numbers equally distant from the extremes are equal}.

This theorem has never, I believe, been stated; it was nevertheless known to several persons, who have even made applications of it.

Let us now see what happens when $h$ is even, for the middle number $P_{\frac{h}{2}}$. Let us suppose that $h$ is a multiple of $4 + 2$, so that $\frac{h}{2}$ is odd.

One knows that when, in a determinant, one makes an odd number of permutations of rows, it changes sign. When one permutes $V$ and $V_{1}$, which have $\frac{h}{2}$ dimensions, the determinant $f(M, M)$\ednote{Printed $f(M, M)$ with no prime; apparently for $f(M, M')$, as on pp. 34--37.} will change sign; one will therefore have
\[
N(V, V_{1}) = -N(V_{1}, V);
\]
one deduces from this
\[
N(V, V) = 0.
\]

The symbol $N(V, V)$ has by itself no meaning since, the two manifolds $V$ and $V$ coinciding, the number of points of intersection is infinite;
\origpage{46}
one must therefore make precise\ednote{Printed ``précisément''; apparently for ``préciser''.} the definition; we shall therefore put
\[
N(V, V) = N(V, V'),
\]
where
\[
V \backsim V'.
\]

This being so, I say that $P_{\frac{h}{2}}$ is odd. Let us suppose, indeed, that it is even, and let
\[
V_{1}, \qquad V_{2}, \qquad \ldots, \qquad V_{\mu}
\]
be the $\mu$ linearly independent manifolds of $h-1$\ednote{Printed $h-1$; apparently for $\frac{h}{2}$, the dimension of the varieties under discussion.} dimensions, where
\[
\mu = P_{\frac{h}{2}} - 1
\]
is odd.

Let us form the determinant in which the $i^{\text{th}}$ term of the $k^{\text{th}}$ column is $N(V_{i}, V_{k})$. This determinant would be skew; that is to say, the terms of the principal diagonal would be zero, and two terms symmetric with respect to this diagonal would be equal and of opposite sign. Since the number of columns would be odd, this determinant would be zero. It would follow, contrary to the hypothesis, that the manifolds $V_{1}, V_{2}, \ldots, V_{\mu}$ would not be linearly independent.

Hence $P_{\frac{h}{2}}$ is odd.

This would no longer be true either if $h$ were a multiple of 4, or if the manifold $U$ were one-sided; for all our reasonings suppose the manifolds two-sided. We shall see examples of this further on.

\subsection*{\S~10. --- Geometric representation.}

There is a way of representing to oneself the manifolds of three dimensions situated in space of four dimensions, a way that singularly facilitates their study.

Let us consider, in ordinary space, a certain number of polyhedra
\origpage{47}
\[
P_{1}, \qquad P_{2}, \qquad \ldots, \qquad P_{n}.
\]

We can suppose that there exist, in space of four dimensions, a certain number of manifolds of three dimensions
\[
Q_{1}, \qquad Q_{2}, \qquad \ldots, \qquad Q_{n}
\]
homeomorphic respectively to $P_{1}, P_{2}, \ldots, P_{n}$.

Let $F_{1}$ be a face of the polyhedron $P_{1}$ and $\Phi_{1}$ the set of points of the boundary of $Q_{1}$ that correspond to the various points of $F_{1}$. Let likewise $F_{2}$ be a face of $P_{2}$ and $\Phi_{2}$ the \emph{image} of this face on the boundary of $Q_{2}$.

It may be that $\Phi_{1}$ coincides with $\Phi_{2}$. In this case the two manifolds $Q_{1}$ and $Q_{2}$ are contiguous, and one passes from the interior of the one to the interior of the other by crossing $\Phi_{1}$.

In this case we shall say that the faces $F_{1}$ and $F_{2}$ are \emph{conjugate}.

It may happen that the faces $F_{1}$ and $F_{2}$ belong to the \emph{same} polyhedron $P_{1}$. Then the manifold of two dimensions $\Phi_{1}$, which does not differ from the manifold of two dimensions $\Phi_{2}$, will separate two portions of the manifold $Q_{1}$ from one another.

This will be better understood if I compare it with an example taking place in ordinary space. Let us consider a rectangle $ABCD$ and a torus, on which we shall trace two cuts, namely: a meridian circumference and a parallel; let $H$ be the point of intersection of this parallel and this meridian. The surface of the torus thus cut up will be homeomorphic to the rectangle; the two lips of the cut formed by the meridian will correspond to the two sides $AB$ and $CD$ of the rectangle; the two lips of the cut formed by the parallel will correspond to the two other sides $AD$ and $BC$. One sees the analogy with what precedes: the rectangle corresponds to the polyhedron $P_{1}$, the cut-up torus to the manifold $Q_{1}$, the sides $AB$ and $CD$ to the two faces $F_{1}$ and $F_{2}$, the two lips of the meridian cut to the two manifolds $\Phi_{1}$ and $\Phi_{2}$, which, as one sees, coincide.

This being so, let us imagine that among the faces of the $n$ polyhedra $P_{i}$ there are
\origpage{48}
a certain number that are conjugate in pairs, and a certain number that remain \emph{free}.

Let us consider the total manifold $V$ formed by the set of the manifolds $Q_{i}$. Since some of these manifolds $Q_{i}$ are contiguous, it may be that the total manifold $V$ is continuous: this is what I shall suppose.

If among the faces of the $P_{i}$ there is none that remains \emph{free}, the manifold $V$ will be closed. In the contrary case, the points that correspond to the faces that remained free will form the complete boundary of $V$.

One conceives that the knowledge of the polyhedra $P_{i}$ and of the mode of conjugation of their faces furnishes us, in ordinary space, with an image of the manifold $V$, and that this image suffices for the study of its properties from the point of view of \emph{Analysis Situs}.

Let us see how the mode of conjugation of the faces must be defined. It is clear first that, for two faces to be able to be conjugate, they must have the same number of sides. Next, to know the mode of conjugation completely, it is not enough to know that such a face is conjugate to such another; one must know, in addition, which vertex or which side of this face corresponds to which vertex or which side of its conjugate.

Only then will the mode of conjugation be completely defined.

To two conjugate faces there corresponds, by definition, one and the same manifold of two dimensions interior to $V$. It may also be that to several edges of the polyhedra $P$ there corresponds one and the same line interior to $V$, or that to several vertices of these polyhedra there corresponds one and the same point interior to $V$. We shall then say that these edges or these vertices belong to one and the same cycle.

Here is how one can form the cycles of edges and the cycles of vertices.

Let $A_{1}$ be an edge, $F'_{1}$ one of the two faces that pass through $A_{1}$, $F_{2}$ the conjugate of $F'_{1}$, $A_{2}$ the edge of $F_{2}$ that corresponds to $A_{1}$; $F'_{2}$ the other face that passes through $A_{2}$; $F_{3}$ the conjugate of $F'_{2}$, $A_{3}$ the edge corresponding to $A_{2}$, etc.

One will stop when one arrives at a free face or when one has
\origpage{49}
returned to the edge $A_{1}$. All the edges $A_{1}$, $A_{2}$, $A_{3}$, \ldots\ will form a cycle.

Likewise for the vertices. Let $S_{1}$ be a vertex, $F'_{1}$ one of the faces that pass through $S_{1}$, $F_{2}$ the conjugate of $F'_{1}$; $S_{2}$ the vertex corresponding to $S_{1}$; $F'_{2}$ \emph{one} of the faces that pass through $S_{2}$, \ldots.

$S_{1}$, $S_{2}$, $S_{3}$, \ldots\ will belong to one and the same cycle.

Only here more than two faces pass through each vertex; then $F'_{2}$, for example, can be chosen in several ways; one must stop only after having exhausted all the possible choices.

The analogy with the formation of cycles in the theory of Fuchsian groups is evident. It becomes still more so if one supposes that there is only a single polyhedron $P_{1}$.

\subsection*{First example.}

The simplest example is that in which one has only a single polyhedron $P_{1}$, and this polyhedron is a cube $ABCDA'B'C'D'$ whose vertices have respectively as coordinates
\[ \begin{matrix} A & 0 & 0 & 0 & \qquad & A' & 0 & 0 & 1 \\[1ex] B & 0 & 1 & 0 & \qquad & B' & 0 & 1 & 1 \\[1ex] C & 1 & 0 & 0 & \qquad & C' & 1 & 0 & 1 \\[1ex] D & 1 & 1 & 0 & \qquad & D' & 1 & 1 & 1 \end{matrix} \]

I suppose that the opposite faces are conjugate, and this in the following manner:
\[ \left\{ \begin{aligned} &ABDC \equiv A'B'D'C',\\ &ACC'A' \equiv BDD'A',\\ &CDD'C' \equiv ABB'A'. \end{aligned} \right. \tag{1} \]
\ednote{In the first row only broken fragments of the sign $\equiv$ printed; the sign is read as $\equiv$, as in the other rows. The second row prints $BDD'A'$, which is not a face of the cube; apparently for $BDD'B'$, the face opposite $ACC'A'$, as in the restatement of this example on p. 65.}

Here is what I mean by this notation: the relation
\[ ABDC \equiv A'B'D'C' \]
signifies:
\origpage{50}

1° That the faces $ABDC$ and $A'B'D'C'$ are conjugate;

2° That the vertices met on the first of these faces are met in the circular order $ABDC$;

3° That the vertices $A$ and $A'$, $B$ and $B'$, $D$ and $D'$, $C$ and $C'$ correspond to one another.

Let us observe in passing that the vertices $ABDC$ are met on the first of these faces in a circular order such that, in traversing the perimeter in the corresponding sense, one has the face on one's left. On the contrary, the corresponding vertices $A'B'D'C'$ are met on the second face in an order such that, in traversing the perimeter in that sense, one has the face on one's right.

This condition must always be fulfilled \emph{if one wishes $V$ to be two-sided}. If two faces $F_{1}$ and $F_{2}$ are conjugate, and if a point describes the perimeter of $F_{1}$ so as to leave the face on its left, the corresponding point on $F_{2}$ must describe the perimeter of that face leaving it on its right.

This being so, let us then suppose that the mode of conjugation is defined by the relations (1). It is then easy to see that all the vertices form a single cycle, and that there are three cycles of edges, comprising respectively those edges that are parallel to the axis of $x$, to that of $y$ and to that of $z$.

\subsection*{Second example.}

Let us still keep our cube, but change the mode of conjugation and define it by the relations
\[ \left\{ \begin{aligned} &ABDC \equiv B'D'C'A',\\ &ABB'A' \equiv DD'C'C,\\ &ACC'A' \equiv DD'B'B; \end{aligned} \right. \tag{2} \]
there will then be two cycles of edges and two cycles of vertices, and I sum up the results by the following relations; I have put the sign $\equiv$ between two edges (or two vertices) to express that they form part of one and the same cycle.
\origpage{51}

Two cycles of edges:
\[ \begin{aligned} &AB \equiv B'D' \equiv C'C \equiv B'A' \equiv AC \equiv DD',\\ &AA' \equiv DC \equiv C'A' \equiv B'B \equiv C'D' \equiv DB. \end{aligned} \]

Two cycles of vertices:
\[ \begin{aligned} &A \equiv B' \equiv C' \equiv D,\\ &B \equiv D' \equiv C \equiv A'. \end{aligned} \]
\ednote{In the second row only two strokes of the last sign $\equiv$, between $C$ and $A'$, printed; it is read as $\equiv$, as elsewhere in these cycles.}

We shall see later for what reasons this mode of conjugation is inadmissible.

\subsection*{Third example.}

Let us keep our cube with the following mode of conjugation:
\[ \left\{ \begin{aligned} &ABDC \equiv B'D'C'A',\\ &ABB'A' \equiv DD'C'C,\\ &ACC'A' \equiv DD'B'B. \end{aligned} \right. \tag{3} \]
\ednote{As printed, system (3) repeats system (2) letter for letter, yet the cycles stated below differ from those of the second example. The stated cycles fit a second row $ABB'A' \equiv C'CDD'$; the printed $DD'C'C$ is apparently a misprint.}

One then finds:

Four cycles of edges:
\[ \begin{aligned} &AB \equiv B'D' \equiv C'A', &&AA' \equiv C'D' \equiv DB,\\ &AC \equiv DD' \equiv B'A', &&CD \equiv BB' \equiv A'C'. \end{aligned} \]
\ednote{In the first row only two strokes of the sign $\equiv$ between $AA'$ and $C'D'$ printed; it is read as $\equiv$. The first cycle prints $C'A'$, so that $A'C'$ occurs twice and $CC'$ never; apparently for $C'C$.}

Two cycles of vertices:
\[ \begin{aligned} &A \equiv B' \equiv C' \equiv D,\\ &B \equiv D' \equiv A' \equiv C. \end{aligned} \]
\ednote{In the first row only two strokes of the sign $\equiv$ between $C'$ and $D$ printed; it is read as $\equiv$.}

\subsection*{Fourth example.}

Let now
\[ \left\{ \begin{aligned} &ABDC \equiv B'D'C'A',\\ &ABB'A' \equiv CDD'C',\\ &ACC'A' \equiv BDD'B'. \end{aligned} \right. \tag{5} \]
\ednote{The print numbers this system (5), directly after (3); no equation (4) precedes it.}
\origpage{52}

One finds:

Three cycles of edges:
\[ \begin{aligned} &AA' \equiv CC' \equiv BB' \equiv DD',\\ &AB \equiv CD \equiv B'D' \equiv A'C',\\ &AC \equiv BD \equiv D'C' \equiv B'A', \end{aligned} \]
and a single cycle of vertices.

\subsection*{Fifth example.}

Let there be a regular octahedron; it will have six vertices, of which four will form a square $BCED$ (the letters are arranged in the circular order in which one meets the four vertices in traversing the perimeter of the square) and two, $A$ and $P$\ednote{Printed $P$; apparently for $F$, the vertex named in the system (6) below.}, are outside the square. Let
\[ \left\{ \begin{aligned} &ABC \equiv FED,\\ &ACE \equiv FDB,\\ &AED \equiv FBC,\\ &ADB \equiv FCE \end{aligned} \right. \tag{6} \]
be the mode of conjugation; one will find six cycles of edges and three cycles of vertices; each edge will form a cycle with the opposite edge, that is to say with its symmetric with respect to the centre of symmetry of the octahedron, and each vertex forming a cycle with the opposite vertex.

It is useless to multiply examples further; I now propose to recognize whether all modes of conjugation are admissible.

Let us call \emph{aster} the figure formed by a certain number of solid angles having the same vertex and arranged around this vertex in such a way that every point of space belongs to one and only one of these solid angles.

In the figure one will therefore distinguish, besides the solid angles themselves and their vertex, a certain number of faces common to two solid angles and a certain number of edges common to several solid angles.
\origpage{53}

I may suppose that the edges and the faces are prolonged indefinitely, or else that they are stopped at the surface of a sphere having its centre at the common vertex of the solid angles. Then these various solid angles will cut out on the sphere a certain number of spherical polygons, so that the surface of the sphere will be subdivided into spherical polygons. This surface thus subdivided may be regarded as homeomorphic to a convex polyhedron whose faces will correspond to the spherical polygons I have just defined.

Then to the solid angles of the aster there will correspond the faces of the polyhedron; to the faces of the aster there will correspond the edges of the polyhedron, and to the edges of the aster its vertices.

Let $S$, $F$ and $A$ be the number of the solid angles, of the faces and of the edges of the aster. Since the polyhedron must satisfy Euler's theorem, one must have
\[ S - F + A = 2. \]

Let us now return to our polyhedra
\[ P_{1}, \qquad P_{2}, \qquad \ldots, \qquad P_{n} \]
and consider a cycle of vertices; to all the vertices of this cycle there will correspond one and the same point of $V$, which I call $M$. Among the manifolds
\[ Q_{1}, \qquad Q_{2}, \qquad \ldots, \qquad Q_{n}, \]
there will be a certain number that have the point $M$ on their boundary; I shall call them the manifolds $Q_{\alpha}$, and they will be those that correspond to polyhedra $P_{\alpha}$ to which the various vertices of the cycle belong.

Let us now consider two of the manifolds $Q_{\alpha}$ that are contiguous, and the manifold of two dimensions that serves as their common boundary. The manifolds of two dimensions thus defined I shall call the $\Phi_{\alpha}$; they will correspond to those faces of the polyhedra $P_{\alpha}$ to which the various vertices of the cycle belong.

We shall consider finally the manifolds of one dimension that are the intersection
\origpage{54}
of two manifolds $\Phi_{\alpha}$, and which I shall call the $L_{\alpha}$. They will correspond to those edges of the polyhedra $P_{\alpha}$ to which the various vertices of the cycle belong.

Let us consider the figure formed by the manifolds $Q_{\alpha}$, $\Phi_{\alpha}$, $L_{\alpha}$, or rather by the points of these manifolds that satisfy the inequality
\[ (x_{1} - x^{0}_{1})^{2} + (x_{2} - x^{0}_{2})^{2} + (x_{3} - x^{0}_{3})^{2} + (x_{4} - x^{0}_{4})^{2} < \varepsilon^{2}. \tag{7} \]

$\varepsilon$ is a very small quantity; $x^{0}_{1}$, $x^{0}_{2}$, $x^{0}_{3}$, $x^{0}_{4}$ are the coordinates of the point $M$.

\emph{This figure is evidently homeomorphic to an aster} whose faces and edges are limited to a sphere. Let $A$ be this aster.

Let us consider any one of the manifolds $Q_{\alpha}$ and, in addition, the manifold $W$ formed by those of its points that satisfy the inequality (7). This manifold will be continuous or will not be. If it is not, I shall decompose it into several continuous manifolds, and the various continuous manifolds that can thus be formed I shall call the $Q'_{\alpha}$. (The number of the $Q'_{\alpha}$ may thus be greater than that of the $Q_{\alpha}$ if one of the manifolds $W$ is not continuous.)

I shall define likewise the $\Phi'_{\alpha}$ and the $L'_{\alpha}$.

Let $q_{\alpha}$, $\varphi_{\alpha}$, $l_{\alpha}$ be the number of the $Q'_{\alpha}$, of the $\Phi'_{\alpha}$ and of the $L'_{\alpha}$; this will also be that of the solid angles, of the faces and of the edges of the aster $A$.

Whence this conclusion:

\emph{For a mode of conjugation to be admissible, it is necessary that, for all the cycles of vertices, one has}
\[ q_{\alpha} - \varphi_{\alpha} + l_{\alpha} = 2. \]

Let us now see how one will be able to calculate the numbers $q_{\alpha}$, $l_{\alpha}$\ednote{Printed $l_{\alpha}$; apparently for $\varphi_{\alpha}$, since the three numbers treated in 1°, 2°, 3° below are $q_{\alpha}$, $\varphi_{\alpha}$, $l_{\alpha}$.}, $l_{\alpha}$ for any cycle of vertices $\alpha$.

1° $q_{\alpha}$ will be the number of vertices of the cycle.

2° To have $\varphi_{\alpha}$, one has only to take the sum of the numbers of faces that pass through the various vertices of the cycle and divide by two. If,
\origpage{55}
for example, the cycle $\alpha$ consists of three vertices $a$, $b$, $c$ belonging respectively to the polyhedra $P_{1}$, $P_{2}$ and $P_{3}$; if the point $a$ is the vertex of a trihedral angle, so that three faces of $P_{1}$ pass through $a$; if four faces of $P_{2}$ pass through $b$ and five faces of $P_{3}$ through $c$, one will have
\[ \varphi_{\alpha} = \frac{3 + 4 + 5}{2} = 6. \]

3° To have $l_{\alpha}$, one will enumerate the edges that end at the various vertices of the cycle in the following manner: all the edges of one and the same cycle of edges will count only as one, if one of the extremities belongs to the cycle $\alpha$; they will count as two if both extremities belong to the cycle $\alpha$.

Let us apply these rules to the six examples\ednote{Printed "six"; the Tableau that follows lists five examples, and the fifth example of the preceding pages is the last one treated there.} with which we were concerned above; we shall arrive at the following Table:
\[ \begin{matrix} \text{Example.} & q_{\alpha}. & \varphi_{\alpha}. & l_{\alpha}. \\[1ex] \text{First}\ldots\ldots & 8 & 12 & 6 \\[1ex] \text{Second}\ldots\ldots & 4 & 6 & 2 \\[1ex] \text{Third}\ldots\ldots & 4 & 6 & 4 \\[1ex] \text{Fourth}\ldots\ldots & 8 & 12 & 6 \\[1ex] \text{Fifth}\ldots\ldots & 2 & 4 & 4 \end{matrix} \]

It is to be remarked that, in the examples 2, 3 and 5, there are several cycles of vertices, but it happens that the three numbers $q_{\alpha}$, $\varphi_{\alpha}$ and $l_{\alpha}$ have the same values for all the cycles of one and the same example.

This Table shows that the relation
\[ q_{\alpha} - \varphi_{\alpha} + l_{\alpha} = 2 \]
is satisfied for all our examples, except for the second. The mode of conjugation of the second example must therefore be rejected.

\subsection*{\S~11. --- Representation by a discontinuous group.}

Here is another mode of representation which can also be employed in certain cases.
\origpage{56}

Let $(x, y, z)$ be a point of ordinary space; let us consider a series of substitutions that change $x$, $y$ and $z$ respectively into
\[ \begin{matrix} \varphi_{1}(x, y, z), & \psi_{1}(x, y, z), & \chi_{1}(x, y, z), \\[1ex] \varphi_{2}(x, y, z), & \ldots\ldots, & \ldots\ldots, \\[1ex] \ldots\ldots, & \ldots\ldots, & \ldots\ldots, \\[1ex] \varphi_{n}(x, y, z), & \psi_{n}(x, y, z), & \chi_{n}(x, y, z), \\[1ex] \ldots\ldots, & \ldots\ldots, & \ldots\ldots \end{matrix} \]

Let us suppose that the set of these substitutions forms a properly discontinuous group. Space will be divided into an infinity of domains
\[ D_{0}, \qquad D_{1}, \qquad D_{2}, \qquad \ldots, \]
and this in such a way that to each domain $D_{i}$ there corresponds a substitution of the group $S_{i}$ which changes $D_{0}$ into $D_{i}$.

Let us consider a surface $\Sigma$ forming a portion of the boundary of $D_{0}$ separating this domain $D_{0}$ from the domain $D_{i}$; the substitution $S^{-1}_{i}$, inverse of $S_{i}$, changes $D_{i}$ into $D_{0}$ and, since the points of $\Sigma$ belong to the boundary of $D_{i}$, the transform of the surface $\Sigma$ will be another part of the boundary of $D_{0}$.

The boundary of $D_{0}$ will thus be divided into portions of surface that will be conjugate in pairs, in such a way that each of them will be the transform of its conjugate by one of the substitutions of the group.

The domain $D_{0}$, with its boundary thus subdivided, will be homeomorphic to a polyhedron whose faces are conjugate in pairs, as in the preceding number. One will then be able, as in the preceding number, to make correspond to this polyhedron, or consequently to the domain $D_{0}$, a closed manifold of three dimensions situated in space of four dimensions, which one will obtain by transporting $D_{0}$ into space of four dimensions and then deforming this domain until one can \emph{glue} the conjugate portions of its boundary against one another.

The analogy with the theory of Fuchsian groups is too evident for it to be necessary to insist on it; I shall confine myself to a single example:
\origpage{57}
\subsection*{Sixth example.}

Let us consider the group derived from the three substitutions
\[ \left\{ \begin{aligned} &(x, y, z; x + 1, y, z),\\ &(x, y, z; x, y + 1, z),\\ &(x, y, z; \alpha x + \beta y, \gamma x + \delta y, z + 1), \end{aligned} \right. \tag{1} \]
$\alpha$, $\beta$, $\gamma$, $\delta$ being integers such that
\[ \alpha\delta - \beta\gamma = 1. \]

I shall call this group the group $(\alpha, \beta, \gamma, \delta)$.

I say that it is properly discontinuous.

To account for this, let us seek how there are distributed in space the transforms of one and the same point
\[ x = a, \qquad y = b, \qquad z = c. \]

First, all the transforms of this point by any combination of the first two substitutions will be comprised in the formula
\[ x = a + k, \qquad y = b + k_{1}, \qquad z = c, \tag{2} \]
$k$ and $k_{1}$ being two integers; it is easy to see, moreover, that these two substitutions are permutable.

Let us now transform the set of the points (2) by the third substitution; we shall find
\[ x = \alpha(a + k) + \beta(b + k_{1}), \qquad y = \gamma(a + k) + \delta(b + k_{1}), \qquad z = c + 1. \tag{3} \]

If we put, for short,
\[ \alpha a + \beta b = a_{1}, \qquad \gamma a + \delta b = b_{1}, \]
so that the point $(a_{1}, b_{1})$ is the transform of the point $(a, b)$ by the substitution
\origpage{58}
\[ (x, y; \alpha x + \beta y, \gamma x + \delta y), \]
which I shall call $s$.\ednote{In the display the first $x$ after $\alpha$ prints like an $r$; apparently damaged type or a misprint for $x$.}

We shall then be able to replace the equations (3) by the following:
\[ x = a + k', \qquad y = b_{1} + k'_{1}, \qquad z = c + 1, \tag{4} \]
$k'$ and $k'_{1}$ being two new integers; by applying the first two substitutions to these points, one therefore always falls back on the same points.\ednote{In (4) the print has $x = a + k'$ with no subscript; apparently for $x = a_{1} + k'$, either a misprint or a subscript that did not print.}

Let us apply the third substitution to them.

Let
\[ a_{2} = \alpha a_{1} + \beta b_{1}, \qquad b_{2} = \gamma a_{1} + \delta b_{1}, \]
so that the point $a_{2}$, $b_{2}$ is the transform of the point $a_{1}$, $b_{1}$ by $s$.

Then the transforms of the points (4) by the third substitution will be comprised in the formula
\[ x = a_{2} + k'', \qquad y = b_{2} + k''_{1}, \qquad z = c + 2, \tag{5} \]
$k''$ and $k''_{1}$ being two integers.

In general, let us suppose that the successive transforms of the point $(a, b)$ by the substitution $s$ are $a_{1}, b_{1}$; $a_{2}, b_{2}$; \ldots; $a_{n}, b_{n}$; \ldots\ and likewise that the successive transforms by the inverse substitution are $a_{-1}$, $b_{-1}$; $a_{-2}$, $b_{+2}$\ednote{Printed $b_{+2}$; apparently a misprint for $b_{-2}$.}; \ldots.

Then all the transforms of the point $(a, b, c)$ by the substitutions of the group (1) will be given by the formula
\[ x = a_{n} + k, \qquad y = b_{n} + k_{1}, \qquad z = c + n, \]
where $n$, $k$ and $k_{1}$ are three arbitrary integers.

One easily sees, moreover, that the group derived from the first two substitutions is permutable with the third.

The fundamental domain $D_{0}$ is a cube whose side is equal to 1, bounded by the six planes
\[ x; y; z = 0; 1. \]
\origpage{59}

The simplest case is that in which
\[ \alpha = \delta = 1, \qquad \beta = \gamma = 0, \]
because then our three substitutions reduce to
\[ (x, y, z; x + 1, y, z; x, y + 1, z; x, y, z + 1). \]

Each of them changes one of the faces of the cube into the opposite face; we therefore fall back quite simply on our first example.

But the way in which the surface of the cube $D_{0}$ is subdivided into conjugate regions is not always so simple.

Let us suppose, for example,
\[ \alpha = \beta = \delta = 1, \qquad \gamma = 0; \]
each of the faces parallel to the axis of $z$ will still be conjugate to the opposite face; but for the faces $z = 0$, $z = 1$ perpendicular to the axis of $z$, it will be more complicated.

Let us suppose that the points $ABCD$, $A'B'C'D'$ have the same coordinates as in our first four examples. Each of the faces $ABCD$, $A'B'C'D'$ will have to be divided into two triangles, namely: $ABC$ and $BCD$ on the one hand, $A'D'C'$ and $A'B'D'$ on the other, and the law of conjugation of the faces will be expressed by the relations
\[ \begin{aligned} &ACC'A' \equiv BDD'A',\\ &CDD'C' \equiv ABB'A',\\ &ABC \equiv A'D'C',\\ &BCD \equiv B'A'D'. \end{aligned} \]
\ednote{The first relation prints $BDD'A'$; apparently for $BDD'B'$, the face opposite $ACC'A'$.}

More generally, the faces parallel to the axis of $z$ will remain conjugate in pairs, but the faces perpendicular to the axis of $z$ will have to be decomposed into smaller polygons, the more numerous the larger the numbers $\alpha$, $\beta$, $\gamma$, $\delta$, and which will be conjugate in pairs according to a more or less complicated law.
\origpage{60}

A simple case is this:
\[ \alpha = \delta = 0, \qquad \beta = 1, \qquad \gamma = -1. \]

In this case the mode of conjugation is the same as in our fourth example.

\subsection*{\S~12. --- Fundamental group.}

We are thus led to the notion of the fundamental group of a manifold.

Let
\[ F_{1}, \qquad F_{2}, \qquad \ldots, \qquad F_{\lambda}, \]
be $\lambda$ functions of the coordinates $x_{1}, x_{2}, \ldots, x_{n}$ of a point $M$ of a manifold $V$ defined by the relations
\[ f_{\alpha} = 0, \qquad \varphi_{\beta} > 0. \]

I do not suppose that the functions $F$ are uniform. When the point $M$, starting from its initial position $M_{0}$, returns to this position after having traversed any path, it may be that the functions $F$ do not return to their original values.

But to fix ideas better, and although this is in no way essential, let us suppose that the functions $F$ are defined in the following manner. They must satisfy differential equations of the form
\[ dF_{i} = X_{i.1}\,dx_{1} + X_{i.2}\,dx_{2} + \ldots + X_{i.n}\,dx_{n}, \tag{1} \]
where the coefficients $X_{i.k}$ will be given functions of the $x_{k}$ and of the $F_{i}$. These functions must be uniform, finite and continuous, as must their derivatives, for all values of the $F$ as well as for all the points of $V$, and even for all the points sufficiently near to $V$.

I suppose likewise that, for all the points sufficiently near to $V$, the equations (1) satisfy the conditions of integrability, which
\origpage{61}
can be written
\[ \begin{aligned} &\frac{dX_{i.k}}{dx_{q}} + \frac{dX_{i.k}}{dF_{1}}X_{1.q} + \frac{dX_{i.k}}{dF_{2}}X_{2.q} + \ldots + \frac{dX_{i.k}}{dF_{\lambda}}X_{\lambda.k}\\ &= \frac{dX_{i.q}}{dx_{k}} + \frac{dX_{i.q}}{dF_{1}}X_{1.k} + \frac{dX_{i.q}}{dF_{2}}X_{2.k} + \ldots + \frac{dX_{i.q}}{dF_{k}}X_{\lambda.k}. \end{aligned} \]\ednote{The last term of the first line is printed $X_{\lambda.k}$, apparently for $X_{\lambda.q}$; the last denominator of the second line is printed $dF_{k}$, apparently for $dF_{\lambda}$. Both are reproduced as printed.}

If then the point $M$ describes on the manifold $V$ an infinitely small contour, the functions $F$ will return to their original values. The same will again be true if the point $M$ describes on the manifold $V$ a \emph{loop}, that is to say if it goes first from $M_{0}$ to $M_{1}$ by any path $M_{0}BM_{1}$, then describes an infinitely small contour, and finally returns from $M_{1}$ to $M_{0}$ by the \emph{same} path $M_{1}BM_{0}$.

But the same may no longer be true if it describes a finite closed contour.

Let us suppose, for example, that we are in ordinary space and that our manifold $V$ is a torus. One can evidently imagine that the functions $F$ return to their original values when the point $M$ describes a loop on this torus, but that the same is no longer true if the point $M$ describes a meridian circumference or a parallel.

The final values of the functions $F$, when the moving point $M$, starting from an initial point $M_{0}$, returns to it after having described a closed contour; these final values, I say, will naturally depend on the initial values.

Let, then, $F^{0}_{a}$ and $F^{1}_{a}$ be the initial and final values of $F_{a}$. The $F^{1}_{a}$ will be functions of the $F^{0}_{a}$, or, in other words, the functions $F$ will undergo, when the point $M$ describes the closed contour considered, a certain substitution, which will change the $F^{0}_{a}$ into $F^{1}_{a}$.

\emph{The set of all the substitutions that the functions $F$ thus undergo, when the point $M$ describes all the closed contours that can be traced on the manifold $V$ starting from the initial point $M_{0}$, will evidently form a group, which I call $g$}.

Now let us consider a closed contour $M_{0}BM_{0}$ traced on $V$ and starting from the initial point $M_{0}$. If this closed contour reduces to a loop, I shall agree to write
\origpage{62}
\[ M_{0}BM_{0} \equiv 0. \]\ednote{Only two strokes of the sign are printed here, so it reads as $=$; the same relation is written $M_{0}BM_{0} \equiv 0$ below and on p. 63, and a stroke of the $\equiv$ apparently did not print.}

If now $M_{0}AM_{1}$, $M_{0}BM_{1}$, $M_{0}CM_{1}$ are three different paths traced on $V$ and going from $M_{0}$ to $M_{1}$, I shall agree to write
\[ M_{0}AM_{1}CM_{0} \equiv M_{0}AM_{1}BM_{0} + M_{0}BM_{1}CM_{0}. \]

It is important to remark that $M_{0}AM_{1}CM_{0}$ is not the same thing as $M_{0}CM_{1}AM_{0}$, nor $M_{0}AM_{1}BM_{0} + M_{0}BM_{1}CM_{0}$ the same thing as
\[ M_{0}BM_{1}CM_{0} + M_{0}AM_{1}BM_{0}; \]
one cannot invert the order of the terms of a sum.

It follows from this convention that one will have
\[ M_{0}BM_{0} \equiv 0, \]
if the closed contour $M_{0}BM_{0}$ constitutes the complete boundary of a manifold of\ednote{Printed ``et''; apparently for ``à'', a variety à deux dimensions.} two dimensions forming part of $V$; and indeed this closed contour can then be decomposed into a very great number of loops.

We are thus led to consider relations of the form
\[ k_{1}C_{1} + k_{2}C_{2} \equiv k_{3}C_{3} + k_{4}C_{4}, \]
where the $k$ are integers and the $C$ closed contours traced on $V$ and starting from $M_{0}$. These relations, which I shall call \emph{equivalences}, resemble the homologies that I studied above. They differ from them:

1° Because, in homologies, the contours may start from any initial point.

2° Because, in homologies, one has the right to invert the order of the terms of a sum.

One will be able to add two equivalences to one another, but on condition of not inverting the order of the terms; if, then, one has
\[ A \equiv B \]
\origpage{63}
and
\[ C \equiv D, \]
one will be able to conclude from this
\[ A + C \equiv B + D, \]
but not
\[ C + A \equiv B + D. \]
Let us add that\ednote{The word ``que'' is printed twice; an apparent misprint.} from
\[ 2A \equiv 0 \]
one has not the right to conclude
\[ A \equiv 0. \]

This being so, it is clear that one can imagine a group $G$ satisfying the following conditions:

1° To each closed contour $M_{0}BM_{0}$ there will correspond a substitution $S$ of the group;

2° The necessary and sufficient condition for $S$ to reduce to the identical substitution is that
\[ M_{0}BM_{0} \equiv 0; \]\ednote{The strokes of the $\equiv$ are only partly printed here.}

3° If $S$ and $S'$ correspond to the contours $C$ and $C'$, and if
\[ C'' \equiv C + C', \]
the substitution corresponding to $C''$ will be $SS'$.

The group $G$ will be called the \emph{fundamental group} of the manifold $V$.

Let us compare it with the group $g$ of the substitutions undergone by the functions $F$.

The group $g$ will be isomorphic to $G$.

The isomorphism may be holohedral.

It may be merihedral if a closed contour $M_{0}BM_{0}$ not decomposable into loops brings the functions $F$ back to their original values.

\subsection*{\S~13. --- Fundamental equivalences.}

The group $G$ will be derived from a certain number of principal substitutions
\origpage{64}
$S_{1}, S_{2}, \ldots, S_{p}$. To each of them there will correspond a closed contour, so that we shall have \emph{$p$ fundamental closed contours} $C_{1}, C_{2}, \ldots, C_{p}$, and any closed contour will be equivalent to a combination of the fundamental contours following one another in a certain order.

These fundamental contours are not in general altogether independent, and there are between them certain relations which I shall call \emph{fundamental equivalences}.

Let us suppose, for example, that one has between the fundamental contours the equivalence
\[ k_{1}C_{1} + k_{2}C_{2} + k'_{1}C_{1} + k_{3}C_{3} \equiv 0. \]

This signifies that the substitution
\[ S^{k_{1}}_{1}S^{k_{2}}_{2}S^{k'_{1}}_{1}S^{k_{3}}_{3} \]
of the group $G$ reduces to the identical substitution.

The fundamental equivalences therefore make known to us the form of the group $G$.

Let us suppose that the manifold $V$ has been defined by the mode of representation of \no 10, and that we have only a single polyhedron $P_{1}$. It is clear that one will obtain all the fundamental contours in the following manner. Let $M_{0}$ be a point interior to $P_{1}$, $A$ a point of one of the faces of $P_{1}$, $A'$ the corresponding point on the conjugate face. One will go from $M_{0}$ to $A$, then from $A'$ to $M_{0}$ without leaving $P_{1}$; the corresponding path on the manifold $V$ will be closed.

There will therefore be as many fundamental contours as pairs of faces.

Here now is how one will form the fundamental equivalences.

Let us consider a cycle of edges. Let there be, for example, an edge, the intersection of the faces $F_{1}$ and $F'_{\mu}$, which for this reason I shall designate by the name of \emph{edge} $F_{1}F'_{\mu}$; let $F'_{1}$ be the face conjugate to $F_{1}$ and $F_{2}$, $F'_{1}$\ednote{A comma is printed between $F_{2}$ and $F'_{1}$; the edge is named $F_{2}F'_{1}$ two lines below.} the edge corresponding to $F_{1}F'_{\mu}$ on this face; let $F'_{2}$ be the face conjugate to $F_{1}$\ednote{Printed $F_{1}$; apparently for $F_{2}$, since $F'_{1}$ is already the face conjugate to $F_{1}$.} and $F_{3}F'_{2}$ the edge corresponding to $F_{2}F'_{1}$ on this face; and so on until one falls back on the face $F'_{\mu}$ and the edge $F_{1}F'_{\mu}$.
\origpage{65}

Let us remark that in carrying out this operation one may fall back several times on the same face.

Let $A_{i}$ be a point of $F_{i}$ and $A'_{i}$ the corresponding point of $F'_{i}$; let $C_{i}$ be the fundamental contour
\[ M_{0}A_{i} + A'_{i}M_{0}. \]

Then one will have the fundamental equivalence
\[ C_{1} + C_{2} + \ldots + C_{\mu} \equiv 0. \]

There will therefore be as many fundamental equivalences as cycles of edges.

When one has thus formed the fundamental equivalences, one will deduce from them the fundamental homologies, which differ from them only in that the order of the terms is indifferent. The knowledge of these homologies will immediately make known the Betti number $P_{1}$.

Let us apply these principles to the examples cited above, and let us remark that, all the manifolds cited in these examples being closed and of three dimensions, one will have
\[ P_{1} = P_{2}. \]

\subsection*{First example.}

\[ \begin{aligned} &(C_{1}) &&ABDC \equiv A'B'D'C',\\ &(C_{2}) &&ABB'A' \equiv CDD'C',\\ &(C_{3}) &&ACC'A' \equiv BDD'B'. \end{aligned} \]

Here is what I mean by the notation
\[ \begin{aligned} &(C_{1}) &&ABDC \equiv A'B'D'C'. \end{aligned} \]

I mean that the face $ABDC$ is conjugate to $A'B'D'C'$ and that, $\alpha$ denoting a point of $ABDC$ and $\alpha'$ a point of $A'B'D'C'$, the fundamental contour $M_{0}\alpha + \alpha'M_{0}$ is denoted by $C_{1}$.
\origpage{66}

\subsection*{Fundamental equivalences.}

\[ C_{1} + C_{2} \equiv C_{2} + C_{1}, \qquad C_{1} + C_{3} \equiv C_{3} + C_{1}, \qquad C_{2} + C_{3} \equiv C_{3} + C_{2}. \]

The fundamental homologies reduce to identities
\[ P_{1} = P_{2} = 4. \]

I pass at once to the third example, since we have seen that the second must be rejected.

\subsection*{Third example.}

\[ \begin{aligned} &(C_{1}) &&ABDC \equiv B'D'C'A',\\ &(C_{2}) &&ABB'A' \equiv DD'C'C,\\ &(C_{2}) &&ACC'A' \equiv DD'B'B. \end{aligned} \]
\ednote{The third row is labelled $C_{2}$ in the print, repeating the second label; apparently for $C_{3}$.}

\subsection*{Fundamental equivalences.}

\[ \begin{aligned} &C_{1} + C_{3} + C_{2} \equiv 0, &&C_{1} - C_{3} - C_{2} \equiv 0,\\ &C_{2} - C_{1} - C_{3} \equiv 0, &&C_{3} - C_{2} - C_{1} \equiv 0, \end{aligned} \]
\ednote{The second minus sign in $C_{3} - C_{2} - C_{1} \equiv 0$ did not print; a blank gap stands in its place.}
which can also be written
\[ 2C_{1} \equiv 2C_{2} \equiv 2C_{3}, \qquad 4C_{1} = 0. \]

\subsection*{Fundamental homologies.}

\[ \begin{aligned} &C_{1} \backsim C_{2} \backsim C_{3} \backsim 0,\\ &P_{1} = P_{2} = 1. \end{aligned} \]

One can give a simple geometric interpretation of this result.

The group $G$ is of finite order and consists only of eight
\origpage{67}
distinct substitutions corresponding to the following contours:
\[ 0, \qquad C_{1}, \qquad C_{2}, \qquad C_{3}, \qquad 2C_{1}, \qquad 3C_{1}, \qquad 3C_{2}, \qquad 3C_{1}. \]
\ednote{The last term reads $3C_{1}$ in the print, repeating the sixth; apparently for $3C_{3}$. Its subscript is faintly printed.}

The group is isomorphic to the following group
\[ \begin{aligned} &(x, y, z, t; -y, x, -t, z;\\ &-t, -z, y, x; z, -t, -x, y;\\ &-x, -y, -z, -t; y, -x, t. -z;\\ &t, z, -y, -x; -z, t, x, -y). \end{aligned} \]
\ednote{The print sets a period after $t$ in the third line, $t. -z$; apparently for a comma.}

This group, which transforms into itself the hypercube of four dimensions whose eight faces have as equations
\[ x = \pm 1, \qquad y = \pm 1, \qquad z = \pm 1, \qquad t = \pm 1, \]
\ednote{The last $\pm$ is imperfectly printed; its strokes are broken apart.}
could be called the \emph{hypercubic} group.

\subsection*{Fourth example.}

\[ \begin{aligned} &(C_{1}) &&ABDC \equiv B'D'C'A',\\ &(C_{2}) &&ABB'A' \equiv CDD'C',\\ &(C_{3}) &&ACC'A' \equiv BDD'B'. \end{aligned} \]

\subsection*{Fundamental equivalences.}

\[ C_{2} + C_{3} \equiv C_{3} + C_{2}, \qquad C_{3} + C_{1} \equiv C_{1} + C_{2}, \qquad -C_{2} + C_{1} \equiv C_{1} + C_{3}. \]

\subsection*{Fundamental homologies.}

\[ C_{2} \backsim C_{3} \backsim 0, \]
whence
\[ P_{1} = P_{2} = 2. \]
\origpage{68}
\subsection*{Fifth example.}

\[ \begin{aligned} &(C_{1}) &&ABC \equiv FED,\\ &(C_{2}) &&ACE \equiv FDB,\\ &(C_{3}) &&AED \equiv FBC,\\ &(C_{4}) &&ADB \equiv FCE. \end{aligned} \]

\subsection*{Fundamental equivalences.}

\[ C_{1} \equiv C_{2} \equiv C_{3} \equiv C_{4}, \qquad 2C_{1} \equiv 0, \]
whence
\[ C_{1} \backsim C_{2} \backsim C_{3} \backsim C_{4} \backsim 0, \]
and
\[ P_{1} = P_{2} = 1. \]

The group $G$ reduces to two substitutions, corresponding to the contours
\[ 0 \qquad \text{and} \qquad C_{1}. \]

\subsection*{Sixth example.}

The group $G$ is evidently holomorphic to the group $(\alpha, \beta, \gamma, \delta)$.

The three substitutions
\[ \begin{aligned} &(x, y, z; x + 1, y, z),\\ &(x, y, z; x, y + 1, z),\\ &(x, y, z; \alpha x + \beta y, \gamma x + \delta y, z + 1) \end{aligned} \]
will correspond respectively to the fundamental contours $C_{1}$, $C_{2}$ and $C_{3}$.

We have first the fundamental equivalences
\[ \begin{aligned} &C_{1} + C_{2} \equiv C_{2} + C_{1},\\ &C_{1} + C_{3} \equiv C_{3} + \alpha C_{1} + \gamma C_{2},\\ &C_{2} + C_{3} \equiv C_{1} + \beta C_{1} + \delta C_{2}, \end{aligned} \]
\ednote{The third row reads $C_{1} + \beta C_{1}$ on the right in the print; apparently for $C_{3} + \beta C_{1}$, as in the row above.}
\origpage{69}
from which it follows first that every combination of the fundamental contours can be put in the form
\[ m_{3}C_{3} + m_{1}C_{1} + m_{2}C_{2}, \]
the $m$ being integers. And since this expression could be equivalent to 0 only if the three integers $m$ were zero, it follows that we have here all the fundamental equivalences.

\subsection*{Fundamental homologies.}

\[ \begin{aligned} &(\alpha - 1)C_{1} + \gamma C_{2} \backsim 0,\\ &\beta C_{1} + (\delta - 1) C_{1} \backsim 0. \end{aligned} \]
\ednote{The second homology reads $(\delta - 1) C_{1}$ in the print; apparently for $(\delta - 1) C_{2}$.}

If these two homologies are not distinct, one will have
\[ C_{1} \backsim C_{2} \backsim 0; \]
whence
\[ P_{1} = P_{2} = 2; \]
this is what happens in the general case, and in particular for our fourth example.

If the determinant of these homologies is zero, that is to say if
\[ (\alpha - 1)(\delta - 1) - \beta\gamma = 0 \]
or
\[ \alpha + \delta = 2, \]
one will have
\[ P_{1} = P_{2} = 3, \]
except in the case where the two homologies reduce to identities. This is what happens for
\[ \alpha = \delta = 1, \qquad \beta = \gamma = 0, \]
that is to say for our first example; one then has
\[ P_{1} = P_{2} = 4. \]
\origpage{70}

\subsection*{\S~14. --- Conditions of homeomorphism.}

One knows that two closed manifolds of two dimensions that have the same Betti number are homeomorphic. This follows, for example, from the study of the periods of abelian functions. Let us consider a Riemann surface $R$ and let $z$ be the corresponding imaginary variable; one will be able to introduce a new imaginary variable $t$, such that $z$ is a Fuchsian function of $t$ and that $t$, considered as a function of $z$, has no singular point on the surface $R$. All the Fuchsian groups corresponding to various Riemann surfaces having the same order of connection will be isomorphic.

This Fuchsian group will, moreover, evidently be nothing other than the fundamental group $g$ relative to the surface $R$ considered as a manifold of two dimensions.

It will be remarked that not all Fuchsian groups are capable of defining in this way a closed manifold of two dimensions. Let us consider the fundamental Fuchsian polygon $R_{0}$, to which, if the Fuchsian function exists in the whole plane, one will have to adjoin its symmetric $R'_{0}$ with respect to the axis of real quantities; but then the domain $R_{0} + R'_{0}$ will no longer always be simply connected. To each point of the closed manifold $V$ there ought to correspond one point of $R_{0}$ (or of $R_{0} + R'_{0}$) and only one, and conversely. Let us suppose that there exist one or more cycles of vertices and that the sum of the angles of this cycle is zero or $\frac{2\pi}{n}$, $n$ being an integer greater than 1; let then $M$ be the point of $V$ that corresponds to this cycle of vertices, and let us consider an infinitely small loop enveloping $M$. By the definition of the group $g$, there must correspond to this loop, in the group $g$, the identical substitution and, in the Fuchsian group, a non-identical substitution. The Fuchsian group therefore cannot be isomorphic to $g$.

There then remain the Fuchsian groups of the first family such that the sum of the angles of any cycle is $2\pi$, and those of the third family. But these must likewise be rejected. Indeed, if the group is of the third family, the domain $R_{0} + R'_{0}$ is no longer simply
\origpage{71}
connected. Let $C$ be a closed contour traced in this domain and such that one does not have
\[ C \backsim 0. \]

To this contour there will correspond, in the Fuchsian group, an identical substitution (since the variable $z$ will have returned to its point of departure) and, in the group $g$, a non-identical substitution. Here again the Fuchsian group cannot be isomorphic to $g$.

There remain therefore only the groups of the first family such that the sum of the angles of each cycle is equal to $2\pi$.

All those of these groups that are of the same genus are isomorphic to one another, and it is for this reason that all closed manifolds of two dimensions that have the same Betti number are homeomorphic.

Is the same true when the number of dimensions is greater; are two closed manifolds of $h$ dimensions ($h > 2$) that have the same Betti numbers homeomorphic?

We are going to see that they are not, and this is what shows us how much the questions of \emph{Analysis Situs} become complicated when the number of dimensions increases.

It is clear first that, if two manifolds are homeomorphic, their groups $g$ will be isomorphic.

Let us now return to our sixth example and inquire whether two groups $(\alpha, \beta, \gamma, \delta)$ can be isomorphic.

Let $(\alpha, \beta, \gamma, \delta)$ and $(\alpha', \beta', \gamma', \delta')$ be these two groups; let $C_{1}$, $C_{2}$, $C_{3}$; $C'_{1}$, $C'_{2}$, $C'_{3}$ be the three fundamental contours of each of these groups;
\[ \left\{ \begin{aligned} &C_{1} + C_{2} \equiv C_{2} + C_{1},\\ &C_{1} + C_{3} \equiv C_{3} + \alpha C_{1} + \gamma C_{2},\\ &C_{2} + C_{3} \equiv C_{3} + \beta C_{1} = \delta C_{2}, \end{aligned} \right. \tag{1} \]
\ednote{The third row prints $= \delta C_{2}$ at its end, with two clear strokes; apparently for $+ \delta C_{2}$, as in the corresponding row of (1 bis).}
\[ \left\{ \begin{aligned} &C'_{1} + C'_{2} \equiv C'_{2} + C'_{1},\\ &C'_{1} + C'_{3} \equiv C'_{3} + \alpha' C'_{1} + \gamma' C'_{2},\\ &C'_{2} + C'_{3} \equiv C'_{3} + \beta' C'_{1} + \delta' C'_{2}, \end{aligned} \right. \tag{1 \textit{bis}} \]
the fundamental equivalences of the two groups\ednote{No final period is printed after ``groupes''; the next page opens a new paragraph.}
\origpage{72}

Let us suppose the two groups isomorphic, and let
\[ \begin{aligned} &a_{3}C_{3} + a_{1}C_{1} + a_{2}C_{2},\\ &b_{3}C_{3} + b_{1}C_{1} + b_{2}C_{2},\\ &c_{3}C_{3} + c_{1}C_{1} + c_{2}C_{2} \end{aligned} \]
be the contours of the first group that correspond respectively to the contours $C'_{1}$, $C'_{2}$ and $C'_{3}$ of the second group; the $a$, the $b$ and the $c$ are integers; we have seen above, indeed, that every contour of the first group can be put in this form.

For the isomorphism to hold, it is necessary that, on substituting in the equivalences (1~\emph{bis}) $a_{3}C_{3} + a_{1}C_{1} + a_{2}C_{2} + \ldots$ in place of $C'_{1}$, $C'_{2}$, $C'_{3}$, one recovers the equivalences (1).

It is therefore necessary first that the substitutions (which I denote by the same signs as the corresponding contours)
\[ \begin{aligned} &a_{3}C_{3} + a_{1}C_{1} + a_{2}C_{2},\\ &b_{3}C_{3} + b_{1}C_{1} + b_{2}C_{2} \end{aligned} \]
be permutable. To simplify the writing, I am going to employ the following notation:

I shall put
\[ a_{3} = h, \qquad b_{3} = k, \qquad a_{1}C_{1} + a_{2}C_{2} = S_{0}; \qquad b_{1}C_{1} + b_{2}C_{2} = T_{0}, \]
so that our first two substitutions reduce to $hC_{3} + S_{0}$, $kC_{3} + T_{0}$.

I shall denote by $S_{1}$ what $S_{0}$ becomes when its coefficients are subjected to the linear substitution $(\alpha, \beta, \gamma, \delta)$, that is to say when one replaces $a_{1}$ and $a_{2}$ by
\[ \alpha a_{1} + \beta a_{2} \qquad \text{and} \qquad \gamma a_{1} + \delta a_{2}. \]

$S_{2}$ will be what $S_{1}$ becomes when its coefficients are subjected to the same substitution, and so on; likewise $T_{1}$, $T_{2}$, \ldots, will be the successive transforms of $T_{0}$.
\origpage{73}

This being so, one will have, by virtue of the equivalences (1),
\[ S_{0} + hC_{3} = hC_{3} + S_{h}. \]

For our substitutions to be permutable, it is therefore necessary that
\[ hC_{3} + S_{0} + kC_{3} + T_{0} = kC_{3} + T_{0} + hC_{3} + S_{0} \]
or
\[ (k + h)C_{3} + S_{k} + T_{0} \equiv (k + h)C_{3} + T_{h} + S_{0} \]\ednote{The subscript of the second $C_{3}$ prints like $1$; apparently a damaged $3$.}
or
\[ S_{k} + T_{0} \equiv T_{h} + S_{0}. \tag{$\lambda$} \]

Let us suppose first that $h$ is equal to $k$ and different from 0; the preceding equivalence can be replaced by the equalities
\[ \begin{aligned} &(\alpha_{h} - 1)(a_{1} - b_{1}) + \beta_{h}(a_{2} - b_{2}) = 0,\\ &\gamma_{h}(a_{1} - b_{1}) + (\delta_{h} - 1)(a_{2} - b_{2}) = 0, \end{aligned} \]
where I have denoted by
\[ \begin{vmatrix} \alpha_{h} & \beta_{h} \\[1ex] \gamma_{h} & \delta_{h} \end{vmatrix} \]
the coefficients of the $h^{\text{th}}$ power of the linear substitution
\[ \begin{vmatrix} \alpha & \beta \\[1ex] \gamma & \delta \end{vmatrix}. \]

These equalities can be satisfied in two ways:

1° Either if
\[ a_{1} = b_{1}, \qquad a_{2} = b_{2}, \]
in which case the two substitutions
\[ hC_{3} + S_{0}, \qquad kC_{3} + T_{0} \]
would be identical;
\origpage{74}

2° Or if the determinant
\[ \begin{vmatrix} \alpha_{h} - 1 & \beta_{h} \\[1ex] \gamma_{h} & \delta_{h} - 1 \end{vmatrix} \]
is zero. But this can happen only if one has
\[ s^{h} = 1, \]
$s$ being one of the roots of the equation
\[ \begin{vmatrix} \alpha - s & \beta \\[1ex] \gamma & \delta - s \end{vmatrix} = 0, \tag{2} \]
that is to say if the substitution $(\alpha, \beta, \gamma, \delta)$ is \emph{elliptic}, which means that the roots of the equation (2) in $s$ are imaginary (in the case in question, they must be equal to an $h^{\text{th}}$ root of unity), or \emph{parabolic}, which means that the roots of the equation (2) are equal.

Let us therefore suppose that $(\alpha, \beta, \gamma, \delta)$ is \emph{hyperbolic}, that is to say that the roots of the equation in $s$ are real; and let us no longer suppose $h = k$.

Then the $k^{\text{th}}$ power of
\[ hC_{3} + S_{0} \]
and the $h^{\text{th}}$ power of
\[ kC_{3} + T_{0} \]
must be permutable; let
\[ \begin{aligned} &h'C_{3} + S'_{0},\\ &h'C_{3} + T'_{0} \end{aligned} \]\ednote{The second line prints $h'C_{3}$; apparently for $k'C_{3}$, since $h' = k'$ is stated just below.}
be these two powers; since one will have
\[ h' = k' = hk, \]
these two powers, by what precedes, would have to be identical. In the group $(\alpha', \beta', \gamma', \delta')$, supposed isomorphic, the $k^{\text{th}}$ power of $C'_{1}$
\origpage{75}
would have to be identical to the $h^{\text{th}}$ power of $C'_{2}$. Since this cannot be so, one must have
\[ h = k = 0, \]
that is to say
\[ a_{3} = b_{3} = 0. \]

Let us pass to the case of elliptic substitutions, with which we shall class the substitution
\[ \begin{vmatrix} -1 & 0 \\[1ex] 0 & -1 \end{vmatrix}. \]

Then the determinant
\[ \begin{vmatrix} \alpha_{h} - 1 & \beta_{h} \\[1ex] \gamma_{h} & \delta_{h} - 1 \end{vmatrix} \]
may vanish for a certain value of $h$, which I call $\nu$.

But then the substitution
\[ \begin{vmatrix} \alpha_{\nu} & \beta_{\nu} \\[1ex] \gamma_{\nu} & \delta_{\nu} \end{vmatrix} \]
reduces to the identical substitution, that is to say that one has
\[ \alpha_{\nu} = \delta_{\nu} = 1, \qquad \beta_{\nu} = \gamma_{\nu} = 0. \]

Then it happens first that $\nu C_{3}$ is permutable with $C_{1}$ and with $C_{2}$, and more generally that two substitutions
\[ \begin{aligned} &a_{3}C_{3} + a_{1}C_{1} + a_{2}C_{2},\\ &b_{3}C_{3} + b_{1}C_{1} + b_{2}C_{2} \end{aligned} \]
are permutable provided that $a_{3}$ and $b_{3}$ are divisible by $\nu$. This sufficient condition is moreover necessary, excluding, as we did above, the case where two powers of our two substitutions would be identical.

To prove it, I am going to employ a still more abbreviated notation; the symbol
\[ m_{1}C_{1} + m_{2}C_{2} \]
\origpage{76}
has a meaning only if $m_{1}$ and $m_{2}$ are integers; but the following
\[ \mu(a'_{1}C_{1} + a'_{2}C_{2}) + \rho(b'_{1}C_{1} + b'_{2}C_{2}) \]
may have a meaning even when $\mu$, $\rho$, the $a'$ and the $b'$ are not integers, provided that
\[ \mu a'_{1} + \rho b'_{1}, \qquad \mu a'_{1} + \rho b'_{2} \]\ednote{The second expression prints $\mu a'_{1}$; apparently for $\mu a'_{2}$.}
are integers. It is evident that what allows us to act thus is that $C_{1}$ and $C_{2}$ are permutable.

Let us then choose
\[ \begin{aligned} &\xi_{0} = a'_{1}C_{1} + a'_{2}C_{2},\\ &\eta_{0} = b'_{1}C_{1} + b'_{2}C_{2}, \end{aligned} \]
in such a way that
\[ \xi_{1} = \xi_{0}s, \qquad \eta_{1} = \eta_{0}s^{-1}; \]
$\xi_{1}$ and $\eta_{1}$ are, by our conventions, the transforms of $\xi_{0}$ and $\eta_{0}$ by the linear substitution $(\alpha, \beta, \gamma, \delta)$; $s$ is one of the roots of the equation (2). By the theory of linear substitutions, one can always choose the numbers $a'$ and $b'$ (which will generally be irrational and even imaginary) so that this is so; let us then put
\[ S_{0} = \mu\xi_{0} + \rho\eta_{0}, \qquad T_{0} = \mu'\xi_{0} + \rho'\eta_{0}. \]

One will have
\[ \mu\xi_{0} + \rho\eta_{0} + hC_{3} \equiv hC_{3} + \mu s^{h}\xi_{0} + \rho s^{-h}\eta_{0} \]
and
\[ k(hC_{3} + \mu\xi_{0} + \rho\eta_{0}) \equiv khC_{3} + \mu\frac{s^{kh} - 1}{s^{h} - 1}\xi_{0} + \rho\frac{s^{-kh} - 1}{s^{-h} - 1}\eta_{1}. \]\ednote{The last subscript prints like $1$; apparently for $\eta_{0}$.}

This being so, the equivalence ($\lambda$) can be written
\[ \begin{aligned} &\mu(s^{k} - 1) = \mu'(s^{h} - 1),\\ &\rho(s^{-k} - 1) = \rho'(s^{-h} - 1). \end{aligned} \]

If $k$ and $h$ are not divisible by $\nu$, these relations are not satisfied
\origpage{77}
identically, and we shall be able to put
\[ \mu = (s^{h} - 1)\varepsilon, \qquad \mu' = (s^{k} - 1)\varepsilon, \qquad \rho = (s^{-h} - 1)\zeta, \qquad \rho' = (s^{-k} - 1)\zeta. \]

But then
\[ \begin{aligned} &k(hC_{3} + S_{0}) \equiv khC_{3} + \varepsilon(s^{kh} - 1)\xi_{0} + \zeta(s^{-kh} - 1)\eta_{0},\\ &h(kC_{3} + T_{0}) \equiv khC_{3} + \varepsilon(s^{kh} - 1)\xi_{0} + \zeta(s^{-kh} - 1)\eta_{0}, \end{aligned} \]
which shows that the $k^{\text{th}}$ power of our first substitution is identical to the $h^{\text{th}}$ power of the second. Since we exclude this case, we conclude that one must have
\[ a_{3} \equiv b_{3} \equiv 0 \qquad (\text{mod.}\ \nu). \]

But one can go further; first, if we call $C'_{1}$ and $C_{2}$\ednote{Printed $C_{2}$ without a prime; apparently for $C'_{2}$.} our first two substitutions, we can replace them by
\[ \omega_{1}C'_{1} + \omega_{2}C'_{2}, \qquad \omega'_{1}C'_{1} + \omega'_{2}C'_{3}, \]\ednote{The last term prints $C'_{3}$ here and again in the parenthesis below; apparently for $C'_{2}$, as in the same expression on p. 79.}
the $\omega$ being integers such that $\omega_{1}\omega'_{2} - \omega_{2}\omega'_{1} = 1$. We can therefore always suppose that
\[ b_{3} = 0 \]
(otherwise one would replace $C'_{1}$ and $C'_{2}$ by $\omega_{1}C'_{1} + \omega_{2}C'_{2}$ and $\omega'_{1}C'_{1} + \omega'_{2}C'_{3}$, choosing the $\omega$ so as to annul the new $b_{3}$).

But the subgroup derived from $C'_{1}$ and $C'_{2}$, because of the isomorphism with $(\alpha', \beta', \gamma', \delta')$, must be permutable with all the substitutions of the group, and in particular with $C_{3}$.

Let us therefore write that the substitution
\[ -C_{3} + C'_{2} + C_{3} \]
forms part of the subgroup derived from $C'_{1}$ and $C'_{2}$. We have
\[ \begin{aligned} &C'_{2} = \mu'\xi_{0} + \rho'\eta_{0},\\ &-C_{3} + C'_{2} + C_{3} = \mu's\xi_{0} + \rho's^{-1}\eta_{0}. \end{aligned} \]
\origpage{78}

For this second substitution to form part of its group, it is necessary, if $a_{3}$ is not zero, that this last substitution be a multiple of $C'_{2}$.

Now this can evidently take place only if $s = s^{-1} = -1$.

If this case is set aside, one must have
\[ a_{3} = b_{3} = 0. \]

If, therefore, we set aside the case where
\[ \alpha + \delta = \pm 2 \]
[parabolic substitutions and the substitution $(-1, 0, 0, -1)$], we must have
\[ a_{3} = b_{3} = 0. \]

I add that $c_{3}$ must be equal to 1, otherwise the combinations of the three fundamental substitutions
\[ \begin{aligned} &a_{1}C_{1} + a_{2}C_{2} \equiv C'_{1},\\ &b_{1}C_{1} + b_{2}C_{2} \equiv C'_{2},\\ &c_{3}C_{3} + c_{1}C_{1} + c_{2}C_{2} \equiv C'_{3} \end{aligned} \]
could not reproduce all the substitutions
\[ m_{3}C_{3} + m_{1}C_{1} + m_{2}C_{2} \]
of the group $(\alpha, \beta, \gamma, \delta)$, but only those in which the integer $m_{3}$ is divisible by $c_{3}$.

Now every substitution of the group $(\alpha, \beta, \gamma, \delta)$ must be capable of being put in the form
\[ m_{3}C'_{3} + m_{1}C'_{1} + m_{2}C'_{2}. \]

For $C_{1}$ to be capable of being put in this form, it is necessary first that $m_{3}$ be zero; and next that one have identically
\[ C_{1} = m_{1}(a_{1}C_{1} + a_{2}C_{2}) + m_{2}(b_{1}C_{1} + b_{2}C_{2}), \]
\origpage{79}
and likewise
\[ C_{2} = m'_{1}(a_{1}C_{1} + a_{2}C_{2}) + m'_{2}(b_{1}C_{1} + b_{2}C_{2}). \]

Since the $m$ and the $m'$ are integers, I shall conclude from this that the determinant
\[ a_{1}b_{2} - a_{2}b_{1} = 1. \]

But I said above that one could replace $C'_{1}$ and $C'_{2}$ by
\[ \omega_{1}C'_{1} + \omega_{2}C'_{2}, \qquad \omega'_{1}C'_{1} + \omega'_{2}C'_{2}. \]

If $a_{1}b_{2} - a_{2}b_{1}$ is equal to 1, we shall be able to choose the $\omega$ in such a way that
\[ \omega_{1}C'_{1} + \omega_{2}C'_{2} = C_{1}, \qquad \omega'_{1}C'_{1} + \omega'_{2}C'_{2} = C_{2}, \]
that is to say that we shall always be able to suppose
\[ a_{1} = b_{2} = 1, \qquad a_{2} = b_{1} = 0. \]

But then the equivalence (1~\emph{bis}):
\[ C'_{1} + C'_{3} \equiv C'_{3} + \alpha'C'_{1} + \gamma'C'_{2} \]
becomes
\[ C_{1} + C_{3} \equiv C_{3} + \alpha'C_{1} + \gamma'C_{2}, \]
whence
\[ \alpha = \alpha', \qquad \gamma = \gamma'; \]
one would find likewise
\[ \beta = \beta', \qquad \delta = \delta'. \]

We must therefore conclude that the two groups $(\alpha, \beta, \gamma, \delta)$ and $(\alpha', \beta', \gamma', \delta')$ can be isomorphic only if one can pass from the one to the other by changing $C'_{1}$ and $C'_{2}$ into
\[ \omega_{1}C'_{1} + \omega_{2}C'_{2}, \qquad \omega'_{1}C'_{1} + \omega'_{2}C'_{2}. \]

This can be stated in another way.
\origpage{80}
Let
\[ S = \begin{vmatrix} \alpha & \beta \\[1ex] \gamma & \delta \end{vmatrix} \]
be a linear substitution with integer coefficients, such that
\[ \alpha\delta - \beta\gamma = 1. \]

Let
\[ T = \begin{vmatrix} \omega_{1} & \omega_{2} \\[1ex] \omega'_{1} & \omega'_{2} \end{vmatrix} \]
be another linear substitution with integer coefficients, such that
\[ \omega_{1}\omega'_{2} - \omega'_{1}\omega_{2} = 1. \]

The substitution
\[ T^{-1}ST, \]
which is called the \emph{transform} of $S$ by $T$, is also linear with integer coefficients and has determinant 1.

If two linear substitutions $S$ and $S'$ with integer coefficients and determinant 1 are transforms of one another by a substitution of the form $T$, I shall say that $S$ and $S'$ belong to the same class.

It is clear first that $S$ and $S'$ can belong to the same class only if the sum $\alpha + \delta$ has the same value for the one and for the other, but this condition is not sufficient, and to one and the same integer value of $\alpha + \delta$ there will correspond several classes of linear substitutions, in the same way as to one and the same determinant there correspond several classes of quadratic forms.

Thus the necessary and sufficient condition for the two groups $(\alpha, \beta, \gamma, \delta)$, $(\alpha', \beta', \gamma', \delta')$ to be isomorphic is that the two linear substitutions $(\alpha, \beta, \gamma, \delta)$, $(\alpha', \beta', \gamma', \delta')$ belong to the same class.
\origpage{81}

We have set aside the case where
\[ \alpha + \delta = \pm 2. \]

If $\alpha + \delta = 2$, the substitution $(\alpha, \beta, \gamma, \delta)$ will be of the same class as $(1, h, 0, 1)$; the group $(\alpha, \beta, \gamma, \delta)$ will be isomorphic to $(1, h, 0, 1)$.

The latter contains a remarkable substitution $C_{2}$, which is a multiple of no other and which is permutable with all the substitutions of the group. One sees moreover without difficulty that, if $h$ is not zero, $C_{2}$ is the only substitution that enjoys this property.

We may moreover set aside the case where $h$ is zero, for the group $(1, 0, 0, 1)$, all of whose substitutions are permutable in pairs, evidently cannot be isomorphic to any other group $(\alpha, \beta, \gamma, \delta)$.

If $\alpha + \delta = -2$, the group $(\alpha, \beta, \gamma, \delta)$ will be isomorphic to $(-1, h, 0, -1)$. The latter contains a remarkable substitution $C_{2}$, which is a multiple of no other, which is not permutable with all the substitutions of the group, but which is permutable with \emph{the double} of all these substitutions. If $h$ is not zero, $C_{2}$ is the only substitution that enjoys this property. If, on the contrary, $h$ is zero, there is an infinity of them, which already proves that $(-1, 0, 0, -1)$ cannot be isomorphic to $(-1, h, 0, -1)$.

Likewise, the presence of this remarkable substitution $C_{2}$ in $(1, h, 0, 1)$ and the absence of any substitution enjoying the same property in $(-1, h', 0, -1)$ show that these two groups cannot be isomorphic.

There remain for me two questions to resolve:

1° Can the groups $(1, h, 0, 1)$ and $(1, h', 0, 1)$, where\ednote{Printed ``ou'' without accent; apparently for ``où''.} $|h| \gtrless 0$, $|h'| \gtrless 0$, $|h| \gtrless |h'|$, be isomorphic?

2° The same question for the groups $(-1, h, 0, -1)$, $(-1, h', 0, -1)$.

Let us begin with the first question.

I observe first that, $C_{2}$ being the only substitution of the first group that enjoys the characteristic property stated above, one must have
\[ C'_{2} \equiv C_{2} \]
\origpage{82}
(or $C'_{2} \equiv -C_{2}$, but then one would change $C'_{1}$ and $C'_{2}$ into $-C'_{1}$ and $-C'_{2}$).

Now, for one to be able, by combining
\[ \begin{aligned} &C'_{1} \equiv a_{3}C_{3} + a_{1}C_{1} + a_{2}C_{2},\\ &C'_{2} \equiv C_{2},\\ &C'_{3} \equiv c_{3}C_{3} + c_{1}C_{1} + c_{2}C_{2}, \end{aligned} \]
to recover all the substitutions of the first group, it is necessary that
\[ a_{3}c_{1} - c_{3}a_{1} = 1. \]

One then easily proves that
\[ C'_{1} + C'_{3} \equiv C'_{3} + C'_{1} + hC'_{2}, \]
which proves that the two groups can be isomorphic only if
\[ h = \pm h'. \]

Let us pass to the second question.

One would see, as just now, that one must have
\[ \begin{aligned} &C'_{2} \equiv C_{2},\\ &a_{3}c_{1} - c_{3}a_{1} = 1 \end{aligned} \]

There next results
\[ C'_{1} + C'_{3} \equiv (a_{3} + c_{3})C_{3} + (c_{1} + a_{1}\varepsilon)C_{1} + (c_{2} + a_{2}\varepsilon - a_{1}c_{3}h\varepsilon)C_{2}, \]
where
\[ \varepsilon = (-1)^{c_{3}} \]
and
\[ \begin{aligned} &C'_{3} - C'_{1} \equiv (c_{3} - a_{3})C_{3} + (c_{1} - a_{1})\varepsilon'C_{1}\\ &+ [(c_{2} - a_{2})\varepsilon' + (a_{1} + c_{1})a_{3}\varepsilon'h]C_{2}, \end{aligned} \]
where
\[ \varepsilon' = (-1)^{a_{3}}. \]
\origpage{83}

If one wishes that
\[ C'_{1} + C'_{3} \equiv C'_{3} - C'_{1} + h'C'_{2}, \] \ednote{The minus sign in this display prints as two short broken strokes; read as a single $-$.}
it is necessary that
\[ a_{3} = 0, \qquad \varepsilon' = 1, \qquad \varepsilon = -1, \qquad a_{1}c_{3} = -1, \qquad h = h'. \]

Hence our two groups can be isomorphic only if
\[ h = \pm h'. \]

The result stated above is therefore general and extends to parabolic substitutions.

If the two linear substitutions $(\alpha, \beta, \gamma, \delta)$, $(\alpha', \beta', \gamma', \delta')$ are not of the same class, the two corresponding groups cannot be isomorphic.

It follows from this long discussion that the different groups $(\alpha, \beta, \gamma, \delta)$ can give rise to an infinity of distinct closed manifolds $V$, that is to say manifolds that are not homeomorphic. Now the number $P_{1}$ can take only one of the three values 2, 3 or 4.

\emph{For two closed manifolds to be homeomorphic, it is therefore not sufficient that they have the same Betti numbers}.

This is what our other examples likewise bring out.

In the third example the group $G$ reduces to eight substitutions, and in the fifth example to two substitutions only.

Let us consider, on the contrary, the hypersphere
\[ x^{2}_{1} + x^{2}_{2} + x^{2}_{3} + x^{2}_{4} = 1; \]
it is a manifold whose group $G$ reduces to a single substitution, the identical substitution.

Here, then, are three manifolds whose groups $G$ are of finite order; but these groups are not isomorphic, so that these manifolds will
\origpage{84}
not be homeomorphic; and yet they will have the same Betti numbers:
\[ P_{1} = P_{2} = 1. \]

It will seem natural to restrict the sense of the word \emph{simply connected} and to reserve it for manifolds whose group $G$ reduces to the identical substitution.

Then a closed manifold of more than two dimensions will be able to have its group $G$ of finite order without being simply connected.

This would not happen with manifolds of two dimensions: their group $G$ could not be of finite order without reducing to a single substitution.

We have seen why a Fuchsian group that admits a cycle of vertices whose sum of angles is equal to $\frac{2n}{n}$\ednote{Printed with an italic $n$ in the numerator; apparently for $\frac{2\pi}{n}$.} $(n > 1)$ cannot be the group $G$ of a closed manifold of two dimensions.

The same is true, for the same reason, of every group of finite order, or more generally of every group containing a substitution of which an integer power reduces to the identical substitution.

It might be interesting to treat the following questions:

1° Given a group $G$ defined by a certain number of fundamental equivalences, can it give rise to a closed manifold of $n$ dimensions?

2° How must one set about forming this manifold?

3° Are two manifolds of the same number of dimensions that have the same group $G$ always homeomorphic?

These questions would require difficult studies and long developments. I shall not speak of them here.

I wish nevertheless to draw attention to one point.

Riemann connected the study of algebraic curves with that of manifolds of two dimensions from the point of view of \emph{Analysis Situs}.

Likewise, the study of algebraic surfaces is connected with that of manifolds of four dimensions. These manifolds have three Betti numbers:
\[ P_{1} = P_{3} \qquad \text{and} \qquad P_{2}. \]
\origpage{85}

Mr.~Picard has shown that if the algebraic surface is the most general of its degree, the number $P_{1}$ reduces to 1; it takes a value greater than 1 only in certain particular cases. Multiple connection thus presented itself as a particular case of simple connection.

This result seems paradoxical; perhaps it will seem a little less so to us now; a group $G$ may be much more complex than another group $G'$ and yet correspond to a smaller value of the number $P_{1}$.

\subsection*{\S~15. --- Other modes of generation.}

One can give other definitions of manifolds which are, so to speak, intermediate between the first two.

For example, one can consider a manifold $V$ formed of the set of points that satisfy certain inequalities and the equations
\[ \left\{ \begin{aligned} &F_{1}(x_{1}, x_{2}, \ldots, x_{n}; y_{1}, y_{2}, \ldots, y_{q})\\ &F_{2}(x_{1}, x_{2}, \ldots, x_{n}; y_{1}, y_{2}, \ldots, y_{q})\\ &\ldots\ldots\ldots\ldots\ldots\ldots\\ &F_{p}(x_{1}, x_{2}, \ldots, x_{n}; y_{1}, y_{2}, \ldots, y_{q}) \end{aligned} \right\} = 0, \tag{1} \]
where the $x$ are coordinates of a point and the $y$ arbitrary parameters; the number of dimensions of $V$ is then $n - p + q$.

One can also adjoin to the equations (1) the relations
\[ \varphi_{\alpha}(y_{1}, y_{2}, \ldots, y_{q}) = 0 \qquad (\alpha = 1, 2, \ldots, \lambda) \]
between the $y$. The $y$ are then no longer entirely arbitrary parameters, and the number of dimensions of $V$ becomes equal to $n - p + q - \lambda$.

One can further consider a manifold $V$ defined by certain inequalities and by the relations
\[ x_{i} = \theta_{i}(y_{1}, y_{2}, \ldots, y_{p}) \qquad (i = 1, 2, \ldots, n), \tag{2} \]
\origpage{86}
the $y$ being parameters bound by $\lambda$ equations
\[ \varphi_{\alpha}(y_{1}, y_{2}, \ldots, y_{p}) = 0 \qquad (\alpha = 1, 2, \ldots, \lambda). \]

The number of dimensions of $V$ is then $p - \lambda$.

Let us consider for a moment $y_{1}, y_{2}, \ldots, y_{p}$ as the coordinates of a point $P$ in space of $p$ dimensions. The equalities
\[ \varphi_{\alpha} = 0 \]
will then define in this space of $p$ dimensions a certain manifold $W$.

To each point of $W$ there will correspond a point of $V$, since the equations (2) express the $x$ as functions of the $y$.

The simplest case is that in which, conversely, to a point of $V$ there corresponds a single point of $W$. But a case that is also very interesting is the following.

Let us suppose that the manifold $W$ remains unaltered when the $y$ undergo the substitutions of a certain group $G$. Let $P$ be a point of $W$, $P_{1}, P_{2}, \ldots, P_{h}$ the transforms of $P$ by the substitutions of $G$.

The set of the points $P, P_{1}, P_{2}, \ldots, P_{h}$ will form what I shall call a \emph{system of points}.

If the functions $\theta_{i}$ are not altered by the substitutions of $G$, it is clear that to the various points of one and the same system there will correspond one and the same point of $V$.

The interesting case is that in which to a point of $V$ there corresponds a single system of points of $W$.

Given a manifold $W$ and a group $G$ that does not alter it, one can always construct a manifold $V$ in such a way that to every point of $V$ there corresponds one system of points of $W$ and only one.

For $V$ to be two-sided, it is necessary and sufficient that the manifold $W$ be two-sided and that all the substitutions of $G$ enjoy the following property.

Let $y_{1}, y_{2}, \ldots, y_{p}$ be the coordinates of $P$, $y'_{1}, y'_{2}, \ldots, y'_{p}$ those of its transform; the functional determinant of the $y'$ with respect to the $y$ must be positive.
\origpage{87}

\subsection*{Seventh example.}

Let
\[ y^{2}_{1} + y^{2}_{2} + y^{2}_{3} = 1 \]
be the equation of the manifold $W$, which will thus be a sphere in ordinary space.

This sphere is not altered when one changes $y_{1}, y_{2}, y_{3}$ into $-y_{1}$, $-y_{2}$, $-y_{3}$. This will be our group $G$.

If we then put, for example,
\[ \begin{aligned} &x_{1} = y^{2}_{1}, &&x_{2} = y^{2}_{2}, &&x_{3} = y^{2}_{3},\\ &x_{4} = y_{2}y_{3}, &&x_{5} = y_{1}y_{3}, &&x_{6} = y_{1}, y_{2}, \end{aligned} \]
\ednote{Printed with a comma in $x_{6} = y_{1}, y_{2}$; apparently for $x_{6} = y_{1}y_{2}$.}
the $x$ will not change when the $y$ change sign, and we shall have defined a manifold $V$ of two dimensions in space of 6 dimensions.

This manifold will be closed; it will be \emph{one-sided}.

Let, indeed, $P$ be a point of $W$, whose coordinates will be $y_{1}$, $y_{2}$, $y_{3}$. To define the position of this point on the sphere $W$, it will suffice to know two of these coordinates, for example $y_{1}$ and $y_{2}$; for the equation of the sphere makes known to us $y_{3}$ as a function of $y_{1}$ and $y_{2}$.

Its transform $P'$, whose coordinates are $-y_{1}$, $-y_{2}$ and $-y_{3}$, is diametrically opposite. But one must not use $y_{1}$ and $y_{2}$ to define the position of the point $P$, because $y_{3}$ is not a uniform function of these two variables. It is better to put
\[ \begin{aligned} &y_{1} = \cos \varphi \sin \theta,\\ &y_{2} = \cos \varphi \sin \theta,\\ &y_{3} = \cos \theta. \end{aligned} \]
\ednote{The second line is printed $y_{2} = \cos \varphi \sin \theta$, identical to the first; apparently for $y_{2} = \sin \varphi \sin \theta$.}

The coordinates of the point $P$ in this new system are $\varphi$ and $\theta$; those of $P'$ are $\varphi + \pi$ and $\pi - \theta$. One then sees that
\[ \frac{\partial(\varphi + \pi, \pi - \theta)}{\partial(\varphi, \theta)} = -1 < 0. \]

The manifold $V$ is therefore one-sided.
\origpage{88}

\subsection*{Eighth example.}

Let $y_{1}, y_{2}, \ldots, y_{q}; z_{1}, z_{2}, \ldots, z_{q}$ be $2q$ parameters bound by the relations
\[ \left\{ \begin{aligned} &y^{2}_{1} + y^{2}_{2} + \ldots + y^{2}_{q} = 1,\\ &z^{2}_{1} + z^{2}_{2} + \ldots + z^{2}_{q} = 1. \end{aligned} \right. \tag{3} \]

If we regard these $2q$ parameters as the coordinates of a point in space of $2q$ dimensions, these equations (3) will represent a closed manifold $W$ of $2q - 2$ dimensions.

If we regard $y_{1}, y_{2}, \ldots, y_{q}$ and $z_{1}, z_{2}, \ldots, z_{q}$ as the coordinates of two points $Q$ and $Q'$ in space of $q$ dimensions, these two points must both lie on the \emph{hypersphere} $S$ which has as equation
\[ y^{2}_{1} + y^{2}_{2} + \ldots + y^{2}_{q} = 1, \]
and which is a closed manifold of $q - 1$ dimensions.

To each pair of points of $S$ there will thus correspond one point of $W$ and only one, and conversely, \emph{provided that I agree to regard the two pairs of points $QQ'$ and $Q'Q$ as distinct}.

Let us now consider the $\frac{q(q + 3)}{2}$ combinations
\[ y_{i} + z_{i}, \qquad y_{i}z_{i}, \qquad y_{i}z_{k} + y_{k}z_{i} \qquad (i, k = 1, 2, \ldots, q), \]
and let us equate them to $\frac{q(q + 3)}{2} = n$ variables
\[ x_{1}, \qquad x_{2}, \qquad \ldots, \qquad x_{n}. \]

We shall thus have defined a manifold $V$ of $2q - 2$ dimensions in space of $n$ dimensions.

When one changes $y_{i}$ into $z_{i}$ and $z_{i}$ into $y_{i}$ (that is to say when one permutes the two points $Q$ and $Q'$), these $\frac{q(q + 3)}{2}$ combinations do not change.
\origpage{89}

To each pair of points of the hypersphere there thus corresponds one point and only one of the manifold $V$, and conversely, but \emph{on condition of not regarding the two pairs $QQ'$ and $Q'Q$ as distinct}.

Is this manifold $V$ closed? I am going to show that it is not for $q = 2$ and that it is for $q > 2$.

One then has
\[ \begin{aligned} &x_{1} = y_{1} + z_{1}, &&x_{2} = y_{1}z_{1}, &&x_{3} = y_{2} + z_{2},\\ &x_{4} = y_{2}z_{2}, &&x_{5} = y_{1}z_{2} + y_{2}z_{1}. \end{aligned} \]

One then recognizes that, for the $y$ and the $z$ to be real, one must have
\[ x^{2}_{1} > 4x_{2}, \qquad x^{2}_{3} > 4x_{4}. \]

One would have analogous inequalities for $q > 2$; but in which cases do these inequalities define a true boundary for our manifold $V$?

To account for this better, I am going first to treat a simpler example.

Let there be first, in ordinary space, the circle
\[ x^{2} + y^{2} = 1, \qquad z = 0. \]

If one agrees to keep only the points of this line for which $y$ is positive, we shall then have the following relations:
\[ x^{2} + y^{2} = 1, \qquad z = 0, \qquad y > 0, \]
which will define a manifold of one dimension (which, in this case, is a semicircumference).

This manifold is not closed; it admits two boundary points:
\[ x = \pm 1, \qquad y = z = 0. \]

Let us consider, on the contrary, the following surface:
\[ x^{2} + y^{2} - z^{2} + (x^{2} + y^{2} + z^{2})^{2} = 0. \tag{4} \]
\origpage{90}

It is the surface generated by the revolution of a lemniscate about its major axis. It consists of two distinct sheets $N_{1}$ and $N_{2}$ having a common conical point, which is the origin; one can go from one sheet to the other only by passing through the origin. Let us then adjoin to the equation (4) the inequality
\[ z > 0. \tag{5} \]

The relations (4) and (5) define a manifold of two dimensions which is nothing other than the sheet $N_{1}$. This manifold must be regarded as closed; it is homeomorphic to a sphere; one must not regard the conical point
\[ x = y = z = 0 \]
as a true boundary point.

In general, if a manifold of $p$ dimensions is not closed, its boundary will be formed by one or more manifolds of $p - 1$ dimensions. If the set of points that one may suspect of being boundary points forms one or more manifolds of fewer than $p - 1$ dimensions, then they are not true boundary points from the point of view that concerns us at present, and the given manifold is closed.

Now one will obtain here the points that one may suspect of being boundary points when one supposes that the points $Q$ and $Q'$ coincide, that is to say that
\[ y_{1} = z_{1}, \qquad y_{2} = z_{2}, \qquad \ldots, \qquad y_{q} = z_{q}. \]

One will thus obtain a manifold of $q - 1$ dimensions. Thus the complete boundary of $V$ will have only $q - 1$ dimensions, while $V$ has $2q - 2$. $V$ will therefore be closed, unless
\[ 2q - 2 = (q - 1) + 1 \]
or
\[ q = 2. \]
\origpage{91}

To account for this better, let us compare the two examples $q = 2$ and $q = 3$.

Let first $q = 2$, and let us consider our manifold in the neighbourhood of the point
\[ y_{1} = z_{1} = 0, \qquad y_{2} = z_{2} = 1, \]
that is to say of the point
\[ x_{1} = 0, \qquad x_{2} = 0, \qquad x_{3} = 2, \qquad x_{4} = 1, \qquad x_{5} = 0. \]
Let us observe first that, for small values of $x_{1}$ and $x_{2}$, the three other variables $x_{3}$, $x_{4}$ and $x_{5}$ can be expanded in powers of $x_{1}$ and $x_{2}$; it will therefore suffice for us to study the variations of $x_{1}$ and $x_{2}$.

We then see that $x_{1}$ and $x_{2}$ can take all the values that satisfy
\[ x^{2}_{1} > 4x_{2}, \]

The plane of $x_{1}, x_{2}$ is thus divided into two regions by the line
\[ x^{2}_{1} = 4x_{2}, \]
which is a true boundary line.

One would obtain the same result if one studied the manifold $V$ in the neighbourhood of any other boundary point. This manifold is therefore not closed.

Let now $q = 3$ and
\[ \begin{aligned} &x_{1} = y_{1} + z_{1}, &&x_{2} = y_{1}z_{1}, &&x_{3} = y_{2} + z_{2}, &&x_{4} = y_{2}z_{2}, &&x_{5} = y_{1}z_{2} + y_{2}z_{1},\\ &x_{6} = y_{3} + z_{3}, &&x_{7} = y_{3}z_{3}, &&x_{8} = y_{1}z_{3} + y_{3}z_{1}, &&x_{9} = y_{2}z_{3} + y_{3}z^{2}. \end{aligned} \]\ednote{The last term is printed $y_{3}z^{2}$, with the 2 set as an exponent; apparently a misprint for $y_{3}z_{2}$, matching the pattern of $x_{5}$ and $x_{8}$.}

Let us study the manifold $V$ in the neighbourhood of the point $P_{0}$ which is such that
\[ y_{1} = z_{1} = y_{2} = z_{2} = 0, \qquad y_{3} = z_{3} = 1, \]
\origpage{92}
whence
\[ x_{1} = x_{2} = x_{3} = x_{4} = x_{5} = x_{8} = x_{9} = 0, \qquad x_{6} = 2, \qquad x_{7} = 1. \]\ednote{Between $x_{2}$ and $x_{3}$ only one short stroke of the $=$ printed.}

We see that, in the neighbourhood of this point, $x_{6}$, $x_{7}$, $x_{8}$, $x_{9}$ can be expanded in powers of $x_{1}$, $x_{2}$, $x_{3}$, $x_{4}$ and $x_{5}$, so that it suffices to study the variations of these last five variables.

To keep only three variables and make a geometric representation possible, I cut my manifold $V$ by the plane manifold
\[ x_{1} = 0, \qquad x_{3} = 0, \]
so that the intersection will be a manifold of two dimensions $V'$.

Let $x_{2}$, $x_{4}$ and $x_{5}$ be the coordinates of a point of $V'$; we shall be able to regard these three variables as the coordinates of a point in ordinary space, and we shall thus have a geometric representation of the manifold $V'$, or rather of the portion of this manifold that is near $P_{0}$.

One then finds
\[ y_{1} = -z_{1}, \qquad y_{2} = -z_{2}, \]
since $x_{1}$ and $x_{3}$ are supposed zero, and consequently
\[ x_{2} = -y^{2}_{1}, \qquad x_{4} = -y^{2}_{2}, \qquad x_{5} = -2y_{1}y_{2}, \]
whence
\[ 4x_{2}x_{4} - x^{2}_{5} = 0. \]

This is the equation of a cone of the second degree, but only one sheet of this cone is suitable, because one must have
\[ x_{2} < 0, \qquad x_{4} < 0. \]

The portion of the cone that is suitable is thus separated from the one that is not by the vertex, which cannot be regarded as a true boundary point. Thus the manifold $V'$, like $V$, is still closed.
\origpage{93}

One would obtain an analogous result by studying $V$ in the neighbourhood of another boundary point, or by cutting $V$ by other plane manifolds than the manifold
\[ x_{1} = x_{2} = 0. \]\ednote{Printed $x_{1} = x_{2} = 0$; apparently for $x_{1} = x_{3} = 0$, the plane variety used on p. 92.}

I wished only to give an example intended to clarify the preceding reasoning.

In sum, \emph{the manifold $V$ is closed if $q > 2$ and is not if $q = 2$}.

Is\ednote{Printed without the hyphen, for ``est-elle''.} the manifold $V$ one-sided or two-sided?

I propose to establish that it is two-sided if $q$ is odd and that it is on the contrary one-sided if $q$ is even.

Let us put
\[ \begin{aligned} &y_{1} = \cos \theta_{1}, &&z_{2} = \cos \theta'_{1},\\ &y_{2} = \sin \theta_{1} \cos \theta_{2}, &&z_{2} = \sin \theta'_{1} \cos \theta'_{2},\\ &y_{3} = \sin \theta_{1} \sin \theta_{2} \cos \theta_{3}, &&z_{3} = \sin \theta'_{1} \sin \theta'_{2} \cos \theta'_{3},\\ &\ldots\ldots, &&\ldots\ldots,\\ &y_{q-1} = \sin \theta_{1} \sin \theta_{2} \ldots \sin \theta_{q-2} \cos \theta_{q-1}, &&z_{q-1} = \sin \theta'_{1} \sin \theta'_{2} \ldots \sin \theta'_{q-2} \cos \theta'_{q-1},\\ &y_{q} = \sin \theta_{1} \sin \theta_{2} \ldots \sin \theta_{q-2} \sin \theta_{q-1}, &&z_{p} = \sin \theta'_{1} \sin \theta'_{2} \ldots \sin \theta'_{q-1} \sin \theta'_{q-1}. \end{aligned} \]\ednote{The first row of the right-hand column is printed $z_{2}$, apparently for $z_{1}$. The last row appears to be printed $z_{p}$, apparently for $z_{q}$, and its next-to-last factor is printed $\sin \theta'_{q-1}$, apparently for $\sin \theta'_{q-2}$ as in the left-hand column.}

The position of a point on $W$ is defined by the $2q - 2$ coordinates
\[ \theta_{1}, \qquad \theta_{2}, \qquad \ldots, \qquad \theta_{p-1}, \qquad \theta'_{1}, \qquad \theta'_{2}, \qquad \ldots, \qquad \theta'_{q-1}. \]\ednote{Printed $\theta_{p-1}$, apparently for $\theta_{q-1}$.}

On the other hand, the group $G$ consists (besides the identical substitution) of the single substitution that permutes $\theta_{i}$ with $\theta'_{i}$. For the manifold $V$ to be two-sided, it is therefore necessary and sufficient that the functional determinant
\[ \frac{\partial(\theta_{i}, \theta'_{i})}{\partial(\theta'_{i}, \theta_{i})}, \qquad (i = 1, 2, \ldots, q - 1) \]
be positive. Now it is equal to
\[ (-1)^{q-1}, \]
\origpage{94}
that is to say to $+1$ if $q$ is odd and to $-1$ if $q$ is even. Hence

$V$ is two-sided if $q$ is odd.

$V$ is one-sided if $q$ is even. Q.~E.~D.

Let us now concern ourselves with determining the Betti number
\[ P_{q-1}. \]

Let us first determine the Betti numbers for $W$.

We can construct on $W$ two manifolds of $q - 1$ dimensions in the following manner. One knows that every point of $W$ corresponds to a pair of points $QQ'$ of the hypersphere $S$. To the pairs $Q_{0}Q'$, where $Q_{0}$ is fixed and $Q'$ describes the whole hypersphere, there will then correspond a closed manifold $U_{1}$ of $q - 1$ dimensions forming part of $W$. Likewise to the pairs $QQ_{0}$, where $Q$ describes the whole hypersphere while $Q_{0}$ is fixed, there will correspond a closed manifold $U_{2}$ of $q - 1$ dimensions forming part of $W$.

These two manifolds will be linearly independent (from the point of view of homologies).

Let us consider, indeed, the integral of order $q - 1$
\[ J = \int \sin^{q-2} \theta_{1} \sin^{q-3} \theta_{2} \ldots \sin^{2} \theta_{q-3} \sin \theta_{q-1}\,d\vartheta_{1}\,d\vartheta_{2} \ldots d\theta_{q-1}, \]\ednote{The factor $\sin \theta_{q-1}$ is printed so, apparently for $\sin \theta_{q-2}$. The first two differentials are printed with the curly $\vartheta$, the last with $\theta$, as reproduced.}
and let $J'$ be the integral of order $q - 1$ obtained by replacing in $J$ the $\theta_{i}$ by the $\theta'_{i}$.

Let us next consider the integral
\[ J + \lambda J', \]
where $\lambda$ is an incommensurable number.

We know the $x$ as functions of the $\theta$ and the $\theta'$; conversely, we can calculate the $\theta$ and the $\theta'$ as functions of the $x$, and this in such a way that the integral $J + \lambda J'$ takes the form
\[ \int \sum X \partial, \]
\origpage{95}
$X$ being an entire function with respect to the $x$, and $\partial$ the product of $q - 1$ differentials $dx_{i}$.

I observe further that the integral $J + \lambda J'$ is zero when it is extended over an infinitely small closed manifold forming part of $W$.

Extended over $U_{2}$, it is equal to $\sigma$ (the surface of the hypersphere $S$): extended over $U_{1}$, it is equal to $\lambda\sigma$; and since there is no linear relation between $\sigma$ and $\lambda\sigma$; it follows that $U_{1}$ and $U_{2}$ are linearly independent.

The Betti number $P_{q-1}$ is therefore at least equal to 3.

To go further, let us consider the manifold $W'$ obtained by suppressing in $W$ all the points of the two manifolds $U_{1}$ and $U_{2}$; I propose to establish that $W'$ is simply connected.

Let us consider likewise the manifold $S'$, by suppressing in the hypersphere $S$ the point $Q_{0}$. To every pair of points $QQ'$ of the manifold $S'$ there will correspond a point of $W'$, and conversely.

But $S'$ is homeomorphic to the domain $D$,
\[ \eta^{2}_{1} + \eta^{2}_{2} + \ldots + \eta^{2}_{q-1} < 1, \]
forming part of space of $q - 1$ dimensions in which the coordinates are denoted by $\eta_{1}, \eta_{2}, \ldots, \eta_{q-1}$.

Thus the surface of an ordinary sphere, when one point is suppressed from it, becomes homeomorphic to the surface of a circle.

To each pair of points of the domain $D$ there will correspond a point of $W'$, from which it follows that $W'$ is homeomorphic to the domain $\Delta$ defined by the inequalities
\[ \begin{aligned} &\eta^{2}_{1} + \eta^{2}_{2} + \ldots + \eta^{2}_{q-1} < 1,\\ &\zeta^{2}_{1} + \zeta^{2}_{2} + \ldots + \zeta^{2}_{q-1} < 1, \end{aligned} \]
and situated in space of $2q - 2$ dimensions, in which the coordinates are denoted by $\eta_{i}$ and $\zeta_{i}$.

Now the domain $\Delta$ is simply connected because it is convex. Let, indeed, $v$ be a closed manifold of any number of dimensions forming part of $\Delta$; I say that one can, by deforming it in a continuous manner,
\origpage{96}
reduce it to a point without its leaving $\Delta$. Let us construct a manifold $v'$ by multiplying by a positive factor $k$ smaller than 1 the coordinates of all the points of $v$ (the manifold $v'$ will be homothetic to $v$, the centre of homothety being the origin and the ratio of homothety being equal to $k$). The manifold $v'$ will be entirely in the interior of $\Delta$ and, when $k$ decreases from 1 to 0, it will first coincide with $v$ and will end by reducing to a point. Q.~E.~D.

Hence $\Delta$, and consequently $W'$, are simply connected. Let us now determine the Betti numbers for $W$, and first the number $P_{h}$, where $h < q - 1$. Let us therefore consider a closed manifold of $h$ dimensions contained in $W$; this manifold, whether or not it meets $U_{1}$ and $U_{2}$, can always be regarded as homologous to a closed manifold of $h$ dimensions meeting neither $U_{1}$ nor $U_{2}$; the latter will be homologous to 0 in $W'$ (since $W'$ is simply connected) and \emph{a fortiori} in $W$.

Thus $P_{h}$ is equal to 1.

Since, on the other hand, the Betti numbers equally distinct\ednote{Printed ``distincts'', apparently for ``distants'', the wording of the theorem elsewhere in the memoir.} from the extremes are equal, $P_{h}$ will again be equal to 1 if $h$ is greater than $q - 1$.

It remains to determine $P_{q-1}$.

Let us consider a closed manifold $v$ of $q - 1$ dimensions forming part of $W$.

If it meets neither $U_{1}$ nor $U_{2}$, it will be homologous to 0; indeed, it forms part of $W'$, and consequently will be homologous to 0 with respect to $W'$ and \emph{a fortiori} with respect to $W$.

Let us suppose that it meets $U_{1}$ and $U_{2}$.

The points of $W$ that correspond to the pair $Q_{1}Q'$, where $Q_{1}$ is fixed and $Q'$ describes the whole hypersphere, will form a manifold $U'_{1}$, which will be homologous to $U_{1}$. Through every point of $W$ there passes one manifold $U'_{1}$ and only one.

The points of $W$ that correspond to the pair $QQ_{2}$, where $Q_{2}$ is fixed and $Q$ describes the whole hypersphere, will form a manifold $U'_{2}$, which will be homologous to $U_{2}$. Through every point of $W$ there passes one manifold $U'_{2}$ and only one.

Two manifolds $U'_{1}$ and $U'_{2}$ have one common point and only one (namely the point that corresponds to the pair $Q_{1}Q_{2}$). On the contrary, two manifolds $U'_{1}$ (or two manifolds $U'_{2}$) do not meet.
\origpage{97}

This being so, let us return to the manifold $v$ and suppose, to fix ideas, that it cuts $U_{1}$ in two points $M'$ and $M''$ and $U_{2}$ in one point $N'$. Through the point $M'$ I can pass a manifold $U'_{2}$, through the point $M''$ an analogous manifold $U''_{2}$; through the point $N'$ a manifold $U'_{1}$.

I can then, by deforming the manifold $v$ infinitely little, arrange that it still cuts $U_{1}$ and $U_{2}$ only at the points $M'$, $M''$ and $N'$, and that it admits around $M'$ a small common part $u'_{2}$ with $U'_{2}$; around $M''$ a small common part $u''_{2}$ with $U''_{2}$; around $N'$ a small common part $u'_{1}$ with $U'_{1}$.

Then the set of the manifolds
\[ U'_{1} - u'_{1}, \qquad U'_{2} - u'_{2}, \qquad U''_{2} - u''_{2}, \qquad v - u'_{1} - u'_{2} - u''_{2} \]
will form a closed manifold which will meet neither $U_{1}$ nor $U_{2}$, and which will consequently be homologous to 0; one will therefore have
\[ U'_{1} - u'_{1} + U'_{2} - u'_{2} + U''_{2} - u''_{2} \backsim v - u'_{1} - u'_{2} - u''_{2}, \]
whence
\[ v \backsim U'_{2} + U'_{1} + U''_{2} \backsim U_{1} + 2U_{2}. \]

It might also happen that, for example, the number that we called $S(M)$ above does not have the same value for the point $M'$ considered as a point of intersection of $v$ and $U_{1}$, and for the point $M'$ considered as a point of intersection of $U'_{2}$ and $U_{1}$. In this case one would have to replace $U'_{2}$ by the opposite manifold, and one will have
\[ v \backsim U'_{1} - U'_{2} + U''_{2} \backsim U_{1}. \]

In all cases, $v$, $U_{1}$ and $U_{2}$ will not be linearly independent, and one will have
\[ P_{q-1} = 3. \]

Let us determine finally the Betti numbers for $V$.

To the pairs $QQ'$ and $Q'Q$ there corresponds one and the same point of $V$; it
\origpage{98}
follows that to $U_{1}$ and to $U_{2}$ there corresponds in $V$ one and the same manifold, so that I can write
\[ U_{1} \backsim U_{2}. \]

On the other hand, the integral $J + J'$, extended over this manifold, is not zero; which shows that one does not have
\[ U_{1} \backsim 0. \]

The number $P_{q-1}$ is therefore at least equal to 2.

To every closed manifold $v$ forming part of $V$ there will correspond a manifold $w$ forming part of $W$; but two cases may present themselves; we know that to each point of $V$ there correspond two points of $W$; and I may say, to shorten the language, that these two points are \emph{symmetric}, since one passes from the one to the other by permuting the $y_{i}$ with the $z_{i}$.

We shall construct $w$ by taking for each point of $v$ \emph{one} of the two points that correspond to it in $w$\ednote{The print has lower-case $w$ here; apparently for $W$, the variety in which each point of $V$ has two corresponding points.}; it may then happen either that $w$ is closed, or that $w$ is not closed but its boundary consists of parts that are symmetric in pairs.

Let us consider first the case where $w$ is closed.

If the number of dimensions is different from $q - 1$, $w$ (and consequently $v$) will be able, by continuous deformation, to reduce to a point, and one will have
\[ v \backsim 0. \]

If the number of dimensions is equal to $q - 1$, one will have, with respect to $W$,
\[ w \backsim mU_{1} + nU_{2}, \]
$m$ and $n$ being integers; but with respect to $V$, $U_{1}$ is homologous to $U_{2}$.

One will therefore have, with respect to $V$,
\[ v \backsim (m + n)U_{1}. \]

Let us now consider the case where $w$ is not closed. We shall admit
\origpage{99}
that the number of dimensions is less than or equal to $q - 1$; the case where this number is greater than $q - 1$ is easily reduced to it, since the Betti numbers are equal in pairs. The boundary $f$ of $w$ will then have fewer than $q - 1$ dimensions.

Let $H$ be the manifold of $q - 1$ dimensions forming part of $W$ which consists of the points that are their own symmetric, that is to say of the points
\[ y_{i} = z_{i}. \]

The boundary $f$ can be deformed in a continuous manner, without its points ceasing to be symmetric in pairs, and in such a way that at the end of the deformation the deformed boundary forms part of $H$.

Hence $w$ can be deformed in a continuous manner, without ceasing to correspond to a closed manifold $v$, and in such a way that at the end of the deformation the deformed manifold $w$ is closed.

Hence $v$ is always homologous to a manifold $v$ corresponding to a closed manifold $w$.

Hence one has, according to the case,
\[ v \backsim 0, \qquad v \backsim (m + n)U_{1}. \]

Hence all the Betti numbers of $V$ are equal to 1, except $P_{q-1}$, which is equal to 2.

From this, two consequences:

1° Let $q$ be odd; $V$ will be two-sided; from which it follows that, for a two-sided manifold of $4k$ dimensions, the number $P_{2k}$ is not necessarily odd.

2° If $q$ is even, $V$ will be one-sided; from which it follows that, for a one-sided manifold of $4k + 2$ dimensions, the number $P_{2k+1}$ is not necessarily odd, as would be the case for a two-sided manifold.

I had announced these results at the end of \no 9.
\origpage{100}

\subsection*{\S~16. --- Euler's theorem.}

One knows Euler's theorem, according to which, if $S$, $A$ and $F$ are the number of the vertices, the edges and the faces of a convex polyhedron, one must have
\[ S - A + F = 2. \]

This theorem was generalized by Admiral de Jonquières to the case of non-convex polyhedra. If a polyhedron forms a closed manifold of two dimensions whose Betti number is $P_{1}$, one will have
\[ S - A + F = 3 - P_{1}. \]

The fact that the faces are plane evidently has no importance; the theorem applies equally to curvilinear polyhedra; it applies further to the subdivision of any closed surface into simply connected regions; these regions then correspond to the faces of the polyhedron; the lines that serve as boundary to two of these regions correspond to the edges, and the extremities of these lines to the vertices.

I propose to generalize these results to any space.

Let, then, $V$ be a closed manifold of $p$ dimensions. Let us subdivide it into a certain number of manifolds $v_{p}$ of $p$ dimensions; these manifolds $v_{p}$ will not be closed, and their boundaries will be formed by a certain number of manifolds $v_{p-1}$ of $p - 1$ dimensions; the boundaries of the $v_{p-1}$ will be formed in their turn by a certain number of manifolds $v_{p-2}$ of $p - 2$ dimensions, and so on; I shall arrive next at a certain number of manifolds $v_{1}$ of one dimension, which will have as boundaries a certain number of isolated points, or manifolds of 0 dimensions, which I shall call $v_{0}$.

The manifold $V$ may have any Betti numbers, but I suppose expressly that the manifolds $v_{p}, v_{p-1}, \ldots, v_{1}$ are simply connected.

I shall call $\alpha_{p}, \alpha_{p-1}, \ldots, \alpha_{1}$ and $\alpha_{0}$ the number of the $v_{p}$, of the $v_{p-1}$, \ldots, of the $v_{1}$ and of the $v_{0}$.
\origpage{101}

The figure formed by all these manifolds may be called a polyhedron; for the analogy with ordinary polyhedra is evident. An ordinary polyhedron is, indeed, a closed manifold $V$ of 2 dimensions, which is subdivided into a certain number of manifolds $v_{2}$, which are the faces. The faces have as boundaries a certain number of manifolds $v_{1}$, which are the edges, and which admit in their turn as boundaries a certain number of manifolds $v_{0}$ called \emph{vertices}.

I propose to calculate the number
\[ N = \alpha_{p} - \alpha_{p-1} + \alpha_{p-2} - \ldots \mp \alpha_{1} \pm \alpha_{0}. \]

I shall introduce here some new denominations, perhaps ill justified, but convenient.

If two polyhedra are obtained by subdividing a manifold $V$ in two different ways, I shall say that they are \emph{congruent}.

Now let there be the polyhedron $P$ formed by the manifold $V$, by the regions $v_{p}$ and by their successive boundaries $v_{p-1}, \ldots, v_{1}, v_{0}$.

Let us subdivide the $v_{p}$ into smaller regions $v'_{p}$; the complete boundaries of the $v'_{p}$ will consist of a certain number of new regions $v''_{p-1}$ and, in addition, of regions $v'_{p-1}$ obtained by subdividing the $v_{p-1}$; the complete boundaries of the $v'_{p-1}$ and of the $v''_{p-1}$ will consist of a certain number of new regions $v''_{p-2}$ and, in addition, of regions $v'_{p-2}$ obtained by subdividing the $v_{p-2}$, and so on; we shall arrive finally at the $v'_{1}$ and the $v''_{1}$, whose complete boundaries will consist of a certain number of new points $v''_{0}$ and, in addition, of the points $v_{0}$.

Let $P'$ be the polyhedron formed by the set of the regions $v'_{p}$, $v'_{p-1}$, $v''_{p-1}$, $v'_{p-2}$, $v''_{p-2}$, \ldots, $v'_{1}$, $v''_{1}$, $v_{0}$, $v''_{0}$.

I shall then say that the polyhedron $P'$ is \emph{derived} from the polyhedron $P$.

I shall clarify this definition by an example borrowed from ordinary Geometry. Let us consider a regular tetrahedron $T$. In each of the faces I join each vertex to the midpoint of the opposite side. Each face will thus be decomposed into six triangles; twenty-four triangles in all. The polyhedron with twenty-four faces thus obtained will be \emph{derived} from $T$.
\origpage{102}

Now let $P$ and $P'$ be two congruent polyhedra, that is to say obtained by two different modes of decomposition of one and the same manifold $V$; there will always exist a polyhedron $P''$ which will be derived at once from $P$ and from $P'$, and which one will obtain by combining the two modes of decomposition; in such a way that, if we call $v_{p}$, $v'_{p}$ and $v''_{p}$ the subdivisions of $V$ in the three modes of decomposition that correspond to the three polyhedra $P$, $P'$ and $P''$, the necessary and sufficient condition for two points to belong to the same region $v''_{p}$ is that they belong at once to the same region $v_{p}$ and to the same region $v'_{p}$.

I propose to establish that the number $N$ is the same for two congruent polyhedra and, since we have just seen that two congruent polyhedra always have a common derivative, it will suffice for us to show that the number $N$ is the same for a polyhedron and for all its derivatives.

If, in the polyhedron $P$, we consider one of the regions $v_{p-1}$, it will always belong to two regions $v_{p}$, and to two only, which it will separate from one another. On the contrary, a region $v_{p-2}$ will belong, in general, to more than two regions $v_{p}$ and to more than two regions $v_{p-1}$. Thus, in ordinary polyhedra, an edge always separates two faces from one another, while a vertex belongs in general to more than two faces and to more than two edges.

We do not, however, exclude the case where a region $v_{p-2}$ would belong to two regions $v_{p-1}$ only. Thus, for an ordinary polyhedron, we should not exclude the case where the midpoint of an edge would be regarded as a vertex, and where this edge would consequently be regarded as formed of two juxtaposed edges.

The regions $v_{p-2}$ which would thus belong only to two regions $v_{p-1}$ will be called \emph{singular}. Let, then, $v_{p-2}$ be a singular region belonging to two regions $v_{p-1}$ which I call $v'_{p-2}$\ednote{Printed $v'_{p-2}$; apparently for $v'_{p-1}$, since the two regions named here are $v_{p-1}$.} and $v''_{p-1}$. It is clear that $v'_{p-1}$ will separate from one another \emph{the same two regions} $v_{p}$ that are separated by $v''_{p-1}$; so that $v_{p-2}$ will also belong only to two regions $v_{p}$.

Likewise I shall say that the manifold $v_{q}$ is \emph{singular} if it belongs only to two manifolds $v_{q+1}$; and in that case it will likewise belong only to two manifolds $v_{q+2}$, $v_{q+3}$, \ldots, $v_{p}$.
\origpage{103}
Now let there be a manifold $v_{h}$; to this manifold there will belong a certain number of manifolds $v_{h-1}$; if one of them is singular, I shall say that the manifold $v_{h}$ is \emph{irregular}; it will be regular in the contrary case.

Let us consider, then, the polyhedron $P$ with the regions $v_{p}$, $v_{p-1}$, \ldots\ and the derived polyhedron $P'$ with the regions $v'_{p}$, $v'_{p-1}$, \ldots. Let us seek to go back from the polyhedron $P'$ to the polyhedron $P$. Let there be two regions $v'_{p}$, which I shall call $\alpha$ and $\beta$; I suppose that they are separated from one another by a region $v'_{p-1}$, which I shall call $\gamma$; that they are consequently contiguous, and that they form part of one and the same region $v_{p}$. (Since $\alpha$ and $\beta$ form part of one and the same region $v_{p}$, the region $\gamma$ is not a subdivision of one of the regions $v_{p-1}$ that separate the regions $v_{p}$ from one another; this region $\gamma$ is therefore one of those I had called $v''_{p-1}$ in the definition of derived polyhedra; but here I no longer make this distinction, and I denote the manifolds I then called $v''_{p-1}$, as well as those I then called $v'_{p-1}$, by the same notation $v'_{p-1}$.)

This being so, let us suppress the region $\gamma$, which serves as boundary to $\alpha$ and to $\beta$, and \emph{annex} the region $\alpha$ to the region $\beta$. We shall thus have suppressed one region $v'_{p}$ and one region $v'_{p-1}$. On the other hand, we shall not have suppressed any region $v'_{p-2}$ if the region $\gamma$ is regular; if, indeed, none of the regions $v'_{p-2}$ is singular, each of them will belong to at least three regions $v'_{p-1}$, and after the suppression of $\gamma$ it will still belong to at least two regions $v'_{p-1}$. Likewise every region $v'_{q}$ (where $q < p - 2$) forming part of $\gamma$ will belong to at least three regions $v'_{p-1}$, and after the suppression of $\gamma$ it will still belong to at least two regions $v'_{p-1}$. The suppression of $\gamma$ therefore suppresses none of the regions $v'_{q}$; \emph{it therefore does not change the value of the number $N$}.

If, on the contrary, the region $\gamma$ is irregular, we no longer have the right to suppress it, for there will then exist a region $v'_{p-2}$ that belongs only to $\gamma$ and to one other region $v'_{p-1}$; after the suppression of $\gamma$, it would belong to only a single region $v'_{p-1}$, which is inadmissible.

What must be done then? The region $\gamma$ separates two regions $v_{p}$ which I have called $\alpha$ and $\beta$; but it does not by itself alone constitute the boundary between $\alpha$ and $\beta$; indeed, since $\gamma$ is irregular, there will exist a singular region $v'_{p-2}$
\origpage{104}
which I shall call $\delta$, and which will belong to $\gamma$ and to another region $v'_{p-1}$, which I shall call $\gamma'$. This region $\gamma'$, by what we have seen above, separates the same regions as $\gamma$, that is to say $\alpha$ and $\beta$.

If the region $\delta$ is regular, we shall be able to suppress it and annex $\gamma$ to $\gamma'$. The region $\gamma + \gamma'$ will then separate $\alpha$ from $\beta$. We shall thus have diminished $\alpha_{p-1}$ and $\beta_{p-2}$\ednote{Printed $\beta_{p-2}$; apparently for $\alpha_{p-2}$, the number of the $v_{p-2}$.} by one unit, while the other numbers $\alpha_{i}$ will not have changed. $N$ will therefore not have changed either.

If $\delta$ is irregular, there will be a singular region $v'_{p-3}$, which I shall call $\varepsilon$ and which will separate it from another region $\delta'$; we shall suppress $\varepsilon$ and annex $\delta$ to $\delta'$, and so on.

We shall thus be able to suppress a region $v'_{q}$ that separates two regions $v'_{q+1}$ from one another, and annex these two regions $v'_{q+1}$ to one another, but on two conditions:

1° If $q$ is smaller than $p - 1$, the region $v'_{q}$ must be singular;

2° And in all cases it must be regular.

This being so, here is the order in which we shall carry out the operations:

I wish to go back from the polyhedron $P'$ to the polyhedron $P$. I may without inconvenience suppose that the polyhedron $P$ admits no singular region, but the polyhedron $P'$ and the intermediate polyhedra will admit some.

By a series of successive suppressions and annexations we shall go back from $P'$ to $P$, passing through a series of intermediate polyhedra which I shall call
\[ P_{0} = P', \qquad P_{1}, \qquad P_{2}, \qquad \ldots, \qquad P_{m-1}, \qquad P_{m}. \]

How shall we pass from the polyhedron $P_{i}$ to the polyhedron $P_{i+1}$?

If in $P_{i}$ there are singular $v'_{0}$, I shall suppress one of them. If there are none, all the $v'_{1}$ will be regular; if there are singular $v'_{1}$, I shall suppress one of them.

If there is no singular $v'_{1}$, all the $v'_{2}$ will be regular; if there are singular $v'_{2}$, I shall suppress one of them.

And so on.

If finally there is no singular $v'_{p-2}$, all the $v'_{p-1}$ will be regular, and one will have the right to suppress any one of them; if one of the regions $v_{p}$ is subdivided into several regions $v'_{p}$, I shall choose two of these
\origpage{105}
regions $v'_{p}$ that are contiguous and separated from one another by a region $v'_{p-1}$ which is their common boundary; I shall annex them to one another by suppressing this common boundary.

None of these operations can alter the number $N$.

One will be stopped only when there is no longer any singular region and when, moreover, none of the regions $v_{p}$ is any longer subdivided into several regions $v'_{p}$. But then one will have arrived at the polyhedron $P$.

None of these operations can alter the number $N$.

This number therefore has the same value for $P$ and for $P'$. Q.~E.~D.

This proof might give rise to certain objections, for one might ask whether, in this series of operations, all the regions will remain simply connected; but, before modifying our proof so as to put myself beyond these objections, I wish to determine what the value of the number $N$ must be for a simply connected polyhedron.

If our theorem is true, the number $N$ must have the same value for two polyhedra obtained by subdividing two homeomorphic manifolds; it therefore has the same value for any two simply connected polyhedra.

It will therefore suffice for us to make this determination for an arbitrarily chosen simply connected polyhedron.

I shall choose the generalized tetrahedron.

I so call the polyhedron formed by the complete boundary of the domain
\[ \begin{aligned} &x_{1} > 0, &&x_{2} > 0, &&\ldots, &&x_{p} > 0,\\ &x_{p+1} > 0, &&x_{1} + x_{2} + \ldots + x_{p} + x_{p+1} < 0. \end{aligned} \]\ednote{Printed $< 0$; apparently for $< 1$, since with every $x_{i} > 0$ the sum cannot be negative.}

One then has
\[ \begin{aligned} &\alpha_{p} = p + 2, &&\alpha_{p-1} = \frac{(p + 1)(p + 2)}{2}, &&\ldots,\\ &\alpha_{q} = \frac{(p + 2)!}{(q + 2)!(p - q)!}, &&\ldots, &&\alpha_{1} = \frac{(p + 1)(p + 2)}{2}, &&\alpha_{0} = p + 2, \end{aligned} \]\ednote{The denominator is printed $(q + 2)!(p - q)!$; the values $\alpha_{p} = p + 2$ and $\alpha_{1}$ printed alongside apparently call for $(q + 1)!(p - q + 1)!$.}
that is to say that the numbers $\alpha_{q}$ are equal to the binomial coefficients.
\origpage{106}
One will therefore have
\[ (1 - 1)^{p+2} = 1 - \alpha_{p} + \alpha_{p-1} - \ldots \pm \alpha_{1} \mp \alpha_{0} \pm 1 = 1 - N \pm 1. \]

One must take before the last term the sign $+$ if $p$ is even, and the sign $-$ if $p$ is odd.

One will therefore have $N = 2$ if $p$ is even, $N = 0$ if $p$ is odd.

I should have arrived at the same result by choosing the generalized cube.

I so call the polyhedron formed by the complete boundary of the domain
\[ -1 < x_{i} < 1, \qquad (i = 1, 2, \ldots, p + 1). \]

One then has
\[ \begin{aligned} &\alpha_{p} = 2(p + 1), &&\alpha_{p-1} = 2^{2}\frac{p(p + 1)}{2}, &&\ldots, &&\alpha_{q} = 2^{p-p+1}\frac{(p + 1)!}{(q + 1)!(p - q)!}, &&\ldots,\\ &\alpha_{1} = 2^{p}(p + 1), &&\alpha_{0} = 2^{p+1}. \end{aligned} \]\ednote{The exponent is printed $p-p+1$, apparently for $p-q+1$; the factor $\frac{(p + 1)!}{(q + 1)!(p - q)!}$ as printed does not reproduce the $\alpha_{p}$ and $\alpha_{p-1}$ printed alongside, which apparently call for $\frac{(p + 1)!}{q!(p - q + 1)!}$.}
whence
\[ (1 - 2)^{p+1} = 1 - \alpha_{p} + \alpha_{p-1} - \ldots \pm \alpha_{1} \mp \alpha_{0} = 1 - N, \]
whence
\[ N = 1 - (-1)^{p+1}, \]
that is to say

$N = 2$ if $p$ is even,

$N = 0$ if $p$ is odd.

Thus, for a simply connected polyhedron, the number $N$ is equal to 2 if $p$ is even and to 0 if $p$ is odd.

This being so, to establish our theorem in a complete and rigorous way, I am going to suppose that it is true for manifolds of fewer than $p$ dimensions.

Let us consider, then, our polyhedron $P$ and a region $v_{q}$ of $q$ dimensions belonging to this polyhedron. This region $v_{q}$ will form part of a certain number of regions $v_{q+1}$, of a certain number of regions $v_{q+2}$, \ldots, of a certain number of regions $v_{p}$. The set of all these regions will form what I shall call the \emph{aster} of $v_{q}$.
\origpage{107}

I shall denote by $\gamma_{h}$ the number of the regions $v_{h}$ $(h > q)$ that thus form part of the aster of $v_{q}$.

Let $(x^{0}_{1}, x^{0}_{2}, \ldots, x^{0}_{n})$ be a point of $v_{q}$; let us consider the hypersphere $S$ which has as equation
\[ (x_{1} - x^{0}_{1})^{2} + (x_{2} - x^{0}_{2})^{2} + \ldots + (x_{n} - x^{0}_{n})^{2} = \varepsilon^{2}, \]
$\varepsilon$ being very small.

Let $\Pi$ be the plane manifold defined by the $q$ equations
\[ \begin{aligned} &A^{1}_{i}(x_{1} - x^{0}_{1}) + A^{2}_{i}(x_{2} - x^{0}_{2}) + \ldots + A^{n}_{i}(x_{n} - x^{0}_{n}) = 0,\\ &(i = 1, 2, \ldots, q), \end{aligned} \]
where the $A^{k}_{i}$ are any constants.

The intersection of $S$, $\Pi$ and $V$ will be a manifold of $p - q - 1$ dimensions, which I shall call $W$, and which will be simply connected. The intersections of $S$ and $\Pi$ with the various regions that form the aster of $v_{q}$ will together form a polyhedron arising from the subdivision of $W$, which will be simply connected.

For this polyhedron the number $\alpha_{h}$ will be equal to $\gamma_{h+q+1}$; since it has fewer than $p$ dimensions, our theorem will be applicable to it, so that we shall be able to write
\[ \gamma_{p} - \gamma_{p-1} + \ldots \pm \gamma_{q+2} \mp \gamma_{q+1} = 2 \text{ or } 0, \tag{A} \]
according as $p - q$ is even or odd\ednote{Printed ``pair ou impair''. Since $W$ has $p - q - 1$ dimensions, the value 2 apparently goes with $p - q$ odd; compare p. 112, ``impair ou pair''.}.

Let us now define a polyhedron $Q$, which will be formed by an operation that may be called a \emph{gridding}.

Let $V$ be a manifold of $p$ dimensions situated in space of $n$ dimensions. Let us construct an infinity of plane manifolds defined by the equations
\[ \begin{aligned} &x_{i} = a_{k.i},\\ &(i = 1, 2, \ldots, n; \qquad K = -\infty, \ldots, -1, 0, +1, +2, \ldots, +\infty). \end{aligned} \tag{B} \]\ednote{Printed capital $K$, apparently for the index $k$ of $a_{k.i}$. The first $\infty$ is printed in an open, damaged-looking form.}
\origpage{108}

These plane manifolds will decompose space into an infinity of domains $D_{n}$ comparable to rectangular parallelepipeds. The boundaries of the $D_{n}$ will be formed by a certain number of domains $D_{n-1}$ of $n - 1$ dimensions forming part of the various plane manifolds $x_{i} = a_{k.i}$ and likewise comparable to rectangular parallelepipeds. The boundaries of the $D_{n-1}$ will be formed by a certain number of domains $D_{n-2}$ comparable to rectangular parallelepipeds in space of $n - 2$ dimensions, and so on.

Then the polyhedron $Q$ will be defined in the following way: the regions $v_{p}$ will be formed by the intersections of $V$ with the domains $D_{n}$, the regions $v_{p-1}$ by the intersections of $V$ with the domains $D_{n-1}$, and so on; finally, the regions $v_{0}$ by the intersections of $V$ with the domains $D_{n-1}$\ednote{Printed $D_{n-1}$; the construction apparently calls for $D_{n-p}$, since $V$ has $p$ dimensions.}.

It follows from this definition that the polyhedron $Q$ admits no singular region.

I shall consider in addition any polyhedron $P$ congruent to $Q$, and a polyhedron $P'$ derived at once from $P$ and from $Q$.

I am going to go back, on the one hand, from $P'$ to $P$ and, on the other, from $P'$ to $Q$, and I shall establish that in these two operations the number $N$ has not changed.

Let us go back first from $P'$ to $P$.

Let $x_{i} = a$ be one of the plane manifolds defined by the equation (B). Let us classify the regions $v'_{q}$, of any number of dimensions, that compose the polyhedron $P'$ into four sorts.

Those of the first sort are those that form part of the manifold
\[ x_{i} = a. \]

Those of the second sort are those that admit points such that
\[ x_{i} = a + \varepsilon, \]
$\varepsilon$ being positive and very small.

Those of the third sort are those that admit points such that
\[ x_{i} = a - \varepsilon. \]

All the others are of the fourth sort.
\origpage{109}

Let $\delta_{q}$, $\delta'_{q}$, $\delta''_{q}$ be the number of the manifolds of $q$ dimensions that are respectively of the first, of the second and of the third sort.

Every manifold of the second sort will be contiguous to a manifold of the third sort, and their common boundary will be a manifold of the first sort which will have one dimension fewer. The manifolds of the first three sorts therefore correspond one to one, and one has
\[ \delta'_{q} = \delta''_{q} = \delta_{q-1}. \]

It is clear moreover that $\delta'_{0} = \delta''_{0} = \delta_{p} = 0$.

If, in the set of plane manifolds (B) that constitute the gridding and have given rise to the polyhedra $Q$ and $P'$, one had suppressed the manifold $x_{i} = a$, one would have obtained two polyhedra $Q_{1}$ and $P'_{1}$ simpler than the first. Let us compare $P'_{1}$ with $P'$.

When one suppresses the plane manifold $x_{i} = a$, one suppresses the manifolds of the first sort and one \emph{annexes} each manifold of the third sort to the corresponding manifold of the second sort. Hence, in passing from $P'$ to $P'_{1}$, the number $\alpha_{q}$ diminishes by
\[ \delta''_{q} + \delta_{q} = \delta_{q} + \delta_{q-1}. \]

In particular, the extreme numbers $\alpha_{p}$ and $\alpha_{0}$ diminish by $\delta_{p-1}$ and $\delta_{0}$. It follows from this that the number $N$ diminishes by
\[ \delta_{p-1} - (\delta_{p-1} + \delta_{p-2}) + (\delta_{p-2} + \delta_{p-3}) - \ldots \pm (\delta_{1} + \delta_{0}) \mp \delta_{0} = 0. \]

Hence $N$ does not change. Thus, by suppressing the manifold $x_{i} = a$, one does not change $N$; but by suppressing in this way all the plane manifolds defined by (B), one will go back to the polyhedron $P$. The number $N$ is therefore the same for $P'$ and for $P$.

Let us now go back from $P'$ to $Q$.

Let $w_{p}$, $w_{p-1}$, \ldots, $w_{1}$, $w_{0}$ be the manifolds whose set constitutes the polyhedron $Q$; let likewise $v'_{p}$, $v'_{p-1}$, \ldots, $v'_{1}$, $v'_{0}$ be the manifolds whose set constitutes the polyhedron $P'$.

We shall distribute the manifolds $v'_{q}$ into $p + 1$ classes.
\origpage{110}

Those of the first class will be those that form part of one of the regions $w_{p}$ without forming part of one of the regions $w_{p-1}$. The polyhedron $P'$ being derived from $Q$, we shall place in this first class all the manifolds $v'_{p}$ (which are all subdivisions of the $w_{p}$); those of the manifolds $v'_{p-1}$ that separate from one another two manifolds $v'_{p}$ forming part of one and the same region $w_{p}$; and their intersections.

Those of the second class will be those that form part of one of the regions $w_{p-1}$; without forming part of one of the regions $w_{p-2}$.

Those of the third class will be those that form part of one of the regions $w_{p-2}$, without forming part of one of the regions $w_{p-3}$, etc.

Those of the $p^{\text{th}}$ class will be those that form part of one of the regions $w_{1}$, without being one of the points $w_{0}$.

Finally, the $(p + 1)^{\text{th}}$ class will comprise the points $w_{0}$.

To go back from $P'$ to $Q$, here is what I am going to do: I shall begin by suppressing all the manifolds of the first class that have fewer than $p$ dimensions, which has the effect of uniting into a single one, by annexation, all the regions $v'_{p}$ that are the subdivisions of one and the same region $w_{p}$.

I say that this operation does not alter the number $N$.

Indeed, I may suppose that the meshes of the gridding that gave rise\ednote{Printed ``naisasnce''; apparently a misprint for ``naissance''.} to $Q$ are close enough that, in the interior of one of these meshes $D_{p}$, that is to say in the interior of one of the regions $w_{p}$, one cannot find points belonging to two different manifolds $v_{p-1}$, except in the case where one finds there points belonging to the intersection of these two manifolds. (I always denote by $v_{q}$ the manifolds whose set forms $P$.) More generally, I may suppose that one cannot find in a region $w_{p}$ points belonging to several manifolds $v_{q} (q < p)$ unless one can find there points belonging to the intersection of these various manifolds.

In a region $w_{p}$ we shall therefore be able to have points of a region $v_{q}$ and of all the regions $v_{h} (h > q)$ that form part of its aster.

But we shall not be able to have points of two regions $v_{q}$, or points of $v_{q}$ and of a region $v_{h} (h > q)$ not forming part of the aster of $v_{q}$, without having, in addition, points of a region $v_{q-1}$.
\origpage{111}

If, then, I suppose that $q$ is the smallest value that can be given to the number of dimensions of a region $v_{q}$ for this region to have points in the interior of $w_{p}$; if I suppose this, I say, we shall have in $w_{p}$ only points of a \emph{single} region $v_{q}$ and of the regions of its aster.

Let us then consider a region $w_{p}$ containing points of $v_{q}$ and of the regions of the aster of $v_{q}$. Let
\[ \gamma_{q+1}, \qquad \gamma_{q+2}, \qquad \ldots, \qquad \gamma_{p} \]
be the numbers that we defined above in defining the aster.

It follows that in the interior of $w_{p}$ we shall have
\[ \begin{matrix} 1 & \text{region} & v'_{q} & \text{of the first class,} \\[1ex] \gamma_{q+1} & \text{»} & v'_{q+1} & \text{»} \\[1ex] \gamma_{q+2} & \text{»} & v'_{q+2} & \text{»} \\[1ex] \ldots\ldots & \ldots\ldots & \ldots\ldots & \ldots\ldots \\[1ex] \gamma_{p} & \text{»} & v'_{p} & \text{»} \end{matrix} \]

By suppressing the regions of the first class of fewer than $p$ dimensions and annexing to one another the $\gamma_{p}$ regions $v'_{p}$ whose set constitutes $w_{p}$, we diminish $\alpha_{p}$ by $\gamma_{p} - 1$, $\alpha_{p-1}$ by $\gamma_{p-1}$, \ldots, $\alpha_{q+1}$ by $\gamma_{q+1}$, $\alpha_{q}$ by 1. Hence, by virtue of the equation (A), the number $N$ does not change.

This done, we shall suppress all the manifolds of the second class of fewer than $p = 1$\ednote{Printed $p = 1$ with two clean strokes; apparently for $p - 1$, as the parallel $p - 2$ below and the $v'_{p-1}$ that follows indicate.} dimensions, annexing to one another all the manifolds $v'_{p-1}$ of the second class that form part of one and the same region $w_{p-1}$. Next, we shall suppress all the manifolds of the third class of fewer than $p - 2$ dimensions, annexing to one another all the manifolds $v'_{p-2}$ that form part of one and the same region $w_{p-2}$; and so on.

We shall thus finally arrive at the polyhedron $Q$.

One would show in the same way as above that none of these operations alters the number $N$.

The number $N$ is therefore the same for $P'$ and $Q$; it is therefore also the same
\origpage{112}
for $P$ and $Q$ and, consequently, for any two congruent polyhedra. Q.~E.~D.

\subsection*{\S~17. --- Case where $p$ is odd.}

I am going to define, with regard to any polyhedron $P$, new remarkable numbers, which I shall call $\beta_{\lambda.\mu}$.

Let first $\lambda > \mu$; I shall consider all the manifolds $v_{\lambda}$; for each of them I shall consider the number of the manifolds $v_{\mu}$ that form part of it; I shall take the sum of all these numbers relative to the various manifolds $v_{\lambda}$, and it is this sum that I shall call $\beta_{\lambda.\mu}$.

Since the manifolds $v_{\lambda}$ are all by hypothesis simply connected, one will have
\[ \beta_{\lambda.\lambda-1} - \beta_{\lambda.\lambda-2} + \ldots \pm \beta_{\lambda.1} \mp \beta_{\lambda.0} = 2\alpha_{\lambda} \text{ or } 0, \]
according as $\lambda$ is odd or even.
Let now $\lambda < \mu$; I shall consider all the manifolds $v_{\lambda}$; for each of them I shall consider the number of the manifolds $v_{\mu}$ of which it forms part (that is to say the number $\gamma_{\mu}$ relative to the aster of $v_{\lambda}$); I shall take the sum of all these numbers relative to the various manifolds $v_{\lambda}$, and it is this sum that I shall call $\beta_{\lambda.\mu}$.

By virtue of the equation (A) of the preceding number, one will have
\[ \beta_{\lambda.p} - \beta_{\lambda.p-1} + \ldots \pm \beta_{\lambda.\lambda+2} \mp \beta_{\lambda.\lambda+1} = 2\alpha_{\lambda} \text{ or } 0, \]
according as $p - \lambda$ is odd or even.

It follows from this definition that
\[ \beta_{\lambda.\mu} = \beta_{\mu.\lambda}. \]

This being so, let us form the following Table
\[ \begin{matrix} +\beta_{p.p-1}, & -\beta_{p.p-2}, & +\beta_{p.p-3}, & \ldots, & \pm\beta_{p.1}, & \mp\beta_{p.0}, \\[1ex] & +\beta_{p-1.p-2}, & -\beta_{p-1.p-3}, & \ldots, & \mp\beta_{p-1.1}, & \pm\beta_{p-1.0}, \\[1ex] & \ldots\ldots, & \ldots\ldots, & \ldots, & \ldots\ldots, & \ldots\ldots, \\[1ex] & & \ldots\ldots, & \ldots, & \ldots\ldots, & \ldots\ldots, \\[1ex] & & & \ldots, & \ldots\ldots, & \ldots\ldots, \\[1ex] & & & & +\beta_{2.1}, & -\beta_{2.0}, \\[1ex] & & & & & +\beta_{1.0}. \end{matrix} \]
\origpage{113}

One sees that in each row and in each column each term is affected alternately with the sign $+$ and the sign $-$, in such a way that $\beta_{\lambda.\mu}$ is affected with the sign $+$ if $\lambda - \mu$ is odd, and with the sign $-$ in the contrary case.

Let us take the sum of the terms of this Table; this sum can be carried out in two ways: by summing first the rows and by summing first the columns.

The sum of the terms of the various rows of the Table is, beginning from the top:
\[ 2\alpha_{p}, \qquad 0, \qquad 2\alpha_{p-2}, \qquad 0, \qquad \ldots, \qquad 2\alpha_{3}, \qquad 0, \qquad 2\alpha_{1} \]
if $p$ is odd, and
\[ 0, \qquad 2\alpha_{p-1}, \qquad 0, \qquad 2\alpha_{p-3}, \qquad \ldots, \qquad 2\alpha_{3}, \qquad 0, \qquad 2\alpha_{1} \]
if $p$ is even. The sum of the terms of the Table is therefore:
\[ 2\alpha_{1} + 2\alpha_{3} + \ldots + 2\alpha_{p} \]
if $p$ is odd, and
\[ 2\alpha_{1} + 2\alpha_{3} + \ldots + 2\alpha_{p-1} \]
if $p$ is even.

The sum of the terms of the various columns of the Table is, beginning from the left:
\[ 2\alpha_{p-1}, \qquad 0, \qquad 2\alpha_{p-3}, \qquad 0, \qquad \ldots, \qquad 2\alpha_{2}, \qquad 0, \qquad 2\alpha_{0} \]
if $p$ is odd, and
\[ 2\alpha_{p-1}, \qquad 0, \qquad 2\alpha_{p-3}, \qquad 0, \qquad \ldots, \qquad 0, \qquad 2\alpha_{1}, \qquad 0 \]
if $p$ is even.

The sum of the terms of the Table is therefore:
\[ 2\alpha_{0} + 2\alpha_{2} + \ldots + 2\alpha_{p-1} \]
if $p$ is odd, and
\[ 2\alpha_{1} + 2\alpha_{3} + \ldots + 2\alpha_{p-1} \]
if $p$ is even.
\origpage{114}

By equating the two expressions of this sum, one will obtain an identity if $p$ is even, and the equation
\[ N = 0 \]
if $p$ is odd.

Whence this consequence:

The number $N$ is zero and independent of the Betti numbers if $p$ is odd; it depends, on the contrary, on the Betti numbers if $p$ is even.

\subsection*{\S~18. --- Second proof.}

The second proof will teach us how it depends on them.

To make it well understood, I am going first to set it out for ordinary polyhedra with $\alpha_{0}$ vertices, $\alpha_{1}$ edges and $\alpha_{2}$ faces.

To each of the $\alpha_{0}$ vertices I make correspond an arbitrary number; to each of the $\alpha_{1}$ edges I make correspond a number $\delta$ equal to the difference of the numbers corresponding to its two vertices.

We shall thus have $\alpha_{1}$ differences $\delta$; but they cannot all be chosen arbitrarily; indeed, they will be determined when one knows the $\alpha_{0}$ numbers assigned to the various vertices, and even when one knows the excesses of $\alpha_{0} - 1$ of these numbers over the last. There will therefore be in all $\alpha_{0} - 1$ differences $\delta$ that can be chosen arbitrarily.

There must therefore be, between the differences $\delta$,
\[ \alpha_{1} - \alpha_{0} + 1 \]
linear relations.

It is clear that one will be able to obtain all these linear relations in the following manner: Let us consider a sequence of edges forming a closed contour; the algebraic sum of the differences $\delta$ relative\ednote{Printed ``relations''; apparently for ``relatives''.} to the various edges of this sequence must be zero.

Let us therefore seek to construct closed contours formed of edges.

We have first the polygonal contours of each face; they are $\alpha_{2}$ in number.
\origpage{115}

If next the polyhedron is not simply connected, we shall be able to trace on its surface $P_{1} - 1$ linearly independent closed contours, in the sense given to that word in the paragraph on homologies. Let $C$ be one of these contours; it will traverse different faces in succession; let $a_{1}a_{2}, a_{2}a_{3}, \ldots, a_{n-1}a_{n}, a_{n}a_{1}$ be the arcs of this contour that lie in each of these faces.

Let us take the first of these arcs $a_{1}a_{2}$, and $F$ the corresponding face.

The point $a_{1}$ and the point $a_{2}$ are on the perimeter of this face; we can then go from $a_{1}$ to $a_{2}$ following this perimeter by a path which I shall call $a_{1}m_{1}a_{2}$. We shall then have
\[ a_{1}m_{1}a_{2} + a_{2}a_{1} \backsim 0, \]
that is to say that, in the contour $C$, one can replace the arc $a_{1}a_{2}$ by the arc $a_{1}m_{1}a_{2}$; let us do the same for the other arcs of $C$; we shall finally have replaced $C$ by the homologous contour
\[ a_{1}m_{1}a_{2} + a_{2}m_{2}a_{3} + \ldots + a_{n}m_{n}a_{1}, \]
which I shall call $C'$. This contour $C'$ will consist of a certain number of edges and of portions of edges. For example, the arc $a_{1}m_{1}a_{2}$ will consist of a segment of an edge joining $a_{1}$ to the nearest vertex, then of a certain number of complete edges, then of a segment $Sa_{2}$ joining $a_{2}$ to the nearest vertex.

But this segment $Sa_{2}$ will be found traversed in the opposite sense in the arc $a_{2}m_{2}a_{3}$. The portions of edges that form part of $C'$ are therefore traversed twice in opposite senses; we can therefore suppress them, and we thus obtain a contour $C''$ formed exclusively of complete edges and homologous to $C$ or to $C'$.

We shall have $P_{1} - 1$ contours analogous to $C''$, which will give us $P_{1} - 1$ relations between the $\delta$.

We thus obtain $\alpha_{2} + P_{1} - 1$ closed contours formed of edges; I say that all the other possible contours are only combinations of them.

Let there be first a closed contour formed of edges and homologous to 0; it
\origpage{116}
will divide the surface of the polyhedron into two regions. Let $R$ be one of them; it will evidently be formed of a certain number $q$ of faces, since the contour is formed exclusively of edges; we shall therefore be able to replace the given contour by $q$ partial contours formed by the perimeters of these $q$ faces.

If now the given contour is not homologous to 0, we shall always be able to replace it by a combination of the contours $C''$ and by a contour homologous to 0.

We therefore have
\[ \alpha_{2} + P_{1} - 1 \]
relations between our $\delta$, and we have no others; but are they all distinct?

To recognize this, one must see whether one can form a linear combination of these relations that reduces to an identity, or, what comes to the same, whether one can form a combination of our $\alpha_{2} + P_{1} - 1$ contours, a combination such that each edge is traversed twice in opposite senses (or, if it is traversed more than twice, that it is traversed as many times in one sense as in the other).

Can one first form such a combination with the $\alpha_{2}$ perimeters $\Pi$? To say that each edge is thus traversed twice in opposite senses is to say that the set of the polygons whose perimeters are thus traversed forms a closed polyhedron.

Now one can evidently construct in this way only a single closed polyhedron, which is the given polyhedron.

With the $\alpha_{2}$ relations corresponding to the $\alpha_{2}$ perimeters $\Pi$ we can therefore form one identical linear combination and only one.

Can one now form such combinations with the perimeters $\Pi$ and the contours $C''$? If this were so, the set of polygons whose perimeters are thus traversed would form a non-closed polyhedral surface whose complete boundary would be formed by a combination of the contours $C''$. Now this is impossible, since the contours $C''$ are linearly independent.
\origpage{117}

One can therefore form with our $\alpha_{2} + P_{1} - 1$ relations one identical combination and only one. There are therefore between the $\delta$
\[ (\alpha_{2} + P_{1} - 1) - 1 \]
distinct relations. The number of arbitrary $\delta$ is therefore
\[ \alpha_{1} - (\alpha_{2} + P_{1} - 1) + 1, \]
so that there results
\[ \alpha_{0} - 1 = \alpha_{1} - (\alpha_{2} + P_{1} - 1) + 1, \]
or
\[ N = 3 - P_{1}. \]

Let us extend this proof to the case where we have a polyhedron of three dimensions, in which one will distinguish:
\[ \begin{matrix} \text{The vertices} & v_{0} & \text{in number} & \alpha_{0}, \\[1ex] \text{the manifolds} & v_{1} & \text{»} & \alpha_{1}, \\[1ex] \text{»} & v_{2} & \text{»} & \alpha_{2}, \\[1ex] \text{»} & v_{3} & \text{»} & \alpha_{3}. \end{matrix} \]

To each of the $v_{0}$ let us make correspond a number, and to each of the $v_{1}$ the difference $\delta$ of the numbers corresponding to its two vertices. We shall thus have $\alpha_{1}$ differences $\delta$, of which $\alpha_{0} - 1$ will be arbitrary.

One will obtain the linear relations that hold between the $\delta$ by constructing all the closed contours formed exclusively with $v_{1}$.

We shall have first the perimeters $\Pi$ of the $v_{2}$, which will be $\alpha_{2}$ in number; next let $C$ be any closed contour not homologous to 0; it will traverse different $v_{3}$ in succession; let $a_{1}a_{2}$, $a_{2}a_{3}$, \ldots\ be the arcs of $C$ that lie in each of these $v_{3}$.

Let us consider the arc $a_{1}a_{2}$, which lies in one of the $v_{3}$, which I call $\Phi$, and whose extremities $a_{1}$ and $a_{2}$ lie in two of the $v_{2}$ that serve as boundary to $\Phi$, namely $a_{1}$ in $\varphi_{1}$ and $a_{2}$ in $\varphi_{2}$.

Let $b_{1}$ be a vertex of $\varphi_{1}$, $b_{2}$ a vertex of $\varphi_{2}$; let us join $a_{1}$ to $b_{1}$ by
\origpage{118}
any line $a_{1}b_{1}$, then $b_{1}$ to $b_{2}$ by a line $b_{1}b_{2}$ formed exclusively by $v_{1}$ belonging to the boundary of $\Phi$, and $b_{2}$ to $a_{2}$ by any line $a_{2}b_{2}$; one will have
\[ a_{1}a_{2} \backsim a_{1}b_{1} + b_{1}b_{2} + b_{2}a_{2}, \]
so that we shall be able to replace the arc $a_{1}a_{2}$ by the arc $a_{1}b_{1}b_{2}a_{2}$; let us do the same for the other arcs of $C$. We shall have replaced $C$ by the homologous contour
\[ a_{1}b_{1}b_{2}a_{2} + a_{2}b_{2}b_{3}a_{3} + \ldots \]
which I call $C'$. The contour $C'$ consists of a certain number of $v_{1}$ and of arcs analogous to $a_{1}b_{1}$, $a_{2}b_{2}$, \ldots, each traversed twice in opposite senses. One can therefore suppress these arcs, and there remains a contour $C''$ homologous to $C$ and formed exclusively of $v_{1}$.

There will be $P_{1} - 1$ contours $C''$.

I say now that every closed contour $K$ formed of $v_{1}$ is a combination of the $\Pi$ and the $C''$. If
\[ K \backsim 0, \]
the contour $K$ will be the complete boundary of a certain region $R$ of two dimensions. This region $R$ can be subdivided into a certain number of manifolds $r$, each of the regions $r$ being the portion of $R$ that lies in the interior of one of the $v_{3}$. Let us consider one of the manifolds $r$, and let $\varphi$ be the region $v_{3}$ in which it lies; the complete boundary of $r$ will be a closed manifold $u$ of one dimension, which will form part of the complete boundary of $\varphi$. Since $\varphi$ is simply connected, $u$ will divide the complete boundary of $\varphi$ into two regions; let $r'$ be one of these regions; it will consist of a certain number of $v_{2}$ that form part of it entirely, and of a certain number of portions of $v_{2}$ (since, among the $v_{2}$ that form the complete boundary of $\varphi$, there are some that are divided into two parts by $u$). One sees that $r$ is homologous to $r'$; we shall be able to replace $r$ by $r'$; if we do the same for all the regions $r$, we shall have obtained a manifold
\origpage{119}
$R'$ homologous to $R$ and which will consist of a certain number of complete $v_{2}$ and of a certain number of portions of $v_{2}$, each taken twice in opposite senses. We shall be able to suppress these portions of $v_{2}$, and we shall thus obtain a manifold $R''$ homologous to $R$ and bounded by the same contour $K$. This manifold $R''$ can be decomposed into a certain number of polygons $v_{2}$, so that the contour $K$ can be replaced by the perimeters $\Pi$ of these polygons.

If $K$ is not homologous to 0, one will be able to replace it by a certain number of contours $C''$ and by a contour analogous\ednote{Printed ``analogue''; apparently for ``homologue'', as in the corresponding sentence on p. 116.} to 0.

We therefore have
\[ \alpha_{2} + P_{1} - 1 \]
relations between the $\delta$, which I shall write
\[ \varepsilon = 0, \]
and we shall have no others. But are they all distinct?

In other words, can we form a linear combination of the $\varepsilon$ that is identically zero, or again a combination of the $\Pi$ and the $C''$ in which each $v_{1}$ appears twice in opposite senses?

If $C''$ had to appear in this combination, the set of the polygons $v_{2}$ whose perimeters $\Pi$ appear in the combination would form a manifold of two dimensions which would not be closed and whose complete boundary would be formed by a certain number of contours $C''$. Now this is not possible, since the $C''$ are linearly independent.

It therefore remains for me to examine the combinations in which only $\Pi$ would appear; the set of the polygons $v_{2}$ whose perimeters $\Pi$ appear in them would then form a closed manifold.

We are thus led to examine the closed manifolds formed exclusively of $v_{2}$.

We have first the complete boundaries of the $v_{3}$, complete boundaries which I shall call $\Phi$, and which are $\alpha_{3}$ in number.

Next let $D$ be any manifold of two dimensions not homologous
\origpage{120}
to 0. Let us treat it as we treated $R$; we shall see that it is homologous to a closed manifold $D''$ of two dimensions, formed exclusively of $v_{2}$.

The number of the $D''$ is $P_{2} - 1$.

I say now that every closed manifold $K$ formed of $v_{2}$ is a combination of the $\Pi$ and the $D''$. If, indeed, it is homologous to 0, it bounds a region $S$ of three dimensions which will consist of a certain number of $v_{3}$, since $K$ consists of a certain number of $v_{2}$; one will therefore be able to replace $K$ by the complete boundaries $\Phi$ of these $v_{3}$. If $K$ is not homologous to 0, one will be able to replace it by a certain number of $D''$ and by a manifold homologous to 0 and formed exclusively of $v_{2}$.

We therefore have, between the $\varepsilon$,
\[ \alpha_{3} + P_{2} - 1 \]
linear relations, which I shall write
\[ \zeta = 0, \]
and we shall have no others. Are they distinct?

To form a linear combination of the $\zeta$ that is identically zero, one must form a combination of the $\Phi$ and the $D''$ such that each $v_{2}$ appears in it twice in opposite senses. We should see, as above, that the $D''$ cannot appear in this combination, and that the set of the $v_{3}$ whose boundaries $\Phi$ appear in it must form a closed manifold of three dimensions. Now we can construct only a single manifold of this sort: it is the given polyhedron itself.

There is therefore one linear combination of the $\zeta$ that vanishes identically.

There will therefore be
\[ (\alpha_{3} + P_{2} - 1) - 1 \]
distinct linear relations between the $\varepsilon$, and
\[ (\alpha_{2} + P_{1} - 1) - (\alpha_{3} + P_{2} - 1) + 1 \]
distinct linear relations between the $\delta$.
\origpage{121}

Among the $\delta$ there will therefore remain
\[ \alpha_{1} - (\alpha_{2} + P_{1} - 1) + (\alpha_{3} + P_{2} - 1) - 1 \]
that will be arbitrary; one thus finds
\[ \alpha_{0} - 1 = \alpha_{1} - (\alpha_{2} + P_{1} - 1) + (\alpha_{3} + P_{2} - 1) - 1 \]
and one would find
\[ \alpha_{0} - 1 = \alpha_{1} - (\alpha_{2} + P_{1} - 1) + (\alpha_{3} + P_{2} - 1) - (\alpha_{4} + P_{3} - 1) + 1 \]
for polyhedra of four dimensions.

One therefore has
\[ N = P_{2} - P_{1} \]
for polyhedra of three dimensions, and
\[ N = 3 - P_{1} + P_{2} - P_{3} \]
for polyhedra of four dimensions.

In general, one will have
\[ N = P_{p-1} - P_{p-2} + \ldots + P_{2} - P_{1} \]
if $p$ is odd, and
\[ N = 3 - P_{1} + P_{2} - \ldots + P_{p-1} \]\ednote{The last sign is printed $+$; since $p-1$ is odd when $p$ is even, apparently for $-$, as in the case $p = 4$ above.}
if $p$ is even.

If we now observe that the Betti numbers equally distant from the extremes are equal, we shall see that one must have
\[ N = 0 \]
if $p$ is odd, as we had already seen in the preceding number\ednote{No final period is printed after ``précédent''; the memoir ends here.}

\end{document}
